\documentclass[10pt]{article}

% Shared presentation layer for the standalone R0--R7 development documents.
% Method equations remain authoritative in the canonical idea sketch.
\usepackage[a4paper,margin=18mm,headheight=14pt]{geometry}
\usepackage{fontspec}
\usepackage{xeCJK}
\xeCJKsetup{CJKspace=true}
\setmainfont{DejaVu Serif}
\setsansfont{DejaVu Sans}
\setmonofont{DejaVu Sans Mono}
\setCJKmainfont{Noto Serif CJK KR}
\setCJKsansfont{Noto Sans CJK KR}
\setCJKmonofont{Noto Sans Mono CJK KR}

\usepackage{amsmath,amssymb}
\usepackage{booktabs,longtable,tabularx,array}
\usepackage{enumitem}
\usepackage{xcolor}
\usepackage{hyperref}
\usepackage{fancyhdr}
\usepackage{microtype}
\usepackage[most]{tcolorbox}

\definecolor{CoreBlue}{HTML}{1F4E79}
\definecolor{CoreLight}{HTML}{EAF3F8}
\definecolor{StopRed}{HTML}{9C0006}
\definecolor{StopLight}{HTML}{FCE8E6}
\definecolor{GateGreen}{HTML}{38761D}
\definecolor{GateLight}{HTML}{EAF4E2}
\definecolor{DecisionOrange}{HTML}{B45F06}
\definecolor{DecisionLight}{HTML}{FFF2CC}

\hypersetup{
  colorlinks=true,
  linkcolor=CoreBlue,
  urlcolor=CoreBlue
}

\setlength{\parindent}{0pt}
\setlength{\parskip}{0.3em}
\setlength{\emergencystretch}{3em}
\setlist[itemize]{leftmargin=1.45em,itemsep=0.1em,topsep=0.2em}
\setlist[enumerate]{leftmargin=1.75em,itemsep=0.14em,topsep=0.22em}
\renewcommand{\arraystretch}{1.15}
\newcolumntype{L}[1]{>{\raggedright\arraybackslash}p{#1}}
\newcommand{\code}[1]{{\small\nolinkurl{#1}}}
\newcommand{\must}[1]{\textcolor{CoreBlue}{\textbf{#1}}}
\newcommand{\haltitem}[1]{\textcolor{StopRed}{\textbf{#1}}}
\newcommand{\passitem}[1]{\textcolor{GateGreen}{\textbf{#1}}}
\newcommand{\decisionitem}[1]{\textcolor{DecisionOrange}{\textbf{#1}}}

\newenvironment{contractbox}[1]
  {\begin{tcolorbox}[colback=CoreLight,colframe=CoreBlue,title={#1},fonttitle=\bfseries,before upper=\raggedright]}
  {\end{tcolorbox}}
\newenvironment{decisionbox}[1]
  {\begin{tcolorbox}[colback=DecisionLight,colframe=DecisionOrange,title={#1},fonttitle=\bfseries,before upper=\raggedright]}
  {\end{tcolorbox}}
\newenvironment{stopbox}[1]
  {\begin{tcolorbox}[colback=StopLight,colframe=StopRed,title={#1},fonttitle=\bfseries,before upper=\raggedright]}
  {\end{tcolorbox}}

\usepackage{fvextra}
\usepackage{xurl}
\DefineVerbatimEnvironment{verbatim}{Verbatim}{breaklines=true,fontsize=\scriptsize}

\hypersetup{
  pdftitle={R1 Teacher, Probe, and Oracle --- Wind3DGS Development Specification},
  pdfauthor={Wind3DGS Project}
}

\pagestyle{fancy}
\fancyhf{}
\lhead{\small Wind3DGS Development --- R1}
\rhead{\small Teacher, Probe, and Oracle}
\cfoot{\thepage}

\title{\textbf{R1 --- Teacher, Probe, and Oracle}\\
Reference·mapping·network-free preflight 명세와 구현·검증 근거}
\author{Wind3DGS Project}
\date{2026-08-22; 구현·검증 근거 갱신 2026-09-22}
\renewcommand{\DevChecklistDate}{2026-09-22}

\begin{document}
\maketitle

\begin{contractbox}{문서 권위와 선행 조건}
이 문서는 canonical sketch의 R1 구현 파생 문서다. 최종 방법 권위는 sketch에 있으며,
\code{eq:rd-common-probe-map}, \code{eq:rd-common-probe-adjoint},
\code{eq:rd-area-mass}, \code{eq:rd-attachment-whitening},
\code{eq:rd-global-response}, \code{eq:rd-exact-zoh},
\code{eq:rd-wind-pullback}, \code{eq:rd-mean-transport},
\code{eq:rd-teacher-period-spectrum}을 텍스트 기준으로 참조한다.
R0 contract/schema와 그 fixture가 통과하기 전에는 이 문서의 실행 상태를 완료로 바꾸지 않는다.
\end{contractbox}

\begin{contractbox}{이 문서에서 현행 계약과 근거를 함께 읽는 법}
R1의 실행 명세, 구현 상태와 논문용 근거를 이 문서 하나에서 관리한다.
수식·계약은 해당 절에서 직접 갱신하고 변경 이유를 인접 설명과 TeX 주석에 남긴다.
\S\ref{sec:r1-contracts}은 공통 계약, \S\ref{sec:r1-implementation}은 구현 상태,
\S\ref{sec:r1-decisions}은 현행 개발 후보와 교체 근거,
\S\ref{sec:r1-artifacts}은 산출물·재현 경로,
\S\ref{sec:r1-fixtures}은 검증 결과, \S\ref{sec:r1-exit}은 최종 채택 조건이다.
실패한 식은 교체 이유를 설명하는 위치에만 당시 적용 범위를 표시해 보존한다.
기존 개발 window 15개와 후속 P3 작은 굽힘 sample 76개의 채택 범위를 구분한다. R1 전체는 미완료다.
\end{contractbox}

\DevChecklistPage{R1}{개발 진단 진행 · R1 미완료}{
\DevProgressRow{0}{\hyperref[sec:r1-chapter-1]{1. 목적과 소유권}}{진행}{Teacher 개발 scope와 canonical acceptance 경계 대조}
\DevProgressRow{0}{\hyperref[sec:r1-chapter-2]{2. 진입 입력}}{대기}{R0 입력과 independent GS 2개 등 진입 artifact 확보}
\DevProgressRow{0}{\hyperref[sec:r1-chapter-3]{3. 이미 동결된 계약}}{진행}{동일 backend 수렴, 고차 map와 force/work 계약 검증}
\DevProgressRow{0}{\hyperref[sec:r1-chapter-4]{4. 구현 체크리스트}}{진행}{개발 경로에서 accepted Teacher·oracle·transport로 연결}
\DevProgressRow{0}{\hyperref[sec:r1-chapter-5]{5. Open Design Decisions}}{진행}{FP64 hi/lo GPU 개발 경로 유지; 접촉·최종 입력 범위는 미결정}
\DevProgressRow{0}{\hyperref[sec:r1-chapter-6]{6. 산출물}}{진행}{개발 evidence는 보존; accepted artifact 생성 필요}
\DevProgressRow{0}{\hyperref[sec:r1-chapter-7]{7. 필수 fixture와 metric}}{진행}{4배 공간 미달·10초 탐색 추천 없음; 독립 기준·oracle 검증}
\DevProgressRow{0}{\hyperref[sec:r1-chapter-8]{8. 종료 조건과 실패 경로}}{대기}{R1 전체 종료 조건과 R2 hand-off 통과}
}{
\begin{contractbox}{별도로 확인된 개발 성과}
\begin{itemize}[label=\DevDoneMark]
 \item Registry/trajectory/sequence와 P1 probe 구현·개발 검증 (4장).
 \item 사용자 GPU smoke 12/12 실행 통과 (7장; 물리 수렴 판정과 구분).
 \item Newton 3개 source/15개 window의 원본·replay·loader 검증 (4장).
 \item 기존 shell의 시간 진단과 P3 압력의 공간 진단 통과 (7장; 서로 다른 범위).
 \item 변화 바람/reset 25개 trace의 재생·개입·공력 검산 (5장; 속도 수렴 실패).
 \item P3 작은 굽힘의 새 wind/BC 수렴·재실행과 sample 76개 검증 (5장).
 \item P3 유한 회전 수식·GPU CPU 대조·짧은 wind/reset 수렴 (4장).
 \item 약한 바람1.5 s의 자연/세 reset 원식·시간·공간·방향 검증 (4장).
\item CPU 풀이+HVP/CUDA graph의 전체1.5 s 비교·검산 근거 확보 (4장).
 \item CPU/GPU 셀프 접촉의 LBVH·미분·CCD 대조, 정밀 기하 인증·실제 퇴화 거절·GPU 프레임 복원과 세 씬 연결 검증 (4장).
\end{itemize}
\textbf{남은 판정:} 큰 변형 범위·독립 비선형 공간 기준 및 oracle/transport.
과거 초기 변위의 실패는 해당 fixture의 반례로 보존한다.
개발 sample은 \code{training_eligible=false}다.
\end{contractbox}
\textbf{다음 조건:} 유한 회전 shell의 수식·수렴 검증 우선; 추가 학습데이터는 보류 (4--5장).
}

\tableofcontents
\newpage

\section{목적과 소유권}
\label{sec:r1-chapter-1}
R1은 network를 학습하기 전에 \textbf{reference가 믿을 수 있는지}, teacher와 GS가 같은 물리량을
비교하는지, 그리고 현재 response-package class가 적어도 network 없이 teacher response를 표현할 수 있는지를
확인한다. 여기서 실패하면 R2의 optimizer나 architecture를 조정하지 않는다.

\begin{longtable}{L{0.24\linewidth}L{0.31\linewidth}L{0.35\linewidth}}
\toprule
소유자 & R1가 소유하는 것 & 소유하지 않는 것\\
\midrule
\endfirsthead
\toprule
소유자 & R1가 소유하는 것 & 소유하지 않는 것\\
\midrule
\endhead
Teacher physics & 구조 model/constitutive/material/damping/solver registry, quadrature, spatial/time convergence & student architecture와 learned result\\
Common probe & canonical probe, teacher/GS map, common-valid mask, mapping-noise floor와 payload별 adjoint & target runtime force path\\
Method identity & teacher와 GS에 동일한 traction law/effective scale/guard/sign convention, reduced evaluator exact-ZOH fixture & teacher 구조 solver를 exact ZOH로 제한하는 것\\
% [논문 작성 참고 -- 수정 이유]
% coefficient를 rho/C_n/C_parallel tuple로 해석하지 않고 식별 가능한 kappa_aero 하나로 축소했다.
Oracle & attachment/mass-constrained network-free Global response preflight & learned held-out evidence 또는 MC1 claim\\
Base transport & teacher surface motion의 independent GS mean/covariance transport, attachment/rest/SPD fixture & optional SH appearance와 perceptual loss\\
Spectrum status & teacher spectrum convergence와 stable-peak availability status & curriculum 채택 여부와 학습 이득 판정\\
\bottomrule
\end{longtable}

\must{R1 성공의 의미:} 두 independent GS variant가 하나의 frozen common-probe denominator에서
converged teacher와 비교 가능하고, corrected oracle package가 unit/map/measure/evaluator error floor 안에서
response를 재현하며, base mean/covariance transport가 physics와 renderer 양쪽 fixture를 통과한다.

\subsection{작업 범위와 논문 근거의 소유권}

Static 3DGS에서 wind-driven passive response package를 학습하려면 먼저 mesh teacher의 응답이
discretization이나 적분기 오차에 지배되지 않는지 확인해야 한다. 이 작업은 training-only reference를
만드는 과정이다. Target-mesh-free inference, student의 positive area/mass 및 Global/Local package
구조는 바꾸지 않았다. 고차 판 자체의 신규성을 이 작업의 논문 contribution으로 주장하지 않는다.

출발점은 \code{code/sessions/2026-09-06_01_teacher_dataset_handoff.md}의 기본 mesh/GUI였다.
후속 요청에 따라 registry, 정량 저장·비교, 실패 진단, 개발 sample까지 단계적으로 구현했다.
GUI 이후의 주된 확장은 물리 reference를 검증하고 데이터 provenance를 보존하는 CPU/CLI 경로다.
기존 Newton viewer가 새 P3 backend나 모든 정량 보고서를 자동으로 표시하는 것은 아니다.
그래프는 sample response, 공간 수렴 곡선과 내부 힘/모드 비교를 별도 artifact로 보존한다.

사용자는 필요한 선택에는 질문하되 연속 진행하도록 승인했다.
2026-09-09 후속 지시로 \textbf{물리 수식을 제대로 구현·검증한 뒤 학습데이터를 생성}하는 순서를
명확히 했다. 추가 sample 발행은 보류하고 기존 76개 작은 굽힘 sample은 역사적 개발 근거로 보존한다.
2026-09-09 사용자는 rest-start 변화 바람과 중간 형상의 velocity-reset 비교를 선택했고,
후속 진행 요청으로 개발 비교 25개 trace를 실행·검산했다 (\S\ref{sec:r1-velocity-reset-plan}).
속도 수렴과 굽힘 상태 coverage의 한계를 확인했으며 새 학습 결과는 없다.
첫 accepted 데이터의 최종 입력 범위와 displaced free-decay의 수렴은 이 탐색 선택만으로 확정하지 않는다.

\section{진입 입력}
\label{sec:r1-chapter-2}
\subsection{R0에서 받아야 하는 artifact}

\begin{itemize}
  \item Frozen \code{DomainContract}, \code{CoordinateContract}와 \code{PhysicalValidPolicy}.
  \item Frozen \code{TensorSchemaRegistry}, \code{SerializationContract}와 method identity.
  \item \code{TeacherStudentVisibility}, source-object group rule, sealed-test hash와 failure denominator registry.
  \item \code{MetricRegistry}/\code{ThresholdRegistry} schema 및 R1에서 채울 threshold owner.
  \item Normal/covariance branch, rank-failure policy, stable Gaussian ID와 package round-trip fixture.
\end{itemize}

\subsection{K0 최소 입력 세트}

\begin{itemize}
  \item Flat-rest strip 1개와 rectangular flag 1개, 우선 material preset 1개.
  \item Object마다 camera/seed/densification이 독립적인 static GS variant 2개. Primitive split/duplicate는 독립 variant가 아니다.
  \item Short impulse, step-on/off, zero-ambient rest/recovery와 aerodynamic-off \(Q=0\) displaced free decay.
  \item Spatial mesh refinement 최소 2단계와 temporal timestep refinement 최소 2단계.
  가능한 한 three-level convergence를 사용하되 실제 level 수와 ratio는 Open Design Decision에 기록한다.
  \item Same rest surface의 metric/material/attachment/wind program과 replay seed/hash.
\end{itemize}

\subsection{Teacher 입력 payload}

Teacher는 mesh vertex/element connectivity, surface quadrature, material, attachment와 prescribed wind를 소유한다.
Reference export는 canonical surface sample의 rest/current position, displacement, velocity, optional normal,
applied traction/total force, external work와 timestamp를 포함한다. Topology와 correspondence는 training/evaluation
artifact에만 남고 target package로 serialize하지 않는다.

\section{이미 동결된 계약}
\label{sec:r1-chapter-3}
\label{sec:r1-contracts}

\subsection{Privilege와 convergence separation}

\begin{itemize}
  \item Teacher mesh는 privileged simulation/evaluation data다. Vertex 하나를 latent anchor 하나로 대응시키지 않는다.
  \item Accepted fine timestep을 고정한 spatial mesh refinement와 accepted fine mesh를 고정한 temporal timestep
  refinement를 별도로 측정한다. 둘을 동시에 바꾼 diagonal difference로 cancellation을 허용하지 않는다.
  두 축은 같은 element/mass/structure/damping/attachment/load identity를 공유해야 한다.
  기존 quadratic patch shell의 temporal pass와 P3 plate의 spatial pass를 합쳐 승인하지 않는다.
  \item 두 convergence는 같은 canonical probe와 physical-time grid에서 velocity, tip trajectory, work 및
  band-integrated spectrum/FRF discrepancy를 측정한다.
  \item Dominant peak의 존재는 teacher convergence 조건이 아니다. Spectrum이 수렴했지만 clear peak가 없으면
  \code{T_dom=unavailable} status이며 teacher failure가 아니다. Spectrum 자체가 수렴하지 않으면 R1 failure다.
\end{itemize}

\subsection{Teacher quadrature와 질량의 현행 계약}
\label{sec:r1-mass}

Teacher quadrature도 canonical \(A_{\mathrm{ref}},M_{\mathrm{ref}}\)를 사용한다.
Shell thickness/density를 별도 SI parameter로 저장하더라도
\(h\rho=M_{\mathrm{ref}}/A_{\mathrm{ref}}\)를 derive/equality-constrain한다.
Constitutive preset은 독립 total mass owner가 아니다.
양의 면적 quadrature \(a_q\)와 그 위치에서의 shape evaluation \(\Psi\)에 대해
\begin{equation}
\begin{aligned}
 W_{A,T}&=\operatorname{diag}(a_q),\qquad a_q>0,\qquad \sum_q a_q=A_{\mathrm{ref}},\\
 m_q&=\frac{M_{\mathrm{ref}}}{A_{\mathrm{ref}}}a_q,\qquad \sum_q m_q=M_{\mathrm{ref}},\\
 M_T^{\mathrm{scalar}}&=\frac{M_{\mathrm{ref}}}{A_{\mathrm{ref}}}\Psi^\top W_{A,T}\Psi,
 \qquad M_T=M_T^{\mathrm{scalar}}\otimes I_3,\\
 \mathbf1^\top M_T^{\mathrm{scalar}}\mathbf1&=M_{\mathrm{ref}}
 \quad\text{(partition of unity를 만족하는 전체 자유도에서)}.
\end{aligned}
\label{eq:r1-consistent-mass}
\end{equation}
Constraint 적용 후 허용 DOF에서는 질량행렬의 양정성과 quadrature rank를 검사한다.
이는 sketch \code{eq:rd-teacher-consistent-mass}의 구현 계약이며 student GS의 positive area와
diagonal mass를 바꾸지 않는다. 기존 선형 teacher의 lumped mass는 별도 law로 구분한다.

% [수식 변경 이유 -- 2026-09-09]
% Vertex/고차 DOF마다 양의 point mass를 붙이는 해석을 positive quadrature와 consistent mass로 교체했다.
% m_q는 quadrature mass다. 고차 shape의 signed weight 또는 DOF row sum과 혼동하면 안 된다.
\textbf{변경 이유:} P3 DOF는 양의 면적을 가진 sample과 같지 않으므로
\(\sum m_q=M_{\mathrm{ref}}\)를 DOF별 positive lumped-mass 제약으로 사용하지 않는다.
고차 registry/validator의 실제 연결은 아직 남아 있다.

\subsection{Mandatory common-probe map}

Source object마다 test open 전에
\[
  \mathcal X_P^0=\{(p_j^0,w_{A,j}^P,w_{M,j}^P)\}_{j=1}^{N_P}
\]
를 동결한다. Sketch의 \code{eq:rd-common-probe-map}에 따라
teacher discretization \(X=T\)와 GS variant \(X=r\) 각각에 대해
\begin{equation}
 T_X\in\mathbb R^{N_P\times N_X},\quad
 \sum_i(T_X)_{ji}=1,\quad S_X=T_X\otimes I_3,\qquad
 u_P=S_Tu_T,\quad v_P=S_Tv_T.
 \label{eq:r1-probe-map}
\end{equation}
% [수식 변경 이유 -- 2026-09-09]
% 공통 식의 전역 nonnegative 조건을 제거하고 아래처럼 map law별 조건으로 교체한다.
% 기존 P1과 GS convex map의 조건은 유지한다. 고차 Teacher만 signed shape를 허용한다.
를 만든다. Ordered support/weight, coverage, kernel/version와 hash를 저장한다.
GS convex map과 기존 선형 barycentric teacher에는 \((T_X)_{ji}\ge0\)를 요구한다.
고차 teacher의 signed shape weight는 별도 map-law/version으로 구분한다.
해당 차수의 polynomial reproduction, conditioning/noise amplification과 virtual-work adjoint를
검증하며 양의 surface quadrature weight와 혼동하지 않는다. 현재 public teacher map v1은
선형/nonnegative 계약만 구현했다. P3 진단 내부의 shape evaluation은 그 API의 완료 근거가 아니다.

\begin{itemize}
  \item 모든 accepted map은 row partition, arbitrary constant vector reproduction과 predeclared
  rest-coordinate affine/smooth analytic field reproduction tolerance를 통과한다.
  \item Teacher와 source-object group의 모든 declared GS variant가 함께 support하는 probe intersection으로
  pre-test frozen common-valid mask 하나를 만든다. Variant별 사후 probe 삭제는 금지한다.
  \item Coverage distance/rate, unsupported reason과 measured mapping-noise floor를 denominator와 함께 저장한다.
  \item 이 map은 training/evaluation-only다. Target runtime GS force는 Gaussian에서 직접 계산한다.
\end{itemize}

\subsection{Payload별 adjoint}
\label{sec:r1-adjoint}

Sketch \code{eq:rd-common-probe-adjoint}에 따라 payload 의미를 구분한다.
\[
  F_X=S_X^\top F_P \quad\text{(probe별 total force)},
  \qquad
  F_X=S_X^\top W_{A,P}\tau_P \quad\text{(probe traction)},
\]
\[
  W_{A,P}=\operatorname{diag}(w_{A,j}^PI_3),\qquad
  \tau_P^\top W_{A,P}S_Xu_X=F_X^\top u_X.
\]
Traction에 \(S_X^\top\)만 적용하거나 area를 두 번 곱하는 구현은 금지한다.
Runtime의 \(F_i=a_i\tau_i\)에는 probe adjoint를 재적용하지 않는다.

\subsection{동일 traction identity와 reduced evaluator}

Teacher quadrature와 GS는 representation만 다르고 다음 traction identity는 같다.
\begin{itemize}
  \item \code{traction_law_id=two_sided_normal_quadratic_v1}: 같은
  domain-wide \(\kappa_{\mathrm{aero}}\), relative-velocity/sign convention과 two-sided normal-only law.
  % [논문 작성 참고 -- 수정 이유]
  % kappa_aero=0.5*rho_air*C_n으로 합쳐 동일 공기 조건에서 구분 불가능한 scalar를 제거했다.
  % 별도 C_parallel/lift는 response distillation이 핵심인 V0의 공력 자유도를 불필요하게 늘리므로 제외한다.
  \item \code{traction_guard_id=vector_norm_before_area_v1}: 같은 vector-norm
  \(\tau_{\mathrm{guard}}\)를 먼저 적용한 뒤 각 discretization의 area quadrature를 곱한다.
  Core quantitative run에서는 두 representation 모두 activation count가 0이어야 한다.
  % [논문 작성 참고 -- 수정 이유]
  % guard threshold는 registry/hash에 남기되 물리계수로 소개하지 않는다. activation은 clipping 성공이 아니라
  % supported envelope 이탈 또는 안정성 실패의 증거다.
  \item Zero ambient와 aero-off를 서로 다른 flag로 기록한다. 전자는 drag를 허용하고 후자만 strict free response다.
  \item 학생 reduced evaluator는 \code{integrator_id=exact_zoh_augexp_v1},
  \code{force_sample_time=frame_start_v1}로 고정한다. Teacher 구조 solver는 별도 converged reference이며
  같은 structural integrator일 필요는 없다. 그러나 exported frame의 aerodynamic traction은 학생과 같은
  frame-start sample을 해당 frame 동안 hold한다. Teacher 내부 structural substep과 work quadrature만 별도
  convention으로 기록할 수 있으며 force sample identity 자체를 바꾸지 않는다.
\end{itemize}

\subsection{Corrected oracle invariant}

Oracle은 attached row가 exact zero인 raw Global field를 common-probe teacher impulse/step/FRF response에 fit하고,
R0 physical-valid policy와 R1에서 동결한 \textbf{oracle-only positive area source}로 만든 mass를 사용해
whitening한다. Network-free 단계이므로 이 area는 student head 출력이 아니며, teacher quadrature에 대한
training-only fit 또는 deterministic GS footprint rule 중 하나를 Open Design Decision으로 동결한다.
어느 경로든 accepted row에서 \(a_i\ge0\), \(\sum_i a_i=A_{\mathrm{ref}}\)와
\(m_i=(M_{\mathrm{ref}}/A_{\mathrm{ref}})a_i\)를 만족해야 한다.
이 oracle measure는 evaluator preflight 전용이며 R3의 learned-measure evidence로 승계하지 않는다.

\(B_c\)를 \(B=B_cR\)로 바꾸면 generalized coordinate도 바뀌므로
teacher의 기존 diagonal \(\Omega,Z\)를 그대로 붙이지 않는다. 다음 두 mathematically consistent 경로 중
\textbf{하나만} canonical로 선택할 수 있으나, 어느 경로인지는 아직 Open Design Decision이다.
\begin{enumerate}
  \item Whitened \(B\)를 고정하고 common-probe time/FRF response로 positive diagonal \(\Omega,Z\)를 다시 fit한다.
  \item Raw full reduced \(M_c,C_c,K_c\)를 whitening congruence-transform한 뒤 일관되게 re-diagonalize한다.
\end{enumerate}
두 경로 모두 student와 같은 positive-diagonal response-package class와 exact-ZOH evaluator로 귀착해야 한다.
Oracle paired 결과는 contract preflight이며 learned held-out evidence가 아니다.

\subsection{Base mean/covariance transport}

Mean은 response field를 통해 transport하고, R0에서 선택한 angular-field 또는 local-affine branch로 current normal과
covariance orientation을 갱신한다. Attachment displacement/angular row는 exact zero다.
Rest state identity, proper rotation, covariance eigenvalue/SPD/finite 보존과 velocity consistency를 만족해야 한다.
Optional higher-order SH, scale/shear, deformation-gradient와 appearance residual은 R1 범위 밖이다.

\section{구현 체크리스트}
\label{sec:r1-chapter-4}
\label{sec:r1-implementation}

\subsection{2026-09-14 후속 GPU·중력·굽힘 구현 상태}
\label{sec:r1-gpu-gravity-update}
현행 개발 솔버는 FP64 hi/lo 상태와 Newmark 평균 가속도 적분을 유지한다.
Newton의 비선형 잔차 보정, 전처리 GMRES, 보조 행렬 갱신/풀이와 독립 검산을 GPU에 유지하며
병렬 합산, 첫 보조 풀이 중복 제거, current 우선, 채택 평가 재사용 및 무압축 저장을 적용했다.
생산 기본 경로는 유지한다. FP32 보조 연산의 사례별 개발 후보 동결은
\S\ref{sec:r1-bounded-closeout}에 별도로 기록하며 전체 FP32 채택과 구분한다.
캐시 무효화·graph/버퍼 수명과 API 계약은 \code{code/docs/gpu_solver_design.md} 및
\code{code/docs/gravity_wrinkle_test.md}를 따른다.

큰 회전에서 기존 평면 투영 충분조건 실패를 실제 자기 교차로 판정하지 않는다.
명시적 \code{local_metric} 모드는 변형률 상한을 통한 국소 비퇴화를 검사하며,
힘·위치 업데이트·에너지·고정점·비유한 값·솔버 검사는 유지한다.
기존 strict 정책을 삭제하거나 전역 비접촉 인증으로 대체한 것은 아니다.

중력은 consistent mass로 적분한 외력이며 프레임 시작에 공력과 함께 고정한다.
중력 ramp 포함2초 후 raw hi/lo 위치·속도를 보존해 무풍/바람 각각4초로 분기한다.
정착 대기·속도 reset·추가 감쇠는 없다. 기본은 왼쪽 고정 직사각형 깃발이다.
굽힘 비율은 Eh와 독립 면밀도를 유지하고 Eh³ 및 경계 굽힘 penalty를 조정한다.
1/100·1/300은 본 계산/검산/재생 완료 근거가 있고1/500은 별도 비교 대상으로
전체 완료 근거를 아직 채택하지 않았다. 각 run의 접촉 두께를 재료 h로 자동 결정하지 않는다.

기존 contact-off GPU 상주 경로는 유지한다. 2026-09-22의 별도 CPU 기준 및 GPU 접촉 후보는
아래 \S\ref{sec:r1-cpu-self-contact}, \S\ref{sec:r1-gpu-self-contact}에 기록하며,
마찰이나 기존 contact-off 연구 계약을 자동 해제하지 않는다.
GPU 기능 구현과 제한된 개발 검산 완료는 최종 동일-backend 공간/시간 수렴,
고차 probe/traction map, 독립 GS 공통 영역, network-free oracle 및 transport의 완료를 뜻하지 않는다.
R0 및 canonical R1 종료 조건은 유지한다.

\subsection{CPU 셀프 접촉 기준 후보 (2026-09-22)}
\label{sec:r1-cpu-self-contact}
사용자 선택에 따라 GPU 이식에 앞선 CPU 기준 경로를 구현했다. IPC Toolkit 1.6.0의
FP64 API를 사용하며 FP32 입력 변환을 FP32 solver 실행으로 해석하지 않는다.
P3 자유도를 고정 선형 collision proxy로 보간하고 힘과 Hessian을 원래 자유도로 adjoint한다.
LBVH의 VF/EE 후보를 정밀 검사하며 같은 macro element나 이웃 element 전체를 제외하지 않는다.
마찰 없는 고정 강성 IPC barrier를 기존 consistent-mass Newmark 잔차에 더하고,
ACCD로 predictor/Newton 보정 경로를 제한한다. C-IPC 전체 재료 모델이나 strain limiting을
채택한 것은 아니다. 같은 Newton 상태의 contact Hessian은 HVP 사이에 재사용한다.

시간 단계의 quadratic 위치 보간은 별도의 chord/tube 분할과 CCD로 검사한다.
유효 시작 간격을 위반한 상태는 CCD 호출 전에 거절하고, 미인증 구간도 실패로 처리한다.
원래 CPU solver 허용오차·재료식과 접촉 OFF 결과를 유지하며 비국소 Hessian 때문에
기존 shell-only current coloring은 접촉 경로에서 금지한다. 작은 기준 경로는 rest 보조 풀이를 쓴다.

관련 단위·회귀 검사는47개 통과, CUDA graph 전용1개 미실행이다. 두 저속 국소 접근의
접촉 ON 경로는 완료·검산됐고 LBVH/전수 탐색이 일치했다. 약한 barrier의 고속1m/s 접근은
dt/64에서도 시간 중간의 최소 간격 위반으로 거절되어 \textbf{전체 샘플 통과가 아니다}.
사용자 승인 후 barrier 계수1,000/10,000을 비교하여 세 사례 각각 LBVH/전수 탐색을
완주하고 각 강도264개 접촉 substep의 힘·시간 CCD를 독립 검산했다. 이는 후속 GPU 경로의
엄격한 \code{local_metric} 기하 승인까지 포함한 통과가 아니다. 시험한 값 중 최소 통과값1,000의
고속 사례를 추가 시간 세분화했으며0.25/0.125ms 끝점의 위치·속도 상대 차이는
각각0.0163\%,0.0303\%로 사전1\% 진단 한도를 통과했다. 전체 시간 궤적의 canonical 수렴은 아니다.
초기 활성 접촉 에너지와 반발도 강도에 따라 변하고,10,000의 시간 이산화 에너지 잔차는 더 크므로
큰 계수를 최적 물성이나 production 기본값으로 자동 채택하지 않는다. 약한100의 실패는 보존한다.
추가로 연결된 말린 천의 한 step과 proxy 세분화의 공간 오차를 확인했다.
시험 최소 간격1mm를 재료 h나 실제 장면의 최종 접촉 두께로 자동 승계하지 않는다.

보장 대상은 선형 proxy다. Bernstein 공간 오차 예산을 명시하면 최소 간격에2배 여유를
더하고 초과 상태를 거절하지만, 이것만으로 인접 고차 곡면까지 포함한 전역 단사성을 증명하지 않는다.
Unweighted IPC의 세분화별 에너지는 표본 수에 의존하므로 접촉 응답의 공간 수렴도 미완료다.
후속 GPU resident 접촉은 다음 절을 따른다. 마찰·실장면 calibration·장기 풍하중·R1 acceptance와 학습 적격성은 남아 있다.
후속 기존 bend500 직사각형/손수건/삼각 깃발의 입력을 보존하는 별도 실행기를 준비했다.
접촉 ON preload부터 다시 계산해 무풍/바람을 분기하며, 원본 GPU hi/lo/current와 다른
CPU FP64/rest 기준 실행임을 명시한다. 세 실제 메시에서 초기 검사와 중력/최대풍의
총6개 짧은 단계 검산을 통과했지만 그 구간 접촉 에너지는0이다. 큰 접힘이나 장기 접촉
완주의 근거로 승계하지 않으며 본 실행은 미시작이다. 이 초기 실행기는 CPU 전용이며
GPU는 후속 별도 실행기로 제공한다. 동결 입력/명령/검사 근거는
\code{experiments/R1_teacher_velocity_reset/self_contact/three_scenes.md}가 소유한다.
구현 계약은 \code{code/docs/p3_self_contact.md}, 설정·전체 수치·source/hash·실패 분모는
\code{experiments/R1_teacher_velocity_reset/self_contact/README.md} 및 연결 report가 소유한다.

\subsection{GPU 전 단계 셀프 접촉 개발 후보 (2026-09-22)}
\label{sec:r1-gpu-self-contact}
CPU 기준과 같은 P3 proxy 및 unweighted IPC barrier를 Warp CUDA로 이식했다.
초기 LBVH 구축·매 질의 refit, VF/EE 거리와 EE mollifier, 정확한 힘/Hessian,
P3 adjoint와 HVP, Newton 이동 CCD 및 quadratic 시간 경로 검사를 GPU에서 수행한다.
고정 topology/초기 구조 준비와 프레임 파일 I/O는 CPU지만 반복 수치 계산·승인·복원에는
CPU IPC를 호출하지 않는다. 실제 거리·미분·CCD는 FP64, 상태는 FP64 hi/lo이고
AABB에만 외향 여유를 둔 FP32를 쓴다. 전체 FP32 또는 M1/M2 접촉 이식 완료가 아니다.

GPU 시간 CCD는 quadratic Bézier control의 상대 속도 상한에 따른 보수 진행이다.
CPU chord/tube Tight Inclusion 구현과 별도로 저장 경로를 대조한다. 후보 버퍼 overflow,
무효 시작점·공간 예산 초과·CCD 반복 한도는 안전 승인으로 바꾸지 않는다.
BVH의 연결은 최초 구축 후 고정하고 경계만 refit하며, 큰 변형에서 주기적 재구축에 따른
속도 개선은 미검증이다. 고정 트리로도 보수 경계를 갱신하므로 후보 누락을 허용하지 않는다.

실제 Newton/Krylov 연산자는 shell+contact다. 비국소 접촉을 shell coloring에 섞지 않고,
current 보조 행렬만 shell-only 근사로 사용한다. 접촉 OFF Gauss/CPU fallback은 없다.
모든 substep을 별도 GPU 객체에서 힘·에너지·위치·핀·국소 기하·시간 CCD로 검산하고,
실패 프레임의 네 raw hi/lo 상태를 GPU에서 시작점으로 복원한다. 조건 분기 본문과
cuDSS 자식을 포함한 그래프에서 host copy/callback0을 확인했다.

초기 v1의 GTX1080Ti 실제 접촉 여섯 사례에서 GPU/CPU 상태와 실패 판정이 일치했다.
계수1,000의 저속 면/엣지8단계씩은 엄격 검산까지 통과했다. 고속12단계의 마지막2단계와
계수10,000의 세 사례 일부는 \code{local_metric} 충분조건 미해결로 거절됐다.
해당 네 사례12구간의 실패 원본은 보존하며 v1 자체를 완주 통과로 승격하지 않는다.
강도를 높이면 간격뿐 아니라 반발·변형도 커지므로10,000을 자동 채택하지 않았다.
GPU 저장 경로의56개 구간은 모두 별도 CPU CCD를 통과했고, 기하 거절 단계도 CPU와 같았다.

사용자 승인 후 v2에 \code{local_metric_refined_v1}을 추가했다. 기존 scalar 충분조건이
실패하면 \(C=F^{\mathsf T}F\)의 Bernstein 행렬75개에서 최소 고유값 하한을 계산한다.
Basis의 비음수성과 partition of unity로 영역 전체의 양의 면적 하한을 확인하고,
미인증 영역만 공간4분할\(\times\)시간2분할한다. GPU 큐와 깊이는 유한하며,
overflow·깊이 한도·비유한 값은 계속 거절한다. 기존 FP64 변환 여유에 행렬 계산·분할
반올림 여유를 더하며 형식적 임의 정밀도 증명이나 연속 ODE 해의 인증으로 주장하지 않는다.
정밀 검사도 같은 수치 표면의 모든 위치·시간에서 비퇴화를 요구한다. 허용 변형률을
높이거나 물성·barrier·dt를 바꾸지 않는다. hi/lo의 큰 공통 원점은 보정 뺄셈으로 제거한다.

같은 여섯 사례56단계의 재계산은 힘·위치·에너지·핀·기하·시간 CCD를 모두 통과했다.
이전 미인증12구간은 추가 행렬 인증 또는 깊이1 분할로 해결됐고,
CPU longdouble 행렬 인증·IPC 시간 CCD와 대조했다. 과거 기하 flags는 별도 저장한다.
관련106개 검사가 skip 없이 통과했으며 실제 공간/시간 붕괴·거의 퇴화·버퍼 한도 거절과
GPU rollback을 포함한다. 상세 수치 여유·큐 계약은
\code{code/docs/refined_metric_certificate.md}가 소유한다.

기존 bend500 직사각형·손수건·삼각 깃발의 GPU 전용 개별 실행 스크립트를 제공한다.
입력 byte/hash·재료·외력·공식 허용오차를 유지하며 새 접촉 ON preload2초에서
raw hi/lo 위치·속도를 보존해 calm/wind 각4초를 독립 분기한다. 세 씬 각각 rest 시작의
중력/최대풍64단계 프레임, 총384단계가 v2의 GPU 검산을 통과했다.
이 구간에는 활성 접촉이 없으며 장기 전체 궤적 완료의 근거로 승계하지 않는다.
RTX5070·긴 접촉 궤적·실장면 calibration·proxy 응답 수렴·곡면 전역 인증 및 R1 acceptance는 남아 있다.
구현 계약은 \code{code/docs/p3_gpu_self_contact.md}, v1 원시 결과·실패 분모는
\code{experiments/R1_teacher_velocity_reset/self_contact/gpu.md}, v2 결과·hash·명령은
\code{experiments/R1_teacher_velocity_reset/self_contact/refined_geometry.md}가 소유한다.

후속 성능 후보 v3는 첫 GMRES 보조 RHS와 채택 trial의 평가를 인스턴스별로 재사용하고,
GMRES 작업 배열의 순차 초기화를 device memset으로 바꾼다. 활성 접촉이 없으면
GPU 분기에서 미분/HVP를 생략하며, 독립 검산에서는 사용하지 않는 Hessian만 제거한다.
힘/에너지와 proxy 오차의 유한성, 후보 overflow, 기하/시간 CCD 및 rollback 조건은 유지한다.
추가 GPU 회귀 검사와 제한 상태 대조는 수치 동작의 검증이며 속도 향상의 확정 근거와 구분한다.
동시 GPU 작업이 확인된 성능 측정은 참고용으로 보존했고, 당시 GPU 단독 실행의
정량 증가율과 최종 가속 배수는 보류했다. 구현 조건과 재측정 명령은
\code{code/docs/p3_gpu_self_contact.md} 및
\code{experiments/R1_teacher_velocity_reset/self_contact/performance.md}를 따른다.
장기 실행이나 R1 acceptance를 완료로 승격하지 않는다.

후속 병렬 후보 v4는 AABB 진단 카운터를 질의별로 합산하고, 접촉 쌍의 공통 미분 항을
GPU 버퍼에 한 번 계산한다. 미분/조립/HVP의 고정 worker는 GPU 후보 수까지만 처리한다.
CCD는128후보 단위의 GPU 작업 큐와 블록별 최소 보폭 합산을 사용하며,
후보 완전성·반복 한도·기하/힘 승인 조건은 유지한다. 추가 버퍼와 launch가 필요하므로
구조 개선만으로 실측 가속을 주장하지 않는다. GPU 회귀와 국소/세 씬 제한 상태 대조를 통과했고,
작은 공유 GPU 측정에서 항상 빨라진 것은 아니므로 초기 단독 성능 채택은 보류했다.
검증 조건·추가 메모리·남은 BVH/atomic 병목과 현재 실행 명령은
\code{experiments/R1_teacher_velocity_reset/self_contact/parallel_v4.md}를 따른다.

이후 사용자가 다른 GPU 작업을 종료했다고 확인한 조건에서 실제 동결v1/v4와 통제된
ON/OFF/재사용/병렬 경로를 새로 측정했다. 세 씬의 rest 시작 중력/최대풍 각1프레임을
예열 후3회 반복한 중앙값에서 v1 대비v4의 계산+독립 검산 시간은16.2--28.2\% 단축됐다.
같은 실행기의 비접촉 첫 프레임에서 접촉 추가 비용은15.7--43.5\%였으며, 과거 공유 부하의
수배 관측 차이를 순수 접촉 비용으로 채택하지 않는다. 마지막v4 병렬화만의 추가 이득은
시간 변동과 비슷하고 증가/감소가 섞여 일관된 가속으로 판정하지 않는다.
기존 승인 기준과 상태 대조를 통과했지만 장기/고밀도 접촉 비용, RTX5070 및 R1 acceptance는
여전히 미완료다. 정확한 표본 분모·클럭/부하 한계·source/hash와 재현 명령은
\code{experiments/R1_teacher_velocity_reset/self_contact/idle_performance.md}가 소유한다.

후속 BVH 후보 v5는 primitive별 트리 순회와 후보별 FP64 거리 검사를 분리한다.
같은 보수 AABB 및 최근접점 수식을 사용하며, 임시 원시 후보 버퍼가 넘치면 GPU 조건 분기에서
기존 탐색 전체를 재실행한다. 임시 버퍼를 새로운 승인 한도로 사용하거나 후보를 잘라내지 않는다.
최종 active/swept 후보 한도, 시간 CCD, 독립 검산과 프레임 복원 기준은 유지한다.
같은 실행기에서 v4 대조와 분리 경로를 예열 후3회씩 순서 교대 측정한 세 씬의 초기 프레임은
계산+검산 중앙값5.0--14.1\% 단축, 여섯 조건 중앙값 합계9.0\% 단축이었다.
호출 수는 유지하며 호출당 탐색 비용을 줄인다. 국소139개 회귀, 동결v5의 접촉6사례56단계와
세 씬384단계 및 기존 상태 대조를 통과했지만 본 시뮬레이션은 실행하지 않았다.
측정 변동, 고밀도 접힘과 임시 후보 재탐색 비용, RTX5070 및 기존 R1 미완료 범위는 남아 있다.
상세 측정 조건·원본/hash·스크립트는
\code{experiments/R1_teacher_velocity_reset/self_contact/broadphase_v5.md}가 소유한다.

사용자가 시작한 v5 직사각형 장기 실행은 preload120프레임과 calm240프레임을 완료했지만,
wind104프레임 승인 뒤 다음 프레임(index104)에서 code99와 flags[0,14]를 기록하고 GPU에서
프레임 시작 상태로 복원한 뒤 종료됐다. Contact/path status는0이었으며 이 결과를 세 씬 장기
완주나 training 적격성 근거로 세지 않는다. 후속 v8 메인 원본은 wind112번째 시도에서 실패했다.
직전 승인 raw hi/lo로 해당 프레임을 재현해60번째 substep의 GMRES code2를 확인했다.
앞59단계 검산은 통과했고 참 선형 잔차가 목표에 미달한 상태에서720회 한도에 도달했다.
미완료 trace의 audit99가 원래 오류를 덮었다. 과거108프레임 외부 종료라는 추정은 정정한다.

v9은 graph 전체 GPU marker를 host의 graph 제출--동기화 wall에 매 프레임 맞춰 내부 solver,
audit과 collision 시간을 자동 정규화하고 raw timer와 배율을 함께 보존했다.
그러나 cycles3에서6으로 늘리는 복구는 해당 실제 실패 프레임을 해결하지 못했다.
후속 v10은 최초 오류와 잔차를 GPU에 보존하고 유한한 code2, contact/path status0,
정상 시간축과 실패 전 prefix의 독립 검산 통과가 확인된 경우에만 GPU 조건 분기로
frame-start hi/lo에서 dt를 절반으로 줄여128단계를 한 번 재실행한다.
같은 held 외력·물성·접촉식·허용오차를 유지하며 선택, 적분, 독립 검산과 채택도 GPU다.
해당 프레임의128단계 검산과 실제 자동 복구를 통과했고 관련 고유 회귀54개가 통과했다.
실패 시도 비용도 전체 시간에 포함하며 다음 프레임은 기본64단계로 복귀한다.
다른 오류나 재실패는 거절하고 Gauss/contact-OFF/CPU fallback은 없다.
이는 단일 실패 프레임의 국소 복구 근거이며 세 씬 장기 완주, RTX5070 복구,
시간 수렴과 training 적격성은 미완료다. 원본 실패 분모·계측 의미·실행 명령은
\code{experiments/R1_teacher_velocity_reset/self_contact/frame_timing_v8.md}와
\code{experiments/R1_teacher_velocity_reset/self_contact/frame112_recovery.md}가 소유한다.

\subsection{현재 개발 상태와 acceptance 구분 (2026-09-11)}

아래 체크박스는 canonical R1 acceptance를 뜻한다. 구현이 존재해도 전체 계약을 아직 통과하지
않았으면 미체크로 둔다. 개발 구현과 국소 fixture의 완료는 다음 표로 별도 기록한다.
각 구현의 근거를 아래 해당 항목에 연결하고, 수식·교체 이유는
\S\ref{sec:r1-contracts}와 \S\ref{sec:r1-decisions}, 결과는 \S\ref{sec:r1-fixtures},
원본·hash·재현 명령은 \S\ref{sec:r1-artifacts}에 함께 보존한다.

\begin{longtable}{L{0.23\linewidth}L{0.20\linewidth}L{0.47\linewidth}}
\toprule
항목 & 개발 상태 & 실제 근거와 남은 경계\\
\midrule
\endfirsthead
\toprule 항목 & 개발 상태 & 실제 근거와 남은 경계\\ \midrule
\endhead
Registry / trajectory / sequence & 구현·개발 검증 완료 & Native SI/hash, force/work와 실패 prefix, CPU replay 및 사용자 GPU 12/12. 전체 R0 동결을 뜻하지 않음.\\
선형 teacher probe & 구현·개발 검증 완료 & Nonnegative barycentric forward/adjoint와 원본 대조. Independent-GS intersection과 고차 public map은 미완료.\\
Native/area hinge & 채택 근거 실패 & Mesh 의존성 및 방향 편향. 결과를 유지하고 isotropic material mapping으로 승격하지 않음.\\
Quadratic patch shell & 제한된 수치 검증 통과 & Objectivity/미분 및 가속도 Newmark 시간 진단 통과. Checkerboard interior force와 공간 속도 실패.\\
P3/독립 spline & 선형 압력·작은 굽힘 fixture 통과 & 각 조건의 1\% 공간/방향/기준 대조. \(x^2\) 초기 변위의 속도 실패 유지. 비선형 shell의 독립 기준으로 자동 승계하지 않음.\\
P3 유한 회전 shell & 수식·CPU/GPU 구현 및 개발 검증 & 막·굽힘 응력, 요소 연결, consistent mass, Newmark와 공력 구현. 약한 바람 통과;4배 공간 미달·10초 풀이 실패/시간 한도는 아래 참조.\\
Teacher 실행 개선 & FP64 hi/lo GPU 상주 개발 경로 & Newton/GMRES·보조 풀이·독립 검산 GPU화 및 후속 최적화. 과거 CPU+HVP 비교는 당시 근거로 보존.\\
P3 셀프 접촉 & CPU/GPU 제한 개발 검증 & GPU 전 단계·프레임 복원·세 씬384단계 검산. 정밀 기하 인증으로 과거12구간 미인증 해결·실제 퇴화 거절. 마찰/장기 접촉/곡면 전역 인증·공간 수렴 미완료.\\
Development samples & 생성·검증 완료 & Newton 3개 source/15 window, 원본·replay·loader 검증. Training eligibility false.\\
Canonical R1/Oracle & 미완료 & 동일 backend의 최종 수렴, 고차 map/traction, spectrum/tip/work, 독립 GS 2개와 oracle/transport 필요.\\
\bottomrule
\end{longtable}

\subsection{TeacherPhysicsRegistry와 reference 생성}

면적·질량은 \S\ref{sec:r1-mass}의 식~\eqref{eq:r1-consistent-mass}를 따른다.

\begin{itemize}
  \item[$\square$] \code{TeacherPhysicsRegistry}에 structure/element/constitutive law, equality-constrained thickness/density,
  elastic parameter, damping, attachment enforcement, nonlinear tolerance와 unit을 등록한다.
  \item[$\square$] Solver/integrator, timestep, nonlinear iteration, linear solver/tolerance, precision, contact-off flag,
  quadrature와 output sampling rule을 version/hash한다.
  \item[$\square$] Material preset 이름을 실제 SI parameter tuple로 해석하는 versioned mapping을 저장한다.
  \item[$\square$] Teacher와 GS의 traction law/effective-scale/guard/sign/sample identity와
  zero guard-activation을 비교하는 validator를 만든다.
  % [논문 작성 참고 -- 수정 이유]
  % 서로 다른 rho/C 값을 맞추는 검사가 아니라 실제 힘을 결정하는 단일 scale과 guard 비활성 상태를 검증한다.
  \item[$\square$] Mesh surface quadrature가 \(\sum a_T=A_{\mathrm{ref}}\),
  \(\sum m_T=M_{\mathrm{ref}}\)를 만족하고 force integral, sign과 external work analytic fixture를 통과한다.
  여기서 \(m_T\)는 quadrature mass이며, consistent mass의 DOF row sum을 양의 point mass로 해석하지 않는다.
  \item[$\square$] Finest fixed timestep의 spatial series와 finest accepted mesh의 temporal series를 별도 command/report로 실행한다.
  \item[$\square$] Accepted mesh/timestep은 결과, tolerance, convergence order diagnostic와 hash로 고정한다.
\end{itemize}

\subsubsection{현재 Registry 구현과 적용 범위}

Pure schema와 Newton adapter를 분리했다. Native material parameter의 SI 해석, rest mesh와
source-object identity, material/solver/traction 설정, realized model 검사를 제공한다.
Registry 생성 자체가 simulation이나 파일 저장을 수행하지 않는다.
\code{validate_against_simulation()} 및 \code{require_valid()}로 실제 mass, pin,
material, rest geometry, gravity, solver 설정과 finite 상태를 검증한다.
Native bending의 단위 N을 연속 판 강성 N\,m로 바꿔 표기하지 않는다.
Native membrane law는 \code{newton_stable_neo_hookean_membrane_v1}이다.
\code{tri_ke}/\code{tri_ka}는 N/m, \code{tri_kd}/\code{edge_kd}는 s,
\code{edge_ke}는 N이다. 이 native tuple을 continuum \(E,\nu,h,\rho\)로 환산하는 경로는 구현하지 않았다.
후속 CPU StVK shell 및 선형 P3 판과 구성식·물성 identity를 구분한다.

Strict JSON은 unknown/missing field, duplicate key, non-finite, unit/shape/version 오류를 거부한다.
Float 표현과 signed zero를 정규화하며 배열 hash는 little-endian C-order의 dtype/shape/unit/content를
포함한다. Effective registry hash, Teacher/GS 공용 traction identity hash, requested/run provenance와
실제 material array/solve 경로의 realized hash를 구분한다.
비활성 명목 설정과 시간 변화 wind speed를 모두 material identity로 섞지 않는다.
Source membership validator는 외부에서 검증한 split ref를 받으며 split 생성·배정·봉인을 대신하지 않는다.
Registry는 teacher-only이고 target package serializer를 새로 만든 것이 아니다.

초기 실제 model 대조는 \(\mathrm{rtol}=5\times10^{-6}\), mass/area absolute \(10^{-8}\),
pin \(10^{-6}\) m를 썼다. 이는 물리 수렴 tolerance가 아니다.
CPU/GPU solve branch와 환경을 기록하며 registry hash 동일성이 장치 간 bitwise trajectory 동일성을
보장하지 않는다. Core import는 Newton/Warp 없이 가능하도록 optional dependency 경계를 검사했다.

\subsubsection{Trajectory, sequence와 초기 상태 구현}

Trajectory는 \(T\) interval과 \(T+1\) state를 구분하고 frame zero, timestamp, position/displacement,
velocity, held aerodynamic force와 interval work를 저장한다. Chunk inventory와 checksum으로 prefix를
검증하고 중간 실패의 완료 prefix, 오류와 미실행 부분을 보존한다. Reader의 무결성 검증과 실제 solver
replay는 별개다. 저장값만 다시 읽어 replay 통과로 표시하지 않는다.

Wind sequence runner는 초 단위 steady/pulse/chirp 등 프로그램을 case별 독립 초기 상태로 실행하고
성공·실패·미실행 inventory를 남긴다. Displaced free-decay 지원은 authored rest를 바꾸지 않고
초기 변위를 적용하며 requested 변위와 float32 realized state를 구분한다. Trajectory v2로 확장하고
v1 읽기 호환을 유지했다. Zero ambient는 상대 속도 drag가 남고, aero-off만 공력을 제거한다.
Material damping off 여부는 또 별도 설정이다.
현재 native 방향은 setup-fixed \code{wind_direction}이며 registry identity에 포함된다.
기존 program은 scalar wind speed의 steady/zero/pulse/chirp를 제어하고,
\code{seed}는 metadata로 저장될 뿐 파형이나 방향을 랜덤화하지 않는다.
시간 변화 wind vector와 재현 가능한 random program/state의 연결은 다음 탐색을 위한 후속 구현이다.

\subsubsection{물리 수식 우선 보완: P3의 3D 요소 내부 연산}
\label{sec:r1-p3-surface-equations}

사용자 지시에 따라 추가 학습데이터를 생성하지 않고, 기존 P3 shape의 3D 기하와 가상 일을
\code{teacher/p3_surface.py}에 분리했다. Flat rest material 좌표 \((s,t)\) [m]에서
10개 signed shape \(N_i\)와 현재 3D DOF \(x_i\)에 대해
\begin{equation}
\begin{aligned}
 r&=\sum_iN_i x_i,& F_a&=r_{,a},& H_{ab}&=r_{,ab},\\
 c&=F_s\times F_t,& J&=\|c\|,& n&=c/J,\\
 \epsilon_{ab}&=\tfrac12(F_a\cdot F_b-\delta_{ab}),&&
 b_{ab}=n\cdot H_{ab}.
\end{aligned}
\label{eq:r1-p3-surface-geometry}
\end{equation}
\(F,J\)는 무차원, \(H,b\)는 1/m다. Reference 곡률은 0이고,
engineering vector는 \(e=(\epsilon_{ss},\epsilon_{tt},2\epsilon_{st})\),
\(b=(b_{ss},b_{tt},2b_{st})\)다. 곡률의 정의는 현재 법선을 사용한다.
% [구현 보완 이유 -- 2026-09-09]
% 기존 normal-only graph는 최종 finite rotation 상태를 표현하지 못한다.
% 아래 기하/가상 일 검사는 공통 요소 연산이며 nonlinear constitutive/edge law 채택이 아니다.
기하와 가상 일은 Verhelst 등 (2024)의 Kirchhoff--Love shell analysis \S2.1--2.3과
대조했으며 원문 링크는 아래 실험 보고서에 보존했다. 두께 방향 strain 전개의 곡률 부호와
여기서 정의한 \(b=n\cdot H\)의 부호를 구분한다.

단위 법선의 위치 의존성을 포함하여
\begin{equation}
 \delta c=\delta F_s\times F_t+F_s\times\delta F_t,\qquad
 \delta n=\frac{(I-nn^\top)\delta c}{J},\qquad
 \delta b_{ab}=\delta n\cdot H_{ab}+n\cdot\delta H_{ab}
 \label{eq:r1-p3-surface-variation}
\end{equation}
를 계산한다. \(c=nJ\)의 혼합 2차 미분에서
\(\delta_1n\,\delta_2J\), \(\delta_2n\,\delta_1J\),
\(n\,\delta_1\delta_2J\)를 모두 포함하여 기하 강성에 필요한 방향 미분을 제공한다.
현재 API의 입력 resultant \(S\) [N/m]와 \(B\) [N]는 각각 \(e,b\)의 공액이다.
\begin{equation}
 \delta E_{\rm elem}=\int_{A_0}(S\cdot\delta e+B\cdot\delta b)\,dA_0,
 \qquad f_{\rm elem}=-\operatorname{adjoint}_x(\delta E_{\rm elem}).
 \label{eq:r1-p3-element-virtual-work}
\end{equation}
Shear resultant에 2를 다시 곱하지 않는다. \code{stress_force}는 주어진 resultant의 음의
adjoint이며 재료 응력을 생성하거나 전역 저장 에너지·접선·solver를 완성하는 함수가 아니다.
양의 rest quadrature의 consistent mass는 식~\eqref{eq:r1-consistent-mass}를 그대로 따른다.

현재 표면에서 \(v_q=(Nv)_q\), \(v_{\rm rel}=v_{\rm air}-v_q\)로 두어
\begin{equation}
 \tau_q=\kappa_{\rm aero}|v_{\rm rel}\cdot n_q|(v_{\rm rel}\cdot n_q)n_q,\quad
 f=N^\top\operatorname{diag}(a_q^0J_q)\tau,\quad
 f\cdot v=\sum_q a_q^0J_q\tau_q\cdot v_q.
 \label{eq:r1-p3-surface-aero}
\end{equation}
현재 면적은 한 번만 곱하고 vector-norm guard는 그 전에 검사한다.
Zero ambient에서 일률은
\(-\kappa_{\rm aero}\sum_q a_q^0J_q|v_q\cdot n_q|^3\le0\)다.
Aero-off만 공력을 0으로 한다. 힘을 합쳤을 때 합력과 모멘트, virtual work가 모두 일치한다.

\textbf{검증 근거:} 신규 10개 검사는 독립 6차 적분, 비직각 요소의 3차 재현,
큰 기울기의 해석적 strain/curvature, 1.7 rad 강체 회전, 1·2차 미분과 대칭,
응력 및 공력의 가상 일·합력·모멘트, 정지 공기의 수동성과 guard 거부를 통과했다.
별도 n4/n8 및 두 대각선의 4개 정적 대조 (총240개 요소)에서 기존 P3와
질량/normal 힘/합력/일률의 최대 상대 차이는 \(1.765\times10^{-15}\) 미만이다.
명령·수치·환경·source hash는
\code{experiments/R1_teacher_velocity_reset/p3_surface/README.md}와
\code{evidence/report.json}에 보존했다. 새 학습 sample은 0개다.

\textbf{후속 결정:} 2026-09-09 사용자가 \textbf{1번 유한 회전 shell}을 선택했다.
아래 전역 구조식은 그 결정에 따라 구현했다. 선행 요소 연산의 원본 검사는 그대로 보존한다.
선행 geometry guard는 구적점의 비퇴화만 검사한다. 아래 전역 shell의 후속 검산에서는
Bernstein 상한을 추가해 전체 요소·수치 시간 보간의 injectivity 충분조건을 별도로 확인한다.

\subsubsection{선택된 P3 유한 회전 shell의 전역 구조식}
\label{sec:r1-p3-shell-equations}

\paragraph{재료·기하 범위와 저장 에너지.}
Flat-rest reference의 \(E=10^6\) Pa, \(\nu=0.3\), \(h=0.01\) m,
면밀도0.1 kg/m$^2$를 유지한다. 자유 영역은 \([0.25,1]\times[0,1]\) m이고
고정 영역의 질량0.025 kg와 자유 영역의 질량0.075 kg를 구분한다.
유한 회전은 허용하되 작은 재료 strain을 전제하는 고정 두께 Koiter/StVK shell이다.
임의의 큰 신장, 소성, 두께 변화, 접촉/self-contact를 포함하지 않는다.
구조 감쇠는 선행 P3와 같은0이며, 정지 공기에서도 상대풍 drag는 유지한다.
Engineering strain/curvature는 식~\eqref{eq:r1-p3-surface-geometry}를 사용하고
\begin{equation}
 \begin{gathered}
 C_\nu=\begin{bmatrix}1&\nu&0\\\nu&1&0\\0&0&(1-\nu)/2\end{bmatrix},\qquad
 D_m=\frac{Eh}{1-\nu^2}C_\nu,\quad
 D_b=\frac{Eh^3}{12(1-\nu^2)}C_\nu,\\
 E_{\rm vol}=\frac12\sum_K\int_K(e^\top D_m e+b^\top D_b b)\,dA_0
 \end{gathered}
 \label{eq:r1-p3-shell-volume}
\end{equation}
로 둔다. \(S=D_me\) [N/m], \(B=D_bb\) [N]이고 굽힘 moment tensor는
\(m=\left[\begin{smallmatrix}B_1&B_3\\B_3&B_2\end{smallmatrix}\right]\)다.
Rest에는 곡률이 없으며 현재 unit normal의 모든 위치 미분을 포함한다.

\paragraph{요소 간 연결과 clamped boundary.}
P3는 위치가 연속인 \(C^0\) 함수이므로 법선 연결의 consistency와 stabilization이 필요하다.
내부 edge의 \(\mu\)는 plus 요소의 rest 외향 단위 conormal이며, 양쪽 flux에 같은 \(\mu\)를 쓴다.
\begin{equation}
 R=n^+-n^-,\qquad
 q=\tfrac12\left(F_a^+m_{ab}^+\mu_b+F_a^-m_{ab}^-\mu_b\right),\qquad
 E_e=\int_e\left(R\cdot q+\frac{\alpha}{2\bar h_e}|R|^2\right)\,ds_0,
 \quad \alpha=2.25Eh^3.
 \label{eq:r1-p3-shell-edge}
\end{equation}
\(\bar h_e\)는 인접 요소 최대 edge 길이의 평균이다. Penalty 계수는 선행 P3 값을 유지했으며
결과를 통과시키기 위해 조정하지 않았다. 고정 경계는 \(R=n-n_0\),
\(q=F_am_{ab}\mu_b\)인 \emph{full flux}를 사용한다. 고정 위치 trace는 정확히 0으로 두고
회전 조건은 위 에너지로 약하게 부과하므로 boundary normal jump도 별도로 보고한다.
연속 법선의 극한에서 \(R=0\)이고
\(n\cdot\delta F_a=-F_a\cdot\delta n\)이므로 첫 항의 변분은 요소 굽힘 적분의 경계 moment 항을 상쇄한다.
Plus/minus 교환과 conormal 반전을 함께 하면 에너지·힘이 동일하다.
% [교체 이유 -- 2026-09-09]
% 기존 scalar slope jump만 current normal jump로 바꾸면 flux의 F/m/n 미분과 경계 full flux가 누락된다.
% 이 에너지 전체에서 gradient/HVP를 유도하며 native hinge의 격자별 계수 보정을 사용하지 않는다.
\href{https://orbi.uliege.be/handle/2268/2093}{Noels (2009), nonlinear DG Kirchhoff--Love shells}
\S3의 normal jump와 consistent moment flux 원리와 대조했다. 이 구현은 위 고정 두께 에너지의 변분이다.
해당 논문의 두께 적분 hyperelastic 재료식이나 nonlinear 안정성 보장을 그대로 구현·승계한 것으로 주장하지 않는다.

\paragraph{힘·접선·지지대 모멘트.}
전체 \(E=E_{\rm vol}+\sum_eE_e\)에서 \(f_{\rm int}=-\nabla_uE\),
\(Hv=D(\nabla_uE)[v]\)를 계산한다. 구현은 해석적 gradient에 forward 방향 미분을 적용하며,
유한 차분은 독립 검사용이다. Flux의 current tangent, moment와 양쪽 unit normal을 모두 미분한다.
CPU와 GPU가 공유하는 구적점 미분을 명시하면 다음과 같다.
\(c=F_1\times F_2\), \(J=\|c\|\), \(P_n=I-nn^\top\)에 대해
\[
 \delta c=\delta F_1\times F_2+F_1\times\delta F_2,\qquad
 \delta J=n\cdot\delta c,\qquad \delta n=P_n\delta c/J.
\]
\(H=(x_{,11},x_{,22},2x_{,12})\), \(t=\sum_i B_iH_i\), \(z=P_nt/J\)로 두면
체적 에너지 밀도의 gradient는
\[
 A_1=S_1F_1+S_3F_2+F_2\times z,\quad
 A_2=S_2F_2+S_3F_1+z\times F_1,\qquad C_i=B_i n
\]
이다. 여기서 \(A_a=\partial\psi/\partial F_a\), \(C_i=\partial\psi/\partial H_i\)다.
Edge에서는 면 수 \(N_e=2\)(내부) 또는1(고정 경계), \(p_e=\alpha/\bar h_e\),
\(\sigma_+=1,\sigma_-=-1\)로 두고 각 면에서
\[
 \begin{aligned}
 g_B&=\frac1{N_e}
 \begin{bmatrix}\mu_1(F_1\cdot R)\\\mu_2(F_2\cdot R)\\
 \mu_2(F_1\cdot R)+\mu_1(F_2\cdot R)\end{bmatrix},\qquad
 g_b=D_b^\top g_B,\\
 t_e&=\sigma(q+p_eR)+\sum_i(g_b)_iH_i,\qquad z_e=P_nt_e/J,\\
 A_1^e&=\frac{\mu_1B_1+\mu_2B_3}{N_e}R+F_2\times z_e,\\
 A_2^e&=\frac{\mu_1B_3+\mu_2B_2}{N_e}R+z_e\times F_1,\qquad
 C_i^e=(g_b)_i n .
 \end{aligned}
\]
Rest 형상함수 미분과 이 gradient를 적분하여 nodal force를 조립한다.
HVP는 위 식의 모든 값에 \((\text{value},\text{directional derivative})\)를 함께 전파해 얻는다.
이 전파에는 normal뿐 아니라 flux의 tangent와 moment 미분도 포함된다.
고정 법선이 받는 gradient와 shell에 가해지는 회전 지지 모멘트는
\begin{equation}
 g_{n_0}=-(q+\alpha R/\bar h_e),\qquad
 \tau_{\rm frame\to shell}=\int_{\Gamma_D}n_0\times g_{n_0}\,ds_0.
 \label{eq:r1-p3-shell-support-torque}
\end{equation}
고정점 반력의 \(\sum x_i\times r_i\)에 이 항을 더해야 전체 지지 모멘트가 된다.
고정 지지의 속도·각속도는0이므로 지지대의 일은0이다.
실제 조립에서 내부 합력0, 자유 shell의 합모멘트0, 고정 shell의
\(\sum x_i\times f_{{\rm int},i}=\tau_{\rm frame\to shell}\)를 검사한다.

\paragraph{동일 물리식의 시간 전진과 실패 보존.}
Consistent \(M\otimes I_3\)의 모든 항을 유지하며, free block에서
\begin{equation}
 Ma_0=f_{\rm hold}+f_{\rm int}(u_0),\quad
 u_1=u_0+\Delta t\,v_0+\frac{\Delta t^2}{4}(a_0+a_1),\quad
 v_1=v_0+\frac{\Delta t}{2}(a_0+a_1),\quad
 Ma_1-f_{\rm int}(u_1)-f_{\rm hold}=0
 \label{eq:r1-p3-shell-newmark}
\end{equation}
을 푼다. Frame 시작의 현재 표면/속도로 상대풍을 계산한 뒤60 Hz frame 동안 nodal force를 고정한다.
하중 또는 reset이 바뀐 경계에서는 새 힘과 상태로 \(a_0\)를 다시 계산한다.
State는 절대 rest 좌표를 반복 차분하지 않는 rest-relative \(u,v,t\)다.
Newton의 정확한 HVP와 \(M+\Delta t^2K_{\rm rest}/4\) preconditioner를 쓰며 GMRES의 실제 residual을 검사한다.
Force tolerance는 \(10^{-10}\) N\(+10^{-9}\) scale, 변위 correction tolerance는
\(10^{-12}\) m\(+10^{-9}\|u\|\)다. Line search는 residual 감소를 요구한다.
실패한 step은 입력을 변경하지 않고 시도 내역을 보존하며, 실행기는 완료한 prefix와 reset 양쪽 상태를 저장한다.
Reset은 \(u,t\)를 보존하고 \(v=0\), 제거 에너지는 \(v^\top Mv/2\)다.
\(\Delta(T+E)-f_{\rm hold}\cdot\Delta u\)는 별도 적분 오차로 기록하며,
비선형 Newmark가 에너지를 정확히 보존한다고 주장하지 않는다.

\paragraph{수식 검사와 해결 과정.}
Kernel/조립/solver의15개 검사는 energy gradient, 정확한 HVP와 대칭,
강체 objectivity, 질량·선형 판 극한, rigid nullspace, 지지 모멘트,
signed P3의3차 재현·force adjoint, 상대풍 수동성, reset 및 실패 상태 보존을 통과했다.
첫 rest-HVP 대조에서 고정 반력 행33개가 불일치했다. 선형 preconditioner의 boundary가
normal-slope 성분만 포함했던 원인이며 full-normal의 두 접선 성분을 포함하도록 수정했다.
수정 후 전체 HVP--sparse K 일치와 허용 free block의 기존 P3 K 일치를 함께 검사했다.
작은 dt의 rest 좌표 상쇄는 solver tolerance 완화 대신 변위 상태와 rest tangent 분리로 줄였다.
원본 검산5개와 Bernstein 상한3개를 포함한 최종 신규23개 검사는 모두 통과했다.
원본의 support torque를 운동방정식에서 재계산하고, manifest를 다시 봉인한 잘못된 모멘트도 거부했다.

정적 원통은 3개 mesh, 2개 대각선, 2개 곡률(1.2/2.0 m$^{-1}$), 3개 회전축의 36개 조합이다.
해석적 isometry 에너지 \(E_*=(D/2)\kappa^2\times0.75\)에 대한 최악 상대 차이는
n4/n8/n16에서 각각 0.628545\%, 0.0642517\%, 0.00747187\%다.
Finest 인공 membrane 에너지 비율은 \(1.082\times10^{-8}\) 미만이다.
이는 주어진 형상의 정적 에너지 수렴이며 외력에 의한 운동의 공간 수렴은 아니다.

독립 DOP853 두 해상도와 Newmark32/64/128 step을 n4의0.0005 s 구간에서 대조했다.
Rest와 최대 회전0.9 rad의 굽은 초기 형상 두 계열이며, 후자는 초기 형상을 지지하는
manufactured preload에 작은 held wind를 더한 조건이다. 큰 초기 변위를 분모로 쓰지 않고
\(u(t)-u(0)\)의 증분을 비교했다. DOP853 자체 최악 차이는 \(6.926\times10^{-10}\),
Newmark finest의 rest/bent 속도 차이는 각각 \(4.078\times10^{-7}\),
\(2.995\times10^{-6}\)로1\% 기준 아래다. 시간 간격을 절반으로 줄일 때 오차가 약1/4로 감소했다.
이 검사는 같은 nonlinear ODE의 시간 적분 정확도이며 독립 공간 reference는 아니다.

\paragraph{방향 변화 바람과 큰 굽힘의 추가 진단.}
같은60 Hz force clock의 실제3 frame 바람에서 n4/n8의 sub64$\to$128 속도 차이는
각각0.275331\%/0.471643\%였다. n16의 sub128$\to$256은0.224287\%로1\% 아래다.
공간 n8$\to$16 at sub128의1.357196\% 실패 이후, n16$\to$32에서는0.385696\%로 줄었다.
n16의 forward/backward 대각선 비교도0.153774\%로1\% 아래다.
모든 오차는 같은 material-area norm의 공통 저장 시각 기준이다. 연속 시간 상한으로 소개하지 않는다.

최대 회전0.9 rad, 초기 nodal 변위0.329974 m의 원통을 지지 하중 없이 놓는 release에서는
n4 sub128$\to$256의 속도 차이13.313118\%, sub256$\to$1024의7.737815\%를 보존한다.
Sub1024$\to$2048에서0.762185\%로 감소하여 \textbf{n4 시간 진단}을 통과했다.
Finest 변위 증분 차이는0.00455078\%다. 구적6/4 대10/8의 같은4개 상태 force mass-dual 차이는
0.00682105\% 이하였다. 큰 시간 오차를 구적 부족만으로 설명하지 않는다.

초기 접선 기저에서 sub256/1024 속도 차이의 제곱 mass norm 중85.8568\%가500--1000 Hz 대역에
모였다. 초기 접선 대칭 오차는 \(1.522\times10^{-16}\), 음의 고유값은0개다.
고정 선형계의 \(\omega_N=2\arctan(\omega\Delta t/2)/\Delta t\) 관계라면876.46 Hz 성분의
0.05 s 위상 지연은 sub256에서2.894 rad, sub1024에서0.184 rad다. 이는 \textbf{초기 접선의 설명용 진단}이며,
실제 비선형 시스템을 고정 모드로 전진시키거나 canonical spectrum을 완료 처리한 근거가 아니다.

원통 release의 n4$\to$8 at sub256은 속도85.2503\%, 변위 증분14.3522\%의 차이가 남았다.
시간 오차가 충분히 제거되지 않은 이 결과를 공간 수렴의 통과/독립 기준으로 사용하지 않는다.
주어진 원통의 free tip/side moment resultant는0.1098901 N/0.0329670 N이므로 자유단 moment=0과
호환되지 않는다. 지지 하중을 제거한 초기 가속도 RMS는 n4/n8/n16에서
368.481/1005.274/2793.453 m/s$^2$로 증가하지만, rest+0.5 m/s wind에서는
1.28553/1.29226/1.29564 m/s$^2$다. 이 값은 고정 영역의 영 가속도까지 포함한 전체1 m$^2$ RMS다.
자유0.75 m$^2$만의 RMS로 표시하면 각각 \(1/\sqrt{0.75}\)배다.
이 강한 경계 과도응답을 바람으로 도달한 상태와 구분한다.
% [정규화 영역 명시 -- 2026-09-09]
% 처음에는 면적 평균이 빠진 단위 오류로 해석했으나 선행 README의 /1m^2 정의를 확인했다.
% 기존 전체 면적 RMS는 유효하다. 새 자유 면적 RMS와 영역을 구분하고 상대 판정은 유지한다.
% [진단에 따른 범위 구분 -- 2026-09-09]
% 임의 analytic 초기 원통은 자유단 조건을 만족하는 wind-reached 상태가 아니다.
% 실패를 삭제하지 않고 별도 stress 결과로 보존하며, 사용자가 고른 rest-start/reset 절차를 따른다.

원래 요청한 rest-start/velocity-reset 경로의 더 큰 변형은 별도5 m/s peak의 수식 진단에서 확인한다.
3 frame 바람 \([0,5,0]\to[2.5,3.5,-2.5]\to[0,0,0]\) m/s,
frame2의 reset에서 n4/sub128은 변위0.0727356 m, 법선 회전0.254037 rad,
strain 성분0.00622455, 최소 area ratio0.995873을 관측했다.
Shape/time 보존, 새 상대풍과 work/reset 장부를 통과했으며 시간·공간 수렴은 별도로 판정한다.
N16의 reset 도달 상태에서는 평면 내 변위 최대7.23639 mm가 함께 발생했다.
같은 상태에서 평면 내 변위만 제거하고 위 유한 회전 에너지를 다시 평가하면,
막 에너지는 \(7.06081\times10^{-5}\) J에서0.902094 J로 증가한다.
Engineering strain 성분의 구적점 최대값도0.00105913에서0.0337355로 변한다.
이는 휘면서 면내 좌표도 변하는 결합의 필요성을 보이는 설명용 대조다.
선형 P3의 재시뮬레이션 결과 또는 별도 수렴 판정으로 소개하지 않는다.
이 조건에서 n4와 n8의 sub128$\to$256 속도 차이는 각각0.103571\%, 0.304620\%다.
공간 n4$\to$8의7.312009\%는 n8$\to$16에서1.849529\%로 줄었지만 아직1\% 위다.
N16의 시간 비교는0.236778\%, 대각선 방향 비교는0.160156\%로 통과했다.
후속 n16$\to$32/sub128 공간 속도 차이는0.391794\%로 통과했다.
Finest 시간 간격과 방향의 수치 보간 검산은 아래 GPU 절의 표에 정리한다.
풍속5 m/s를 데이터 입력 범위로 동결하거나 원통 release 실패를 이 결과로 대체하지 않는다.
기존0.5 m/s n8/sub128 reset의 독립 재실행은17개 배열/1,109,544 scalar와
step diagnostics가 정확히 일치했다. 제거 kinetic은 \(1.86857643438\times10^{-5}\) J다.

\paragraph{구적점 밖의 기하·strain 상한.}
구적점의 양의 면적 검사만으로 요소 사이 또는 저장 시각 사이의 fold를 배제할 수 없다.
이를 보완해 각 P3 요소의 변위 gradient를 공간2차 Bernstein 계수로 변환하고,
각 Newmark 구간의 시간2차 control을
\(u_0,\;u_0+\Delta t\,v_0/2,\;u_1\)로 둔다.
이는 식~\eqref{eq:r1-p3-shell-newmark}의 평균 가속도에 해당하는 수치 보간이다.
공간/시간의 양의 Bernstein product를 적용하면 Green strain은 공간4차/시간4차 계수의
최대 절댓값으로 상한을 얻는다. Reference 평면으로 투영한 변위 gradient의
Frobenius norm 상한을 \(r\)라 할 때, convex 사각형에서
\[
 r<1\ \Longrightarrow\
 \|\pi x(X)-\pi x(Y)\|\ge(1-r)\|X-Y\|,\qquad
 J\ge(1-r)^2>0 .
\]
고정부의 영 변위 연장까지 연속이므로 이 충분조건은 해당 수치 표면의 전역 단사성을 보장한다.
큰 회전 중 이 조건을 만족하지 않는 경우가 곧 self-intersection이라는 뜻은 아니다.
곱셈·미분 계수의 float64 누적 여유를 명시적으로 더하며, 정식 interval arithmetic 인증이나
정확한 연속 ODE 해의 상한으로 주장하지 않는다.

Hessian의 공간1차/시간2차 control의 vector norm으로 engineering curvature 상한도 얻는다.
따라서 선형 두께 보간의 표면 strain은
\(\|e\|_\infty+(h/2)\|b\|_\infty\)로 상한을 둔다.
이는 두께 방향 고차항까지 포함한3D hyperelastic strain 검증이 아니다.
5 m/s reset의 n16/sub128에서는 \(r\le0.034626\), \(J\ge0.931948\),
mid-surface strain 성분 상한0.00113629, 선형 표면 strain 성분 상한0.0170444였다.
전체 요소·구간에서 위 충분조건을 만족했다. 새 재료 허용 strain 기준을 이 관측값에 맞춰 정하지 않는다.

실제3 frame 방향 변화 바람과 지지 하중 없이 굽은 형상을 놓는 후속 진단의 raw/config/source hash,
수치와 실패 여부는 \code{experiments/R1_teacher_velocity_reset/p3_shell/README.md}에서 함께 관리한다.
완료된 실행의 \code{status=completed}와 물리 acceptance를 구분한다.
\textbf{채택 경계:} 고차 public registry/map의 직렬화와 common-valid probe, 충분한 기간의
동적 공간/시간/방향·독립 기준, spectrum/FRF, oracle/GS 및 장기 안정성은 별도 조건이다.
위 수식 검사로 R1을 완료 처리하거나 학습데이터 발행을 재개하지 않는다.

\paragraph{응답 차이의 수치 보간 시간 상한.}
RMS의 정규화 영역을 명시한다. 선행 단기 보고서는 고정 영역의 영 응답을 포함해
\(\|f\|_{\rm whole}^2=(1\,\mathrm{m}^2)^{-1}\int_{A_{\rm free}}\|f\|^2\,dA_0\)를 사용했다.
새 장시간 비교 law v2는 자유 면적 평균
\(\|f\|_{A_0}^2=(0.75\,\mathrm{m}^2)^{-1}\int_{A_{\rm free}}\|f\|^2\,dA_0\)를 보고한다.
둘 모두 변위 m/속도 m/s의 RMS다. 같은 영역의 numerator/denominator에 공통 인자가 적용돼
상대 오차와1\% 판정은 유지되며, 아래 상한 식은 두 정규화 모두에 성립한다.
원본 배열과 solver 허용오차는 변경하지 않았다. 두 절대값의 변환과 초기 해석 정정은
\code{p3_shell/evidence/metric_units_v6.json}에 보존했다.
공통 저장 시각만의 비교에서 누락될 수 있는 fine 시각과 reset 직전 속도를 함께 비교한다.
각 구간의 변위는 위2차 Newmark 보간, 속도는 저장 양 끝점의 선형 보간으로 정의한다.
Reset 시각의 좌극한은 별도 저장한 pre-reset 속도이며 우극한은 두 run 모두0이다.
공통 fine 구간에서 속도 차이는 선형이므로 면적 norm의 최대는 끝점에 있다.
저장 속도와 변위 보간의 미분 사이 roundoff 결함을
\(\eta=\max\|2(u_1-u_0)/\Delta t-v_0-v_1\|_{A_0}\)로 계산하면,
\[
 E_u^{\mathrm{all}}\le E_u^{\mathrm{end}}+
 \frac{\Delta t_{\mathrm{fine}}}{2}
 \bigl(E_v^{\mathrm{end}}+\eta_a+\eta_b\bigr).
\]
변위 증분에도 같은 식을 적용하고, 기준 응답의 끝점 최대값을 분모 하한으로 사용한다.
이 상한은 저장 수치 궤적의 시간 보간을 비교하며 정확한 연속 ODE 해에 대한 인증은 아니다.
별도의 raw 검산에서 힘/Newmark/work/reset 식을 재계산하고 각 원본 manifest에 연결한다.
약한 바람의 n32/sub128$\to$256에서는 속도 상한0.231300\%, 변위 상한0.002233\%였다.

\paragraph{동일 수식의 GPU 실행.}
사용자는 장시간 검증의 비용을 줄이기 위해 GPU 이식을 우선하도록 선택했다.
첫 이식은 Warp1.17.0 float64의 요소/edge 기하·구성식·힘·정확한 방향 HVP,
요소 적분·고정 순서 nodal gather와 현재 면적 공력이다.
Atomic force 합산, fast math 또는 생산 HVP의 유한 차분을 사용하지 않는다.
Newton/GMRES 제어, consistent-mass 희소 풀이와 rest-K preconditioner,
마지막 총에너지 합산은 CPU에 유지한다. 전체 solver의 GPU 상주를 주장하지 않는다.

GTX1080Ti에서 신규5개 검사와 n4/8/16/32의 rest/원통/부드러운 변형12개 연산자 대조를 통과했다.
N16/sub128, 5 m/s peak의3 frame/reset 궤적은 CPU 대비 속도 상대 차이
\(7.17524\times10^{-13}\), 변위 \(7.33627\times10^{-16}\)이었다.
HVP 호출9501회와 reset 제거 kinetic도 일치했다. CPU/GPU 구현 일치와 공간/시간 수렴은 구분한다.
추가 n32/sub128 strong의 CPU 실행은7255.413 s에 완료했고, reset 직전 속도까지 포함한
CPU/GPU 차이는 속도 상대 \(2.89993\times10^{-12}\), 변위 \(9.52711\times10^{-16}\)이었다.
이 n32의 HVP 호출은 CPU14329/GPU14330회로1회 달랐다. 동일 정수 반복 횟수 대신 동결한
허용오차의 실제 잔차와 궤적 일치를 검사한다. CPU v5의 wind24개/비교20개도 완료했다.
같은 GPU source의 강한 바람에서 다음 수치 보간 상한을 얻었다. 각 원본의 모든 step은
CPU 원식으로 별도 검산했다. 속도·변위/증분 모두 기존1\% 기준을 통과했다.
\begin{center}
\begin{tabular}{lrr}
\toprule
0.05 s strong/reset 비교 & 속도 상한(\%) & 변위/증분 상한(\%)\\
\midrule
n16$\to$32, sub128 & 0.3917943 & 0.0182097\\
n32, sub128$\to$256 & 0.2437471 & 0.00335713\\
n16$\to$32, sub256 & 0.4000651 & 0.0176079\\
n32, sub256, forward/backward & 0.0406189 & 0.00361927\\
\bottomrule
\end{tabular}
\end{center}
약한 바람의 n32/sub256 방향 비교는 속도0.0397466\%, 변위0.00376005\%다.
Finest strong 궤적의 Bernstein 상한은 \(r\le0.03415697\), \(J\ge0.93285276\),
mid-surface strain0.001006379, 선형 표면 strain0.017308061이다.
4개 실제 상태의 구적6/4 대10/8 force 차이는0.000959886\% 이하였다.
독립 공간 기준과 재료 허용 strain의 최종 채택은 여전히 별도 조건이다.
격리한 project bootstrap40개에서 실제37개 import 및 producer22개 identity를 확인하고,
n16 재실행의17개 배열/4,228,422 scalar와 step diagnostics가 정확히 같은지도 확인했다.
실행 script·원본·성능 측정 범위는
\code{experiments/R1_teacher_velocity_reset/p3_shell_gpu/README.md}에 보존한다.

\paragraph{선택한 장시간 검증 조건.}
사용자는 다음1.5 s 검증을 기존0.25--0.5 m/s knot 목표 속도로 먼저 진행하고 결과에 따라
확대하기로 선택했다. Seed20260909/PCG64의60 Hz,90 frame,12 frame 간격 knot와 벡터 선형 보간,
마지막18 frame zero ambient를 유지한다. 보간된 실제 풍속은 knot 하한보다 작을 수 있다.
Rest natural과18/42/66 frame에서 위치·시각을 보존하고 속도만0으로 둔 독립 분기를 계산한다.
Natural parent의 완결된 checkpoint를 사용해 공통 prefix의 중복 전진을 줄이며,
parent manifest·source·material·grid·wind identity와 reset 제거 kinetic을 보존한다.
프레임별 원본 저장, 모든 interval의 원식 검산, 시간/공간/방향 비교와 재시작 검사를 분리한다.
CUDA 검산을 사용할 때에는 모든 interval을 CUDA로 재계산하고 frame별3개 상태와 공력을
CPU로 추가 대조했음을 명시한다. 이를 모든 상태의 CPU 검산이라고 소개하지 않는다.
실행 및 후속 결과는 \code{experiments/R1_teacher_velocity_reset/p3_shell_random/README.md}에 기록한다.
추가 학습데이터 생성과 R1 완료 판정은 계속 보류한다.

\paragraph{1.5초 개발 검증 완료와 적용 범위.}
N8/sub128, n16/sub128, n8/sub256, n16/sub256 forward/backward의 다섯 조건에서
natural과 세 reset을 완료했다. 총20개 물리 원본/239,616 interval의 원식·에너지·고정 조건과
Bernstein 기하 검산을 통과했다. CUDA의 전 interval 재계산과 별도로1,170 frame의
CPU 상태3,510회 및 공력1,170회를 대조했다. 상태 수는 frame 경계 중복을 포함한다.
다섯 비교의 네 분기 모두 아래 전체 수치 보간 속도 상한과 변위/증분에서1\%를 통과했다.
\begin{center}
\begin{tabular}{lrrrr}
\toprule 비교 / 속도 상한(\%) & Natural & Reset 0.3 s & Reset 0.7 s & Reset 1.1 s\\\midrule
공간 n8$\to$16, sub128 & 0.699080 & 0.910953 & 0.848095 & 0.720628\\
공간 n8$\to$16, sub256 & 0.694414 & 0.904299 & 0.842757 & 0.711993\\
시간 sub128$\to$256, n8 & 0.108164 & 0.136703 & 0.119953 & 0.169299\\
시간 sub128$\to$256, n16 & 0.096664 & 0.119147 & 0.096074 & 0.148473\\
방향 n16, sub256 & 0.044940 & 0.055049 & 0.045901 & 0.063805\\
\bottomrule\end{tabular}\end{center}
변위의 최대 상한은0.071417\%, 분기 시작 상태 대비 증분은0.072494\%다.
공간 차이가1\%에 가까워 n8/sub256을 추가했고, 시간 간격을 절반으로 줄여도 공간 최대 속도
상한이0.910953\%에서0.904299\%로 비슷하게 유지됨을 확인했다.
이는 이번 약한 바람 조건의 세분 비교다. 독립 비선형 공간 기준 통과로 소개하지 않는다.

Finest n16/sub256 natural의 최대 nodal 변위는9.067314 mm다.
0.3/0.7/1.1 s reset에서 위치·시각을 유지하고 제거한 kinetic은
4.196988/7.363604/16.355723 $\mu$J다.
마지막0.3 s zero ambient의 누적 외력 일은 natural에서
\(-1.378833\,\mu\mathrm{J}\)이고,20개 원본 모두 음수였다.
Relative-velocity drag를 유지한 결과이며 별도 구조 감쇠를 추가하지 않았다.
모든 조건의 checkpoint42 무초기화 재시작에서9개 배열과 step diagnostics가 정확히 같았다.
N8/n16 sub128의 scalar 수는369,848/1,409,474개이며,
sub256은735,416/2,802,626개다.
N8/sub128은 격리 runtime에서도 같은 재시작과 원식 검산을 확인했다.

Finest natural의0.3/0.7/1.1/1.5 s와 최대 nodal 변위 시각0.603385417 s의5개 실제 상태를
CPU 구적6/4 대10/8로 대조했다. 최대 상대 차이는 energy \(5.05927\times10^{-12}\%\),
탄성력 mass-dual0.00000742146\%, 움직이는 공력0.000185059\%,
같은 위치에서 velocity0인 공력 \(2.02035\times10^{-13}\%\)였다.
이는 선택 상태의 구적 진단이며 전체 시간·독립 공간 인증은 아니다.

실행·검산 producer26개와 물리 허용오차는 그대로 유지했다. 재현 runtime v4는47개 project
source이며 추가4개 저장 검사와4개 비교 검사를 포함한 관련 고유 검사60개가 통과했다.
전체 repository test suite를 실행했다는 의미는 아니다. RMS 정규화 영역을 명시하면서
비교 wrapper만 제어 중단한 v1 기록은 solver 실패와 구분해 보존했다.
수치 근거·그래프·해시 연결은 위 보고서와 \code{evidence/verification/final_summary_v1.json}에 있다.

도달 형상에서 속도를 제거한 뒤 같은 미래 바람으로 전진하는 절차는 이번 약한 바람에서 성립했다.
그러나 약9 mm 변위는 큰 변형의 장시간 coverage를 보장하지 않는다. 사용자는 후속1--3단계로
동일 파형의4배(knot1--2 m/s), 같은 reset과 원식·공간·시간·방향 비교를 승인했다.
Seed·시각·재료·고정 조건·물리식·허용오차를 유지하고 필요한 수치 세분만 보완한다.
후속 풍속 확대와 학습 적격성 채택은 별도 단계다. 독립 비선형 공간 기준·재료 허용 strain·
공개 Registry/고차 GS/probe/oracle의 남은 조건을 유지한다. 추가 학습데이터0개, \code{training_eligible=false}, \code{r1_complete=false}다.

\paragraph{4배 파형의 입력 계약과 보완 결과.}
같은 seed·방향·시각의 바람 벡터를 \(\mathbf W_k^{(4)}=4\mathbf W_k^{(1)}\)로 확장한다.
\code{--wind-scale}~\code{4}는 knot 목표 속도1--2 m/s를 뜻하며 벡터 보간 중의 풍속 하한은 아니다.
60 Hz의 frame-start held force와 재료·고정 폭·물리식·solver 허용오차는 유지한다.
Chunk schema v2는 \code{wind_scale}과 재구성한 전체 wind program의 identity를 검사한다.
Parent 분기·원식 검산·비교·그림도 이 배율을 읽으며 서로 다른 배율의 원본을 섞지 않는다.
물리/GPU source22개는 그대로이고, 입출력 연결을 갱신한 producer26개와 재현 source47개는
runtime v5에 동결했다. 선행 schema v1 원본은 격리한 runtime v4로 계속 재현한다.
배율과 기존 경로의 관련10개 검사, 실제 CUDA smoke24 interval의 CUDA/CPU 검산과
격리한 v5의 CPU 검산을 통과했다. 이는 전체 repository test suite의 실행을 뜻하지 않는다.

N8/sub128은 natural과 세 reset의29,952 interval, 별도 checkpoint128 interval의 검산을 완료했다.
N16/sub256 정방향도 natural과 세 reset의59,904 interval 및 checkpoint256 interval의
원식·기하 검산을 마쳤고 checkpoint42의 정확한 재시작을 통과했다.
역방향의 네 원본과 checkpoint 검산도 완료했다. 처음 네 조건(n8/n16의 sub128,
n16/sub256 정·역방향)은 주원본179,712 interval과 checkpoint768 interval을 모두 검산했다.
N16/sub256 natural의 최대 nodal 변위는109.084864 mm다. 선택한5개 실제 상태의 CPU 구적6/4 대10/8 최대 차이는
탄성력 mass-dual0.000931460492\%, moving aero0.000646301517\%였다.
선택 상태의 구적 진단과 전체 시간의 수렴 판정은 구분한다.

전체 수치 보간 비교의 속도 상대 상한은 다음과 같다. N16과 n32의 시간·방향 비교는
모두 통과하나 n16$\to$32의 마지막 reset 공간 비교는 미달이다. 후속 값은2026-09-11 저장 report와
summary의 일치를 확인한 결과이며, 이번 문서 감사에서 전체 원본/물리 재검산을 다시 수행하지 않았다.
\begin{center}
\begin{tabular}{lrrrr}
\toprule 비교 / 속도 상한(\%) & Natural & Reset 0.3 s & Reset 0.7 s & Reset 1.1 s\\\midrule
공간 n8$\to$16, sub128 & \textbf{1.087962} & \textbf{1.211015} & \textbf{1.193883} & \textbf{3.495261}\\
시간 sub128$\to$256, n16 & 0.144732 & 0.158862 & 0.143669 & 0.731830\\
방향 n16, sub256 & 0.059589 & 0.068246 & 0.062644 & 0.274527\\
공간 n16$\to$32, sub256 & 0.444132 & 0.467933 & 0.519905 & \textbf{1.855258}\\
시간 sub128$\to$256, n32 & 0.140173 & 0.165289 & 0.163890 & 0.678680\\
방향 n32, sub256 & 0.022578 & 0.025601 & 0.023865 & 0.106118\\
\bottomrule\end{tabular}\end{center}
굵은 값은 기존1\% 기준 미달이다. Natural의 최대 속도 차이는1.4578125 s에서
0.003749052115 m/s RMS이며, 오차 제곱합의98.682105806\%가 초기 법선 방향 성분에 있다.
이는 오차 방향의 진단이며 단독으로 원인을 확정하지 않는다. N16 시간 비교 통과를 n8의 시간
오차까지 직접 측정한 결과로 소개하지 않는다.
Reset 1.1 s의 절대 속도 차이는0.001871508541 m/s RMS로 natural보다 작지만,
reference peak도0.053544172024 m/s로 작아 상대 오차가 더 크다. 기준을 완화하지 않으며
이 거친 공간 비교의 네 분기 모두 미달로 보존한다.
Coarse 공간 비교는 sub128, 보완 공간 비교는 sub256이다. 두 값만으로 공간 수렴 차수를
추정하지 않으며 n16/n32의 시간 비교를 별도로 제시한다. 보완 공간 판정 자체는 동일 sub256에서 수행했다.

허용치를 유지하고 n16$\to$32/sub256, n32의 시간 sub128$\to$256 및 대각선 방향 비교를
추가했다. 모든 natural/reset 분기와 원식·기하를 확인하는 판정 계획을 결과 전에 저장했다.
후속 계산은 기존 원본을 보존하고 아래의 별도 producer 구간을 잇는 실행 묶음으로 마쳤다.
N32/sub128 및 sub256 natural은 각각11,520/23,040 interval의 원식·기하 검산을 통과했다.
Sub256 natural의 최대 nodal 변위는109.091615 mm, Bernstein 면적비 하한은0.960399302,
mid-surface/선형 fibre strain 성분 상한은 각각0.000294309290/0.003948190025다.
최대 frame 에너지 결산 차이는 \(7.51789\times10^{-12}\) J, zero-ambient tail 일 합은
\(-0.000765905530\) J다. N32/sub128의 checkpoint42에서9개 배열의5,502,710개 값과
step diagnostics가 정확히 일치했다.
후속 저장 요약은 필수20개 비교 중19개 통과와 선택 상태 구적 통과를 기록한다.
Reset1.1 s의 보완 공간 속도 절대 상한은0.000993330 m/s, 기준 peak 하한은0.053541309 m/s다.
작은 분모를 이유로1\% 기준을 바꾸지 않으며, 거친 격자의 네 실패도 보존한다.
실행 종료와 물리 수렴 채택을 구분하고 실패 원인·전체 원본 검토는 남긴다.
실행·검산·형상과 입력 근거는
\code{experiments/R1_teacher_velocity_reset/p3_shell_random/scale4/README.md},
최신 보완 결과와 확인 범위는 그 아래 \code{fast_handoff/README.md} 및
\code{evidence/handoff_20260911/}이 소유한다.

\paragraph{Teacher 실행 병목 보완과 채택 경로.}
\label{sec:r1-teacher-execution-cost}
동일 수식·float64·허용오차에서 HVP에 불필요한 힘/에너지 출력과 회수를 분리하고 CUDA graph로
호출 비용을 줄였다. CuPy 기본 GMRES가 작은 검사에서 HVP8--10회를126회로 늘려 조기 종료를
직접 구현했다. 공유 비기본 stream과 작업 벡터 복사로 graph capture/aliasing 문제를 보완했으나,
현재 GTX1080Ti의 반복 희소 LU는 CPU보다 느려 최종 채택하지 않았다.
2026-09-11 당시에는 \textbf{CPU 반복 풀이+GPU HVP/CUDA graph}를 선택했고 추가 GPU 이식을 보류했다.
이는 과거 상태다. 2026-09-14 현행은 아래 후속 GPU 상주 개발 경로이며, 당시 CuPy 시제품의 미채택을 현재 cuDSS 경로에 적용하지 않는다.

32격자/sub256/1.5 s 전체 비교에서 계산 시간은9,991.43 s에서8,651.35 s로13.4124\% 줄었고,
기록된 준비·저장·검산 포함 합계는11,472.26 s에서10,172.67 s로11.3281\% 줄었다.
양쪽23,040 interval의 자동 검산과 전체 궤적 대조를 확인했다. 이는 단일 순차 비교이며
다른 GPU 부하와 재시작 영향이 있어 잠정 성능이다. Student 실시간 실행 속도의 근거가 아니다.
정확한 계측·허용치·source/원본 검토 범위는
\code{experiments/R1_teacher_velocity_reset/p3_shell_random/profiling/full_run/README.md}를 따른다.

기존 계산은 strict source 계약을 완화하지 않고 별도의 segmented view로 보존했다.
원본과 새 HVP/graph 구간의 producer identity, 경계 상태·재시작·원식 검산을 각각 연결한다.
CPU/CUDA의 작은 전체 흐름 검사와 원본 변조 거부를 시험했으며, 다른 producer의 짧은 일치를
모든 궤적의 정확한 동등성으로 해석하지 않는다. 실행 구조와 근거는 위 \code{fast_handoff/README.md}가 소유한다.

\paragraph{10초 시간 간격 탐색의 계약과 미확정 결과.}
\label{sec:r1-ten-second-timestep}
사용자는 장면별10--20 s 길이를 목표로 우선10 s의 큰 시간 간격을 찾기로 했다.
32격자/rest/동일 재료·고정 조건에서60 Hz frame-start held force를 유지하며
분할 수1/4/16/64/256, 즉 기존 대비 \(\Delta t\)256/64/16/4/1배를 순서대로 시험한다.
같은 후보를0.5/2/10 s로 연장하되 실패 시 남은 구간을 생략한다. 다음 후보는 같은 rest에서 시작한다.
안정 후보의 경계를 세분한 뒤 같은10 s의 \(\Delta t,\Delta t/2,\Delta t/4\) 세 쌍 비교를
모두1\% 미만으로 확인해야 정확도 추천군에 넣는다. 실패 상태를 작은 간격으로 이어 큰 후보의 성공으로 처리하지 않는다.

Wind는 seed20260909,0.2 s knot의 목표1--2 m/s 벡터 보간과 마지막2 s의 zero ambient다.
10 s용 새 배열이므로 앞선1.5 s와 동일 prefix로 취급하지 않는다. 후보끼리는 실제 wind 배열을
공유하지만 힘은 각 후보의 현재 형상·속도에서 계산한다. Reset·공간 수렴·새 재료는 이 탐색에 추가하지 않았다.
최대256배는 공력 갱신 경계를 보존하기 위한 상한이며 검증된 정확도 상한이 아니다.
시간에 따른 wind 벡터 보간과 각 frame의 형상 의존 힘을 hold하는 계약을 구분한다.

저장 최종 상태는 \code{no_stable_candidate}다. 256배는80 frame 뒤 \code{linear_solve},
64/16/4배는 각각157/158/158 frame 뒤 \code{line_search}로 실패했다.
기존 간격은82 frame에서 후보 활성 실행2시간 한도로 중단되었다.
10 s 통과 후보와 시간 정확도 추천은 없으며, 모든 간격의 물리적 불안정을 입증한 결과가 아니다.
비슷한 실패 시점의 원인은 아직 분석하지 않았다. \code{geometry_unresolved}는 기하 충분조건 미확정,
\code{resource_limit}는 운영 한도이며 실제 뒤집힘·수치 폭발과 구분한다.

GPU가 사용 중이어도 기다리지 않고 시작하지만 후보는 순차 실행하고 같은 출력의 중복 실행을 막는다.
한도는 후보24개·전체 활성12시간·후보별2시간·단계 응답300초다. 기존 물리 허용오차는 완화하지 않았다.
실행기10개 검사와 작은 CPU 중단/재개 흐름은 통과했으며, 이번 저장 결과 확인은
동결 runtime/plan/wind170개와 report5개 hash 대조에 한정한다. 전체 raw NPZ/원식 재검산은 별도다.
설정·실행 명령·실패 원본과 검증 범위는
\code{experiments/R1_teacher_velocity_reset/timestep_search/README.md}와 그
\code{evidence/handoff_20260911/}을 따른다. R1 전체·학습 적격성·추가 발행 보류는 유지한다.

\subsection{Canonical probe와 mapping}

\begin{itemize}
  \item[$\square$] Probe ID/rest position/area--mass weight와 source-object scope를 serialize한다.
  \item[$\square$] Teacher 및 각 independent GS map의 ordered support와 law별 weight를 deterministic하게 생성한다.
  GS/선형 teacher는 nonnegative, 고차 teacher는 별도 version의 signed shape evaluation과 conditioning 검사를 요구한다.
  \item[$\square$] Partition, constant, coordinate-linear/smooth-field reproduction, coverage와 unsupported rate를 측정한다.
  \item[$\square$] 모든 declared variant의 common-valid intersection을 test open 전에 한 번 동결하고 hash한다.
  \item[$\square$] Displacement/velocity forward map과 total-force/traction adjoint의 shape/unit/work identity를 property test한다.
  \item[$\square$] Spatial convergence, temporal convergence와 GS mapping noise를 별도 error component로 보고한다.
\end{itemize}

\subsubsection{현재 선형 probe와 비교기 구현}

Rest surface의 fixed probe, positive area/mass quadrature와 barycentric support를 저장한다.
Partition, constant/affine reproduction, coverage, unsupported reason 및 map hash를 검사했다.
현재 public map은 선형 3-support/nonnegative 계약이다.
Forward 식~\eqref{eq:r1-probe-map}과 \S\ref{sec:r1-adjoint}의 total-force/traction adjoint를 검사했다.
고차 Teacher에는 동일 식의 별도 shape law와 검증을 구현해야 한다.
Total force N과 traction N/m\(^2\)를 구분하고 area를 두 번 곱하지 않는다.
Probe trajectory를 raw manifest에 연결하고 v1/v2 원본 재계산을 지원한다.
비교기는 같은 object/material/rest/input/BC 등 고정 조건을 검사한 뒤 velocity/tip/work/대역 PSD를
SI 및 무차원 값으로 보고한다. 실행 성공만으로 물리 수렴을 선언하지 않는다.
Independent GS 두 개의 reconstruction, common-valid intersection과 network-free oracle는 미완료다.

\subsection{Traction/exact-ZOH fixture}

Frame 시작의 \((\mu_i^{t-1},v_i^{t-1},n_i^{t-1})\)에서
\[
  v_{{\rm rel},i}=W_t^{\rm eff}(\mu_i)-v_i,\quad
  v_{n,i}=v_{{\rm rel},i}^\top n_i
\]
를 만들고 sketch \code{eq:rd-wind-pullback}의 two-sided normal-only traction, numerical guard, area force와
\(Q=B^\top F\) 순서를 그대로 사용한다.
% [논문 작성 참고 -- 수정 이유]
% V0 fixture에서 v_parallel을 제거해 접선항이 구현에 조용히 남지 않도록 한다. 필요 시 optional law/version에서
% 별도 fixture와 ablation으로 다시 추가해야 한다.

\begin{itemize}
  \item[$\square$] \(n\mapsto-n\)에서 traction/rollout identity를 통과한다.
  \item[$\square$] \(v_{\rm rel}\mapsto-v_{\rm rel}\)에서 traction sign이 반전된다.
  \item[$\square$] Zero ambient에서 drag work \(\sum_iF_i^\top v_i\le0\) analytic fixture를 통과한다.
  \item[$\square$] Random state와 piecewise-constant load의 exact-ZOH 결과를 very-fine reference와 비교한다.
  \item[$\square$] Aero-off \(Q=0\)에서 positive pole/damping의 energy decay와 deterministic replay를 통과한다.
  \item[$\square$] Hidden corrector/substep, implicit alternative, force-level clamp와 second update가 trace에 없고
  quantitative core의 \(\tau_{\mathrm{guard}}\) activation count가 0임을 검사한다.
  % [논문 작성 참고 -- 수정 이유]
  % guard가 활성인 rollout을 정상 표본으로 받아들이면 solver instability를 공력 saturation으로 숨길 수 있다.
\end{itemize}

\subsection{Network-free oracle preflight}

\begin{itemize}
  \item[$\square$] Independent GS마다 attached/invalid policy를 적용한 raw Global field와 fit provenance를 저장한다.
  \item[$\square$] R0 valid policy와 frozen oracle-measure ID로 mass Gram, eigenvalue/rank floor 및 whitening transform을 기록한다.
  \item[$\square$] Open decision에서 선택한 canonical oracle path 하나만 실행하고 path/config/hash를 기록한다.
  \item[$\square$] Arbitrary whitening+unchanged diagonal pole, mapped teacher mode-only와 post-hoc asset pole tuning을 validator가 거부한다.
  \item[$\square$] Attachment zero, \(B^\top M_xB=I\), positive rad/s pole/damping, impulse/step/FRF/work와
  exact-ZOH free-decay fixture를 통과한다.
  \item[$\square$] 두 GS의 oracle response를 같은 probe/mask에 올리고 teacher convergence와 mapping-noise floor에 대조한다.
\end{itemize}

\subsection{Base mean/covariance transport}

\begin{itemize}
  \item[$\square$] Teacher surface displacement/velocity/normal을 training/evaluation correspondence로 각 GS에 transport한다.
  \item[$\square$] Rest/zero-state mean과 covariance가 bitwise 또는 declared tolerance 안에서 input identity다.
  \item[$\square$] Known proper rigid rotation에서 normal/rotation 방향, covariance eigenvalue와 SPD를 보존한다.
  \item[$\square$] Attached Gaussian mean/angular response가 numerical zero이고 invalid covariance가 stable failure reason으로 집계된다.
  \item[$\square$] Transported position finite difference와 reported kinematic velocity가 tolerance 안에서 일치한다.
  \item[$\square$] Appearance/SH loss 없이 base renderer가 finite image를 만들고 실패 frame을 denominator에 유지한다.
\end{itemize}

\subsection{개발용 sample 생성·검증과 학습 적격성}
\label{sec:r1-samples}

기존 Newton native backend에서 CPU로 새 source 세 개를 만들었다.
1 m XZ square n8, \(M=0.1\) kg, fps60, substeps16, iterations10, 1 s/61 state,
5$\times$5 trapezoid probe 25개다. Native stretch/area stiffness \(10^3\), material damping 0.1,
bending 10 N, bending damping 0, gravity 0을 사용했다.
\begin{itemize}
 \item Wind pulse: 0.5 m/s, 0.2 s half-sine 뒤 0.8 s zero ambient.
 \item Step on/off: 0.5 m/s의 켜짐/꺼짐. Off에도 air drag가 남는다.
 \item Aero-off decay: \(\Delta Y=0.01x^2\), 공력 off, material damping 유지.
\end{itemize}
각 source의 60 interval을 12 interval/13 state window 5개로 나눠 총 15개를 만들었다.
Target array는 \([13,25,3]\), batch 크기는 \([4,4,4,3]\)이다.
Single development source-object group이며 train/dev/test 독립 객체 분할이나 GS pairing을 완료한 것이 아니다.

Raw/probe inventory와 label/input/work를 모두 대조하고 각 CPU source를 독립 재실행했다.
세 replay의 position/velocity/aerodynamic force/work 최대 차이는 모두 0,
pin drift와 guard count도 0이었다. 각각 maximum nodal speed는
0.202946126/0.471042246/0.133365527 m/s,
interval work 절댓값 합은 0.000789484374/0.006596602165/0 J다.
마지막 값은 순 외력 일이나 구조 energy drift가 아니다.
생성·대조·replay·batch 검증은 48.370 s가 걸렸다.

Dataset에는 static probe/weight/attachment, registry, split, per-window NPZ와 checksum manifest를 둔다.
\code{TeacherSampleDataset.open(..., allow_development=True)}로만 개발 데이터를 열 수 있고,
default는 \code{training_eligible=false} 데이터를 거부한다.
\code{verify_sources}는 manifest뿐 아니라 raw/probe에서 label을 재구성해 대조한다.
Tamper, incomplete output, rehashed corruption과 static geometry 오류를 검사했다.
저장된 NumPy batch 읽기는 Torch/Newton/Warp 없이 가능하다.

Raw source는 \code{experiments/artifacts/runs/teacher_sample_dataset/20260908_sample_v1},
dataset은 \code{experiments/artifacts/datasets/teacher_samples/20260908_sample_v1}다.
Source inventory 93개/1,150,138 byte, dataset manifest 포함 20개/125,293 byte다.
Compact evidence와 response PNG를 Git에 보존하고 NPZ 원본은 ignored artifact에 둔다.
\textbf{이 데이터는 P3 데이터가 아니며 후속 P3 통과로 eligibility를 올리지 않는다.}

\section{Open Design Decisions}
\label{sec:r1-chapter-5}
\label{sec:r1-decisions}

\begin{decisionbox}{해석 규칙}
이 표는 결정을 대신하지 않는다. 권장 기본값은 K0 구현을 단순화하기 위한 제안이며,
선택한 option은 registry/version, 근거 fixture와 동결 시점을 기록한 뒤에만 canonical R1 실행에 사용한다.
\end{decisionbox}

\begin{longtable}{L{0.18\linewidth}L{0.27\linewidth}L{0.29\linewidth}L{0.16\linewidth}}
\toprule
결정 항목 & 허용 선택지 & 권장 기본값(미동결) & 동결 시점\\
\midrule
\endfirsthead
\toprule
결정 항목 & 허용 선택지 & 권장 기본값(미동결) & 동결 시점\\
\midrule
\endhead
Teacher 구조 model & 사용자 선택: 유한 회전 shell & P3 Koiter/StVK 막·굽힘과 normal-jump 연결의 수식 및 solver 검증을 진행한다. 작은 재료 strain·큰 회전 범위이며 acceptance는 별도다 & 모델 선택 2026-09-09; 채택은 검증 후\\
Constitutive preset & \((E,\nu,h,\rho)\) 계열 / 동등한 명시 SI shell parameter & opaque class ID 대신 실제 SI tuple를 저장하되 \(h\rho=M_{\mathrm{ref}}/A_{\mathrm{ref}}\)를 derive/equality-constrain & teacher dataset 생성 전\\
Damping & Rayleigh \(\alpha M+\beta K\) / verified material damping / mode-target calibration & 좁은 target band에 calibration한 Rayleigh damping과 measured modal ratio report & convergence config freeze 전\\
Teacher integrator & implicit Newmark/generalized-\(\alpha\) / 충분히 작은 explicit reference & 개발 P3의 acceleration-form Newmark와 약한 바람 시간 진단 완료.10초 입력의 안정·정확 간격은 미확정 & canonical backend/입력 채택 전\\
Refinement ladder & 2-level / 3-level 이상, uniform / structured refinement & 가능한 경우 3-level과 일정 refinement ratio; 최소 2-level은 smoke only로 표시 & K0 run 전\\
Probe construction & canonical UV grid / rest-surface samples / geometry landmarks+grid & fixed rest-surface probe grid와 area/mass weights, tip landmark를 보조로 포함 & map 생성 전\\
Map kernel/support & 선형 barycentric / 고차 shape teacher, local weighted GS support & GS convex map 유지. P3는 signed 10-shape map의 별도 schema, polynomial reproduction과 work 검사 필요 & mapping-noise threshold freeze 전\\
첫 데이터 입력 범위 & rest-start 변화 바람 / 도달 형상의 velocity-reset / displaced free-decay & 탐색 순서는 아래 절로 결정. 최종 입력 범위는 미동결이며 기존 decay 요구 유지 & 최종 데이터 scope freeze 전\\
Effective traction scale & domain-wide fixed \(\kappa_{\mathrm{aero}}>0\) 후보 범위 & material/object와 독립인 scalar 하나를 dev에서 전체 domain에 동결하고 runtime에는 \(g_W\)만 노출 & teacher final run 전\\
% [논문 작성 참고 -- 수정 이유]
% F가 kappa_aero*g_W^2에 비례하므로 둘을 함께 tune하지 않는다. kappa_aero는 teacher/student identity,
% g_W는 prescribed wind strength control이라는 서로 다른 owner를 갖는다.
Numerical traction guard & fixed high \(\tau_{\mathrm{guard}}\) / guard 없이 hard reject & supported envelope 위의 fixed global guard와 core activation zero; activation은 OOD/failure로 유지 & 첫 stress fixture 전\\
% [논문 작성 참고 -- 수정 이유]
% tau_guard를 coefficient tuple에서 분리하면 논문에서 물리 가정과 수치 안전장치를 명확히 구분할 수 있다.
Oracle area source & teacher-quadrature 기반 oracle-only positive fit / deterministic GS covariance-footprint rule & 가장 단순한 positive normalized rule을 두 GS에 동일 적용하고 source/hash 공개; student measure claim으로 승계 금지 & oracle field fit 전\\
Oracle canonical path & whitened-field pole refit / full \(M,C,K\) congruence+re-diagonalization & reliable full reduced operators가 있으면 후자, 아니면 response-fit; 두 결과를 asset별 혼용하지 않음 & oracle acceptance run 전\\
Oracle rank/fit band & fixed rank와 impulse/time/FRF weight 후보 & R2와 같은 intended Global rank, converged teacher band 안에서 fit & oracle config freeze 전\\
Convergence metric/threshold & absolute SI / dimensionless reference-scale / 둘 다 & dimensionless primary+SI diagnostic, velocity/FRF/work를 함께 사용 & finest result를 보기 전 dev freeze\\
Peak analysis status & area/mass aggregate, band/window/prominence 후보 & R1은 convergence와 availability status만 출력; curriculum 채택은 R2b 소유 & spectral report 전\\
\bottomrule
\end{longtable}

\subsection{Rest-start 바람과 변형 상태의 velocity-reset: 설계와 개발 결과}
\label{sec:r1-velocity-reset-plan}

2026-09-09 사용자는 rest에서 시뮬레이션한 뒤 중간의 변형 position에서 속도를 0으로 만들고
계속하여 적절성을 판단하는 순서를 선택했다. Canonical sketch의
\code{sec:rd-initial-state-exploration}을 따른다. \textbf{아래 개발 비교는 실행됐으며 학습 결과는 없다.}
실행 전에 vector 보간·seed·horizon·checkpoint와 해상도 ladder를 선언했다.
이번 입력은 vector 선형 보간이며 급격한 반전, 다중 seed 또는 최종 supported envelope까지 검증하지 않았다.

\begin{enumerate}
 \item \textbf{원본 생성:} 원래 rest geometry/material/attachment를 고정하고 시간 변화하는 바람으로
 전체 물체의 연속 궤적을 만든다. Wind backend의 개발/accepted 상태와 forcing identity를 먼저 확인한다.
현재 선형 P3 처방 압력의 통과를 nonlinear wind 검증으로 승계하지 않는다.
 기존 scalar-speed runner의 seed만 바꾸어 방향 랜덤화를 구현했다고 표시하지 않는다.
 Vector wind program/version과 RNG 상태 또는 사전 생성한 시간별 입력의 hash를 연결해야 한다.
 \item \textbf{Checkpoint와 대조군:} 사전 지정한 물리 시각 \(t_j\)에서 위치·속도와 backend가 필요로 하는
 내부/적분 상태, wind 시각·seed를 저장한다. 먼저 개입 없는 checkpoint 재시작이 원본의 후속 궤적을
 정해진 수치 허용오차 안에서 재현하는지 확인한다.
 \item \textbf{속도 초기화 분기:} 각 checkpoint에서 위치와 rest 기준은 유지하고 전체 자유 DOF 속도를
 0으로 만든다. 같은 원본에서 각 시점의 독립 분기를 먼저 비교한다. 한 궤적에서 반복 초기화하는 경우에는
 누적 개입으로 별도 식별한다. 패치별 독립 reset은 이번 기본 비교에 포함하지 않는다.
 \item \textbf{같은 이후 바람으로 계속:} 바람의 시간이나 seed를 처음으로 돌리지 않는다.
 Frame 시작의 개입 후 각 분기의 geometry/velocity로 traction을 한 번 평가하여 hold한다.
 Newmark 등에서는 새 속도·힘과 일치하도록 acceleration과 적분 이력을 구성하고,
 개입 전 acceleration이나 stale force cache를 무조건 재사용하지 않는다.
 \item \textbf{적절성 판정:} 위치/고정부 보존, finite 상태와 guard, 운동에너지 변화,
 이후의 displacement/velocity·위상·회복 및 공간·시간 refinement를 확인한다.
 같은 물리 시각의 reset 규칙을 각 해상도에 적용하며, 두 분기의 차이 자체를 모델 오차로 해석하지 않는다.
\end{enumerate}

분기점 직전/직후를 \(-,+\)로 쓰면 해당 개입은
\begin{equation}
 \begin{aligned}
  x_T(t_j^+)&=x_T(t_j^-), & v_T(t_j^+)&=0,\\
  \Delta K_{\mathrm{reset}}
    &=-\tfrac12 v_T(t_j^-)^\top M_Tv_T(t_j^-)\leq 0 .
 \end{aligned}
 \label{eq:r1-velocity-reset}
\end{equation}
여기서 고정 DOF의 속도도 0이며 \(M_T\)는 해당 backend의 물리 mass다.
형상과 재료 상태를 보존하므로 개입 순간의 탄성 에너지는 유지한다.
제거한 kinetic energy는 별도 intervention ledger에 기록하고 aerodynamic work나 자연 damping loss로
합산하지 않는다. 동일한 wind field도 속도 변화로 relative-wind force가 달라질 수 있다.
% [논문 작성 참고 -- 개입과 초기 조건의 구분]
% Velocity reset은 형상을 새로운 rest로 지정하는 연산이 아니고, wind-off나 aero-off도 아니다.
% 자연 연속과 reset 분기는 서로 다른 초기 운동 상태를 가지므로 후속 궤적의 일치를 요구하지 않는다.

\textbf{후속 데이터/학습 판단.} 원본과 분기의 전체 위치·속도·풍장·실제 힘·개입 기록을 보존하고
그 뒤 common probe/patch와 시간 window를 추출한다. Reset을 가로지르는 불연속은 일반 학습 window에서
분리한다. 자연 연속 분기와 reset 분기를 서로 다른 train/dev/test에 넣지 않으며,
객체 일반화에는 기존 source-object-disjoint 규칙을 추가로 적용한다.
분기 시작의 teacher 상태가 student의 generalized state로 자동 변환된다고 가정하지 않는다.
Teacher가 수렴한 뒤 network-free oracle와 명시적 student 초기화 계약을 확인하고,
같은 학습 예산의 rest-start-only/혼합 데이터 비교로 추가 이득을 판단한다. 현재는 데이터 비율을 정하지 않는다.
기존 \(x^2\) 초기 변위의 실패와 canonical aero-off displaced decay 요구를 유지하며,
이 탐색만으로 \code{training_eligible}를 올리거나 R1/R2를 통과시키지 않는다.


\paragraph{실행 조건과 재시작 구현 (2026-09-09).}
기존 native Newton VBD의 \code{demo} 모드에 explicit gravity-off/rest 초기화를 적용했다.
Teacher 모드의 방향 고정/registry 계약은 바꾸지 않았다. 별도
\code{wind3dgs.velocity_reset_development.v1}이며 \code{training_eligible=false}다.
1 m XZ flag, x=0 위치 고정, 질량 0.1 kg, native stretch/area 1000/1000 N/m,
material damping 0.1 s, bending 10 N, bending damping off, \(\kappa=0.6\)을 사용했다.
CPU, iterations10, fps60, 90 frame(1.5 s)이다. PCG64 seed20260909로 12 frame마다
random unit direction와 0.25--0.5 m/s target를 생성하고 vector를 선형 보간했다.
보간 중 속력은 target 하한보다 작을 수 있다. 처음은 wind0, 마지막 18 frame은 ambient0/air drag 유지다.
Checkpoint18/42/66(0.3/0.7/1.1 s)에서 각각 원본으로부터 독립 분기했다.
공간 mesh4/8/16 at substeps32와 시간 substeps8/16/32 at mesh8의 합집합인 5조건이다.

\code{deterministic_prefix_replay_v1}은 새 solver에 저장한 wind prefix를 재생해 내부 상태를 재구성한다.
빠른 full-state deserialization은 아니다. 원본 prefix의 위치/속도를 \(10^{-7}\) m, m/s 이내로 대조하고,
두 state buffer의 velocity만 0으로 만든다. 다음 VBD step은 현재 입력 위치/속도로 inertia와 previous
position을 다시 구성하며, frame-start 공력을 새 상대 속도로 평가한다. Rigid body/contact 없는 fixture다.
Checkpoint 배열에는 개입 후 속도를 저장하고 개입 전 속도는 별도 배열로 보존한다.
따라서 reset 경계를 가로지르는 이전 interval의 끝 속도를 일반 smooth state처럼 읽으면 안 된다.

\paragraph{실행 무결성과 분기 응답.}
5개 원본/5개 독립 replay/15개 reset의 25개 trace, 2,250 interval/2,275 state를 생성했다.
54,000 VBD substep, CPU 373.876 s다. 모든 replay의 위치/속도/force/work 최대 차이는 0이다.
Pin drift=0, guard=0, 모든 상태가 유한하며 위치 보존과 식~\eqref{eq:r1-velocity-reset}을 검산했다.
현재 geometry/velocity/wind에서 NumPy로 독립 공력을 재계산하고 work와 제거 kinetic ledger를 대조했다.
Edge stretch 최대 1.000017469, rest 대비 face area 최소 0.999981978을 관측했다.
신규 8개와 Newton 관련 85개 검사, 새 프로세스의 raw/report 재계산이 통과했다.
이는 repository 전체 test discovery나 새 GPU 검사가 아니다.

\begin{longtable}{rrrr}
\toprule Reset [s] & 최대 nodal 변위 [mm] & 제거 kinetic [$\mu$J] & 이후 원본 대비 probe RMS 최대 [mm]\\\midrule
0.3 & 2.038361 & 9.689417 & 8.651398\\
0.7 & 21.667045 & 52.273766 & 16.597835\\
1.1 & 38.493980 & 67.954399 & 13.712264\\
\bottomrule\end{longtable}
위 표는 mesh8/substeps32다. 마지막 열은 서로 다른 물리 상태의 차이이며 모델 예측 오차가 아니다.
모든 reset 뒤 다시 움직였지만 1.5 s 관측으로 장기 안정성이나 평형 회복 완료를 주장하지 않는다.

\paragraph{공통 probe 해상도 진단: 속도 기준 실패.}
공통 25개 vertex와 총면적 1 m$^2$의 trapezoid weight를 사용한다.
각 비교는 공통 frame에서 area-weighted RMS 차이의 최댓값을 fine의 peak response로 나눈다.
변위는 원래 rest 기준이고 reset은 checkpoint 이후 구간이다. 1\%는 이 실행의 개발 진단 기준이며
연속 시간 상한·방향·독립 기준이나 canonical R1 acceptance를 대체하지 않는다.
\begin{longtable}{llrrr}
\toprule Finest pair & 분기 & 변위 [\%] & 속도 [\%] & 속도 절대 차이 [mm/s]\\\midrule
공간8$\to$16 & 자연 & 0.520181 & 12.040068 & 4.657035\\
 & 0.3 s reset & 0.667281 & 13.827964 & 4.743307\\
 & 0.7 s reset & 0.887863 & 20.503919 & 4.435040\\
 & 1.1 s reset & 0.654738 & 81.148231 & 3.389095\\
시간16$\to$32 & 자연 & 0.077190 & 2.176690 & 0.835145\\
 & 0.3 s reset & 0.069383 & 2.400700 & 0.816840\\
 & 0.7 s reset & 0.158086 & 3.546817 & 0.757634\\
 & 1.1 s reset & 0.223625 & 17.756006 & 0.555413\\
\bottomrule\end{longtable}
공간의 substeps는 32, 시간의 mesh는 8이다. 모든 finest pair가 속도 기준을 넘었다.
자연 연속의 공간 속도 차이는 6.2530\%에서 12.0401\%로 커졌고,
시간은 8.2148\%에서 2.1767\%로 감소했으나 기준에 못 미쳤다.
1.1 s reset의 fine peak 속도는 공간 4.17643 mm/s, 시간 3.12803 mm/s로 작으므로
절대 오차와 분모를 함께 해석한다. 공간 pair에도 시간 오차가 남으므로 mesh만의 오차로 단정하지 않는다.

\paragraph{굽힘 상태 coverage에 대한 추가 반례.}
x=0 모서리의 위치만 고정하면 원래 XZ 평면의 Z축 주위 강체 회전도 허용된다.
이는 고정부와 원소 rest 길이/면적/dihedral을 보존하므로 위치 변화의 크기만으로 탄성 굽힘을 판단할 수 없다.
25개 probe에서 면적 가중 최소제곱으로 이 한 자유도 회전을 맞추었다.
0.7 s는 회전 1.241208$^\circ$, 전체 변위 RMS 12.700918 mm 중 회전 잔차가 0.002380 mm이고,
1.1 s는 전체 22.540597 mm 중 잔차 0.035869 mm다. 이는 위치 분해이며 탄성 에너지 비율이 아니다.
% [논문 작성 참고] 큰 position 변화와 큰 internal elastic deformation을 구분한다.
% 현재 fixture는 pin-line rigid rotation을 허용한다. Reset 후 정지는 굽힘 복원력 검증과 같지 않다.
따라서 이번 성과는 분기·공력·개입 처리와 실패 발견까지다. 단일 seed/짧은 horizon/회전 위주의 상태로
실제 굽힘 상태의 데이터 적절성이나 학습 이득을 확정하지 않는다.
후속 후보는 고정 모서리의 기울기도 제한하거나 동일한 물리 폭의 고정 영역으로 굽힘을 유도하는 비교다.
해상도마다 두 vertex 열을 고정해 물리 폭이 달라지는 조건을 같은 boundary로 취급하지 않는다.
같은 backend의 시간 오차를 먼저 줄이고 공간 및 입력/상태 coverage를 다시 판단한다.
현재 native 결과를 P3 또는 nonlinear accepted Teacher로 승계하지 않는다.

\paragraph{재현과 보존.}
실행은 \code{code/scripts/check_teacher_velocity_reset.sh}, 읽기 전용 재검산은 같은 launcher의
\code{--verify}다. Raw root는 \code{experiments/artifacts/runs/teacher_velocity_reset/20260909_velocity_reset_v1}이며
manifest 제외 inventory55개/6,791,538 byte다. Producer source25개와 설치 Newton VBD source hash를 대조했다.
\code{experiments/R1_teacher_velocity_reset/README.md}에 정확한 명령과 결과, config/environment,
report/verification/rotation diagnostic, 원본 manifest/provenance와 응답·refinement 그림을 보존한다.
새 학습 patch/window dataset을 발행하지 않았으며 기존 15개 development window와 별도다.

\paragraph{후속 고정 폭 비교와 완료 조건의 변경.}
사용자는 왼쪽 0.25 m 영역 고정을 선택했다. 같은 1 m flag의 \(x\leq0.25\)를 고정하고
자유 길이 0.75 m, 원래 rest/material/mass/wind를 유지한다. Mesh4/8/16에서 경계가 정확히 일치한다.
별도 \code{wind3dgs.velocity_reset_development.v2}이며 기존 v1을 덮어쓰지 않았다.
실패한 개발 sample만 발행하는 안 대신 \textbf{실패 원인을 해결한 뒤 sample을 생성}하도록 요청했다.
수렴 기준과 \code{training_eligible=false}를 유지하며 아직 새 dataset을 발행하지 않았다.

첫 strip 비교는 25 trace/CPU320.758 s다. Finest 자연 연속의 시간16$\to$32 차이는
변위22.623233\%/속도29.463330\%, 공간8$\to$16은 변위117.752331\%/속도140.233689\%였다.
Mesh8/substeps32의 자연 연속에서 전체25 probe를 면적 가중 최적 평면에 맞춘 잔차는
0.3, 0.7, 1.1, 1.5 s에 각각 0.099051, 0.034441, 0.291320, 0.214125 mm다.
즉 비평면 굽힘을 만들었지만, 형상 진단 통과가 물리 수렴 통과는 아니다.

\paragraph{시간 간격과 solver 반복 분리.}
자연 연속만 검사한 세 후보는 CPU110.384 s다. Substeps32에서 iterations10$\to$20의 속도 차이는
0.436875\%, iterations10에서 substeps32$\to$64는1.799880\%, 64$\to$128은1.257692\%다.
이어 mesh8/substeps64,128에서 원본/독립 replay/세 reset을 실행했다.
10 trace/900 interval/910 state/CPU372.303 s이며 상태·공력·work·개입 검산을 통과했다.
공통 probe 해상도를4로 명시하여 같은25 probe를 유지했고, 추가 run에는 공간 pair가 없다.
\begin{longtable}{lrrrr}
\toprule 분기 & 변위 [\%] & 속도 [\%] & 속도 차이 [mm/s] & Fine peak [mm/s]\\\midrule
자연 연속 & 1.045758 & 1.257692 & 0.232614 & 18.495303\\
0.3 s reset & 1.050965 & 1.252818 & 0.218555 & 17.445073\\
0.7 s reset & 0.959697 & 1.469637 & 0.261337 & 17.782454\\
1.1 s reset & 1.005909 & 0.974783 & 0.179629 & 18.427539\\
\bottomrule\end{longtable}
변위와 속도가 모두1\% 이하여야 하므로 네 분기 모두 실패다. 분모는 각 fine 분기의
해당 구간 최대 probe RMS이며 reset suffix만 비교한다. 시간 차이는 감소했지만 통과하지 않았고,
초기 공간 실패가 해결된 것도 아니다. Float32가 주원인이라는 판단도 이 결과만으로 확정하지 않는다.

\paragraph{현재 고정 조건에서도 남는 native 굽힘의 격자 의존성.}
동일 형상 \(w(X)=10^{-3}[\max(X-0.25,0)/0.75]^2\) m에 대해
\(E_b=\tfrac12(10\,\mathrm{N})\sum_e\ell_e\theta_e^2\)를 float64 기하학으로 재계산했다.
Mesh4/8/16/32의 에너지는 각각8.888856504/5.185163554/2.777765415/1.435178592 $\mu$J다.
독립 strip slope 해석식과 대조했으며 n32는 n4의 약16.15\%다.
이 계산에는 시간 적분·질량·감쇠·바람·solver 반복이 들어가지 않는다.
따라서 시간 refinement만으로 이 물성 이산화의 문제를 없앨 수 없다.
다만 전체 동역학 오차의 기여율 모두를 정적 에너지로 설명하지는 않는다.
기존 native 감사와 일치하는 반례이며, 방향 편향으로 탈락한 면적 가중 보정을 재채택하지 않는다.
% [논문 작성 참고] 고정 폭으로 회전 위주 상태는 보완했지만 native 물성의 격자 의존성은 남는다.
% 과거 감사에서 확인한 문제를 이번 fixture의 선행 조건으로 연결한다. 수렴 기준 완화로 해결하지 않는다.

\paragraph{준비된 sample 계약과 미결정 해결 범위.}
Exporter/NumPy loader는 임시 CPU 자료로 검사했다. 12 interval/13 state와
5$\times$5 probe의 네 개 겹치는3$\times$3 patch를 사용하며, 겹침 probe 면적을 분배하여
총면적1 m$^2$/질량0.1 kg을 보존한다. Natural 전체와 reset 이후 suffix만 사용한다.
중복 prefix/replay·개입을 가로지르는 interval·불완전 tail은 제외한다.
이 native 단계의 계획은19 window$\times$4 patch=76 sample이며 native source의 발행 수는0이다.
각 창의 초기 displacement/velocity, 절대 시각·풍장·전체 coupled context와 source group을 보존한다.
전역 aero work의 patch별 복사본을 합산하지 않으며, teacher context는 target inference 입력이 아니다.
관련17개 검사/10.979 s는 경계·원본 대조·변조 거부·면적 보존·batch·의존성 분리에 대한 통과다.

해결 선택지는 기존 P3 선형 판을 작은 굽힘의 actual wind/reset에 연결하거나,
Newton 비선형 굽힘 구성·힘 계산을 보완하는 것이다. 전자는 작은 변형으로 범위가 제한되며
새 BC/load/mass/map/time/space 검증이 필요하다. 후자는 큰 변형 목표와 연결되지만
구성식·미분·방향·동역학 검증으로 이번 sample 작업 범위가 커진다. 사용자에게 장단점을 질문했으며
후속 사용자 결정은 \textbf{기존 P3 활용}이다. 아래에서 새 BC/load를 다시 검증했으며,
이전 pressure 통과를 새 wind/BC 통과로 승계하지 않았다.
상세 조건·원본 hash·producer snapshot·정적 진단과 재현 명령은
\code{experiments/R1_teacher_velocity_reset/clamped_samples/README.md}에 보존한다.

\subsection{검증된 작은 굽힘 P3 wind/reset sample}
\label{sec:r1-p3-wind-samples}

\paragraph{선택 범위와 물리 계약.}
2026-09-09 사용자가 작은 굽힘 P3 활용을 선택했다. Native 계수를 고쳐 동등한 재료로 만든 것이 아니라,
기존 KL/P3 구성과 consistent mass를 재사용한 별도 범위다.
\(E=10^6\) Pa, \(\nu=0.3\), \(h=0.01\) m, 면밀도0.1 kg/m$^2$다.
전체1 m$^2$ 중 \(X\leq0.25\)는 움직이지 않는 지지 영역이며,
자유 판 \([0.25,1]\times[0,1]\)의 질량은0.075 kg, 지지 영역은0.025 kg이다.
구조 감쇠는 기존 P3처럼0이며 정지 공기의 relative-wind drag는 유지한다.
원래 seed/60 Hz/90 frame/12 frame vector 보간/reset18,42,66을 유지하되
target 풍속 상한을0.05 m/s로 낮춰 작은 굽힘을 검사한다. \(|w|\leq0.1h=1\) mm,
\(\|\nabla w\|\leq0.01\)을 넘어가면 실패다. 큰 변형·새 풍속/seed로 검증 범위를 확대하지 않는다.

\paragraph{고정 폭을 판 경계에 연결.}
P3 자유 영역 경계 \(X=0.25\)의 변위0은 강하게 부과하고 기울기0은 다음 에너지로 약하게 부과한다.
\begin{equation}
 E_{\Gamma_D}=-\int_{\Gamma_D}(\partial_n w)M_{nn}(w)\,ds
   +\frac12\int_{\Gamma_D}\frac{\alpha}{h_E}(\partial_n w)^2\,ds,
 \qquad \alpha=Eh^3(3/2)^2 .
 \label{eq:r1-p3-clamp}
\end{equation}
내부 edge의 half-average 대신 \textbf{full boundary moment}를 쓴다. 기존 P3 penalty를 튜닝하지 않았다.
독립 tensor cubic spline은 같은 자유 사각형에서 첫 두 x coefficient를 제거하여 변위·기울기를 강하게 고정한다.
고정 영역을 제외하면서 기존 조립기의 질량을0.075 kg으로 환산하여 면밀도를 보존한다.
% [변경 이유] 기존 pressure fixture는 position-only여서 새 고정 폭 조건과 같지 않다.
% 기존 P3 source/pressure evidence는 그대로 두고 자유 영역과 weak slope BC adapter를 별도로 추가했다.
공식 FEniCS-Shells의 clamped KL boundary 항을 웹으로 확인했으며 정확한 URL은 새 실험 README에 보존한다.

\paragraph{Signed 공력 map과 전체 모드 시간 전진.}
형상은 \(r(x,y)=(x,w(x,y),0.5-y)\), 면적 법선은
\(nJ=(-w_x,1,w_y)\)다. Current normal/area와 상대풍에서 동일 \(\kappa=0.6\)의 양면 traction을 계산한다.
P3 shape의 transpose로 법선 DOF의 generalized force를 모으고 quadrature power와 대조한다.
면내 제약·지지 영역의 반력 몫은 위치 변화가0이므로 work에 기여하지 않는다.
매60 Hz 시작의 force를 hold하고 consistent-mass 정규직교 \(\Phi\)의 \textbf{전체 모드}를 유지한다.
\begin{equation}
 \begin{aligned}
  \Phi^\top M\Phi&=I,\quad \Phi^\top K\Phi=\Omega^2,\quad
  \ddot q+\Omega^2q=f_q,\quad f_q=\Phi^\top f,\\
  q^+&=\cos(\Omega\Delta t)q+\Omega^{-1}\sin(\Omega\Delta t)\dot q
       +2\Omega^{-2}\sin^2(\Omega\Delta t/2)f_q,\\
  \dot q^+&=-\Omega\sin(\Omega\Delta t)q+\cos(\Omega\Delta t)\dot q
       +\Omega^{-1}\sin(\Omega\Delta t)f_q,\\
  W_k&=f_q^\top(q^+-q),\qquad E_k=\tfrac12(\|\dot q\|^2+\|\Omega q\|^2).
 \end{aligned}
 \label{eq:r1-p3-exact-held}
\end{equation}
Reset은 \(q\)를 보존하고 \(\dot q\)만0으로 만든다. 제거 kinetic을 별도 ledger에 둔다.
구조 내부1/2/4분할의 최대 상대 차이는 \(8.876\times10^{-15}\) 미만이고,
모든 모드를 독립 scipy expm과 비교한 energy-scaled 차이는 \(1.248\times10^{-11}\) 미만이다.
이는 지정한 frame-start held-force 구조 전진의 검증이며 force clock 자체의 연속 시간 수렴 인증은 아니다.
Low-pass·modal truncation으로 속도 차이를 숨기지 않았다.

\paragraph{새 물리 조건의 수렴 결과.}
P3 n4/n8/n16, n16 반대 대각선, n16 Duffy6 공력, spline16/32의7개 모델에서
자연 연속과 세 reset을 실행하고 n16 내부2/4분할을 추가했다.
각 최종 run은36 trace/3,240 interval/3,276 state다. 비교는 고정 영역의0응답을 포함한 전체1 m$^2$ norm이다.
각 hold 구간128등분의 최대 차이에 velocity/acceleration Lipschitz 상한을 더하고,
fine sampled peak로 나누어 보수적인 연속 시간 상대 상한을 만든다. Reset은 개입 후 suffix다.
\begin{longtable}{lrrl}
\toprule 네 분기 중 최악값 & 변위 상한 [\%] & 속도 상한 [\%] & 판정\\\midrule
P3 4$\to$8 & 0.489434 & 1.423535 & coarse 실패 보존\\
P3 8$\to$16 & 0.250302 & 0.882329 & 통과\\
P3 n16 대각선 & 0.224174 & 0.224285 & 통과\\
공력 Duffy4$\to$6 & 0.209477 & 0.197062 & 통과\\
Spline16$\to$32 & 0.227642 & 0.505603 & 통과\\
P3 n16/spline32 & 0.221238 & 0.386575 & 통과\\
\bottomrule\end{longtable}
기준1\%를 완화하지 않고 n16을 sample source로 사용했다. 끝점 sampled 속도 차이도 공간 비교에서
최대0.9408\%다. 끝점 sampled 진단과 전체 면적 norm의 시간 상한은 구분한다.
P3--P3/P3--spline/spline--spline 교차 면적은 공통 overlay의 Duffy4/6/8로 계산했다.
두 대각선과 knot를 공통 분할하며 tensor의 축별 degree를 전체 degree와 혼동하지 않는다.

\paragraph{범위·에너지·독립 재실행.}
실제 검사 지점 최대 변위는0.109697 mm다. P3 Bernstein/spline control point의 convex hull과
전체 modal 진폭을 결합한 모든 위치·시간의 변위 상한은0.421778 mm 미만,
기울기 norm 상한은0.000791 미만으로 설정 범위를 지켰다.
Modal step work 잔차는 \(7.523\times10^{-24}\) J 미만, 누적 work/reset ledger는
\(3.309\times10^{-23}\) J 미만이다. 직접 consistent \(M,K\) 이차형식과의 에너지 차이는
float64 상쇄를 포함해 \(3.763\times10^{-17}\) J 미만이며 이를 앞의 잔차와 혼동하지 않는다.
세 reset에서 제거한 n16 kinetic은0.508396, 0.878272, 1.205717 nJ다.
Verified/replay CPU91.145/91.445 s의 \textbf{618개 배열/44,291,978개 scalar와 물리 report가 정확히 일치}했다.
각 run의 manifest 제외 inventory는54파일/167,594,035 byte다. 앞선77.781 s initial run은
추가 envelope/직접 ledger/독립 expm 검사 이전의 진단 자료로 보존한다.

\paragraph{실제 발행 sample과 채택 경계.}
\code{wind3dgs.validated_p3_patch_sample.v1}로19 window$\times$4 patch=\textbf{76 sample}을 생성했다.
자연 연속7개와 세 reset의6/4/2개 window, 각12 interval/13 state/9 probe다.
원래 displacement/velocity·초기 상태·절대 clock·wind·전역 work·coupled context와 source group을 보존한다.
Natural 마지막6 interval의 불완전 tail, 분기 prefix/replay 및 개입을 가로지르는 interval은 제외한다.
전체25 probe의 네 겹침 patch 가중치는 총면적1 m$^2$/label mass measure0.1 kg을 보존한다.
이 positive measure를 consistent \(M\)의 대체로 쓰지 않으며 전역 work를 patch별로 합산하지 않는다.
별도 프로세스의 raw 상태/map/풍장/work 대조와 NumPy loader가 통과했고,
batch8의10개 batch(마지막4개), 응답 shape \([8,13,9,3]\)를 확인했다.
신규11개 검사는 경계·signed map·독립 시간 기준·에너지·convex hull·원본 연결·변조/승격/덮어쓰기 거부와
SciPy/Newton/Warp/Torch 없는 reader를 검증했다.

\code{sample_scope_eligible=true}는 이 작은 굽힘 조건만 뜻하며
\code{canonical_training_eligible=false}, \code{r1_complete=false}를 유지한다.
큰 변형·새 object/seed·장기 실행·independent 3DGS/oracle·student 학습 이득은 검증하지 않았다.
Raw와 실제 sample은 각각 \code{experiments/artifacts/runs/teacher_velocity_reset/} 및
\code{experiments/artifacts/datasets/teacher_p3_reset_samples/20260909_small_bending_v1/}에 보존한다.
정확한 명령·report/hash·당시 source ZIP·대표 그림은
\code{experiments/R1_teacher_velocity_reset/p3_samples/README.md}를 따른다.

\subsection{현행 구조 검증 후보: P3 판과 독립 spline}
\label{sec:r1-current-plate}

\begin{decisionbox}{현행 후보와 채택 범위}
공간 보완의 현행 후보는 P3 \(C^0\) interior-penalty 선형 판이다. 아래 식과 고정 penalty로
처방 압력 fixture를 검사했다. 이 선형 진단만으로 후속 선택된 nonlinear membrane+bending/wind Teacher의 수렴을 보장하지 않는다.
후속 작은 굽힘 actual wind/reset의 별도 고정 경계와 채택 범위는 \S\ref{sec:r1-p3-wind-samples}를 따른다.
Native hinge와 quadratic patch를 대신 검토하는 이유는 뒤의 교체 근거에 보존한다.
\end{decisionbox}

Engineering curvature를 \(\kappa=[\kappa_{xx},\kappa_{yy},2\kappa_{xy}]^\top\)로 두고,
SI elastic parameter의 KL constitutive matrix를 다음과 같이 사용한다.
\begin{equation}
 D=\frac{Eh^3}{12(1-\nu^2)},\qquad
 D_b=D\begin{bmatrix}1&\nu&0\\\nu&1&0\\0&0&(1-\nu)/2\end{bmatrix}.
 \label{eq:r1-plate-constitutive}
\end{equation}
\(D\)는 N\,m이고 물리 두께 \(h\)는 m다. Quadrature/DOF mass는
식~\eqref{eq:r1-consistent-mass}로 구분한다.

독립 기준은 unit square의 tensor-product \(C^2\) cubic B-spline, 4/8/16/32 spans,
Gauss4 tensor quadrature의 \(H^2\)-conforming KL 판이다. Consistent \(M,K\)를 조립하고
x0 첫 coefficient만 제거해 slope-free 경계를 유지한다. Quadratic 초기 상태는 Greville collocation으로
정확히 표현하고 generalized eigenbasis 전체를 사용한다.
Simply-supported analytic frequency를 별도 구현 검증으로 사용했으며 주 비교의 경계를 바꾸지 않았다.

P2는 triangle당 6 DOF, P3는 10 DOF의 Lagrange shape다. P3는 local Vandermonde로 shape/gradient/Hessian을
계산한다. Volume KL energy에 interior edge moment consistency와 slope-jump penalty를 더한다.
\begin{equation}
 E=\frac12\sum_t\int_t\kappa^\top D_b\kappa\,dA
 -\sum_{e\,\mathrm{interior}}\int_e[\partial_n w]\langle M_{nn}\rangle\,ds
 +\frac12\sum_{e\,\mathrm{interior}}\int_e\frac{\alpha}{h_{\rm avg}}[\partial_n w]^2\,ds.
 \label{eq:r1-c0ip-energy}
\end{equation}
% [수식 변경 이유 -- 2026-09-09]
% Curvature fit의 정적 에너지 일치만으로 interior weak-force consistency가 보장되지 않았다.
% 현행 공간 후보에는 interior moment consistency와 slope-jump penalty를 함께 둔다.
% 경계는 position-only/slope-free다. Exterior rotation penalty를 추가하지 않는다.
Normal orientation과 jump/average는 source에 고정돼 있다.
\(h_{\rm avg}\)는 이웃 cell diameter의 평균이며 물리 두께와 다른 길이다.
P2는 \(\alpha=Eh^3\), P3는 \(\alpha=Eh^3(p/2)^2\), \(p=3\)의 고정 factor 2.25다.
P3 결과를 본 뒤 penalty를 sweep하거나 유리한 값을 골라낸 것이 아니다.
P2의 volume/edge는 Duffy3/Gauss2, P3는 Duffy4/Gauss3를 사용한다.
Exterior rotation penalty는 추가하지 않았다.

식의 reference는 선행 구현 기록에 연결한 공식 FEniCS-Shells Kirchhoff--Love demo다.
해당 demo의 clamped 경계를 이번 position-only 경계로 복사하지 않았다.
Reference URL은 \code{experiments/R1_teacher_spatial_remediation/README.md}, 당시 적용 판단은
\code{code/sessions/2026-09-08_07_teacher_spatial_remediation.md}에 보존돼 있다.
이번 문서 통합에서는 외부 문헌을 새로 조회하지 않았으며 신규성 조사를 수행한 것으로 표시하지 않는다.

P2/P3는 constant-curvature interior force가 각각 약 \(10^{-15}/10^{-14}\) N,
unconstrained affine null 3개/left-pin null 1개, proper rotation, gradient, area/mass,
patch polynomial force/energy를 검사했다.
Independent oscillator matrix exponential에는 zero mode와 resonance도 포함했다.
P3 n16 forward frequencies는
\([6.35772073,14.26003426,24.28380007,24.88148287,46.36581911]\),
spline32는 \([6.35770535,14.25999311,24.28325170,24.88117461,46.36373132]\) rad/s다.

\subsection{현행 CPU 시간 적분: acceleration-form Newmark}
\label{sec:r1-current-integrator}

\textbf{적용 범위:} 기존 quadratic-patch/lumped-mass CPU shell의 정밀도 보정 구현이다.
최종 P3 backend의 temporal acceptance는 별도로 필요하다.

\(\beta=1/4,\gamma=1/2\)인 같은 Newmark 식을 acceleration unknown으로 다시 풀었다.
\begin{equation}
\begin{aligned}
 x_{n+1}&=x_{\rm pred}+\beta\Delta t^2a_{n+1},\qquad
 v_{n+1}=v_{\rm pred}+\gamma\Delta t a_{n+1},\\
 R(a_{n+1})&=Ma_{n+1}+\nabla E(x_{n+1})-f_{n+1}.
\end{aligned}
\label{eq:r1-acceleration-newmark}
\end{equation}
% [수식 변경 이유 -- 2026-09-09]
% 위치 차분을 dt^2로 나누는 residual 평가를 acceleration unknown의 식으로 교체했다.
% 같은 Newmark beta/gamma와 물성·tolerance를 유지한다. P3 backend의 적분 검증으로 승계하지 않는다.
Undamped tangent는 \(M+\beta\Delta t^2H\)다. 가까운 두 position을
\(\Delta t^2\)로 나누는 잔차 평가를 피했다. 물리 stiffness/damping, tolerance를 바꾸거나
고주파 필터·숨은 viscosity를 넣지 않았다. 기존 실패를 재현하고 새 3,457 step을 완료했다.
Short 차이는 2.20620\%로 개선됐지만 full N2560 11.7491\%는 남아 시간 해상도 문제임을 확인했다.



\subsection{구조 후보의 교체 이유와 이전 식의 적용 범위}
\label{sec:r1-rejected-structures}

\textbf{변경 근거:} 다음 식과 결과는 native hinge, 면적 보정 hinge 및 quadratic patch를
최종 Teacher로 채택하지 못한 이유다. 새로운 구조식은
\S\ref{sec:r1-current-plate}에서 직접 제시했고, 여기서는 이전 식의 단위·반례·실패를 보존한다.

\subsubsection{고정 native bending 계수}

설치된 Newton 1.3.0의 에너지 의미를 독립 NumPy로 계산했다. Flat XY/XZ square에
\(w=A(X/W)^2\), \(W=H=1\) m, \(A=0.01\) m, \(k_e=10\) N와
\(n=4,8,16,32\)를 사용했다. Boundary edge는 bending 합에서 제외한다.
\[
 E_b=\frac12 k_e\sum_{e\text{ interior}}l_e^0\theta_e^2,
 \qquad K_E=\frac{2E_b}{A^2}.
\]
\(K_E\)는 에너지 등가 scalar이지 nonlinear force/A나 tangent HVP가 아니다.
Strip 기울기 \(s_j\)의 해석식
\[
 E_b=\frac12 k_eH\sum_j[\arctan(s_{j+1})-\arctan(s_j)]^2,
 \qquad K_{\rm small}=\frac{4k_eH}{W^2}\frac{n-1}{n^2}
\]
과 최대 상대 차이 \(1.792\times10^{-15}\)였다.
Energy는 \(0.000374911,0.000218695,0.000117157,0.000060531\) J,
\(K_E\)는 약 \(7.498,4.374,2.343,1.211\) N/m로 줄었다.
Finest/coarsest 비는 0.161454다. Native 계수를 고정하면 같은 연속 변형의 강성이 사라지는
mesh 의존성이 있으므로 그 계수 자체를 연속 판 강성으로 채택할 수 없다.
아주 작은 진폭의 float32 입력 문제가 있어 asymptotic fixture는 이진 정확 진폭 \(2^{-14}\)로
분리했다. 기본 \(A=0.01\) 결과나 tolerance를 사후 완화한 것이 아니다.

\subsubsection{Area-weighted hinge}

후보 \(k_e=B_hl_e^0/(A_L+A_R)\)에서 \(B_h\)의 단위는 N\,m다.
이는 native coefficient N과 구분되며 자동으로 isotropic plate \(D\)가 되지 않는다.
7개 변형, 두 대각선, 두 곡률, 네 mesh의 112개 사례를 계산했다.
작은 곡률의 \(\pm45^\circ\) 원통 bending에서 방향 에너지 비가 3으로 접근했고,
\(n=32\)에서 2.9375001853이었다. 한 방향에 맞춘 scalar calibration으로 해소되지 않는다.
\textbf{등방 Teacher mapping으로는 채택하지 않았다.}

\subsubsection{정적 선형 판}

\code{flat_plate_quadratic_patch_v1}은 central triangle의 edge-neighbor BFS patch를
1--4 ring에서 찾고 rank 6인 첫 design을 쓴다.
\([1,s,t,\tfrac12s^2,st,\tfrac12t^2]\)를 SVD로 fit하며 relative cutoff \(10^{-12}\),
condition limit \(10^8\)이다. Local length scale을 되돌린 curvature operator \(C\)와
engineering curvature \(\kappa=[\kappa_{xx},\kappa_{yy},2\kappa_{xy}]^\top\)를 사용한다.
당시 동일한 constitutive 식~\eqref{eq:r1-plate-constitutive}를 다음 patch 연산자에 적용했다.
\[
 E_b^{\mathrm{patch}}=\frac12\sum_t A_t(C_tw)^\top D_b(C_tw),\qquad
 f_b^{\mathrm{patch}}=-\sum_t C_t^\top A_tD_bC_tw.
\]
% [교체한 식의 기록 -- 2026-09-09]
% 이 patch energy/force는 실패 분석을 재현하기 위한 이전 후보다.
% 현행 공간 보완 후보의 식은 eq:r1-c0ip-energy다. 이전 static pass를 동적 pass로 읽지 않는다.
HVP는 이 선형 operator에서 정확하다. \(\nu=0,0.3\), 12방향 cylinder와 twist/dome/saddle,
quartic/sine, \(n=4\ldots32\)의 408개 정적 사례가 당시 기준을 통과했다.
Affine null 3개와 일반 변형 finest energy 오차 0.77\% 미만을 확인했다.
후속 공간 진단은 이 정적 통과가 동적 consistency를 보장하지 않음을 보여 주었다.

\subsubsection{3D 구조 연산자}

\code{flat_shell_stvk_projected_curvature_v1}은 StVK membrane의 Green strain과
current central normal에 투영한 bending을 합친 CPU 진단이다.
정확한 gradient/HVP에 normal derivative와 geometric-stress 항을 포함했다.
Finite area는 \(10^{-8}\)보다 커야 하고 proper rigid rotation을 단순 normal-flip으로 거부하지 않는다.
Rest Hessian의 PSD와 일반 current tangent의 가능한 indefiniteness를 구분한다.

\(\nu=0,0.3\), \(h=0.01,0.001\) m, 네 mesh의 1,200개 prescribed-geometry 사례에서
objectivity, rigid motion, gradient/HVP finite difference와 선형 극한은 통과했다.
그러나 8개 total-energy ladder는 \(n=16\to32\)에서 오차 감소에 실패했다.
Membrane/bending 오차 상쇄와 thin coarse mesh의 artificial membrane energy를 따로 남겼다.
대표 thin case의 ratio는 n4에서 약 185, n32에서 0.0447이었다.
이 현상만으로 뒤의 rest-normal checkerboard 결함을 설명하지 않는다.

\subsubsection{Quadratic patch의 내부 weak-force 결함}

상수 quadratic curvature를 넣어도 checkerboard mesh의 안전한 내부 vertex에 약
\(9.5238095238\times10^{-4}\) N의 force가 n8/16/32에서 남았다.
경계에 닿는 patch의 영향을 받는 모든 vertex를 제외해 safe interior 9/121/729개를 검사했다.
Forward/backward는 약 \(10^{-17}\)--\(10^{-15}\) N의 roundoff 수준이었다.
Polynomial curvature/energy 재현과 그 adjoint가 만드는 약형 내부 힘의 consistency는 다른 조건이다.

Safe interior에만 \(O(h_{\rm mesh}^2)\) parity perturbation을 가해 에너지를 최소화하면
8.1794\%, 21.6639\%, 31.0786\% 감소했다. Amplitude는
\(3.495\times10^{-6},6.886\times10^{-7},1.640\times10^{-7}\) m였다.
작아지는 요철로도 mesh-dependent energy를 낮출 수 있는 반례다.
Ring을 2--4로 늘려도 잔류 force가 각각 약
\(2.97\times10^{-5},8.64\times10^{-5},6.11\times10^{-5}\) N로 남았다.
Patch ring 확대나 mass 교체만으로 해결됐다고 판단하지 않았다.

Lumped mass를 consistent mass로 바꾼 checker n16의 두 번째 positive frequency는
11.13933에서 11.239343 rad/s로 바뀌지만 독립 기준 14.259993과의 차이가 남는다.
Old n16의 첫 5개 positive frequency는 다음과 같다 (rad/s).
\[
\begin{aligned}
\text{forward}:&\quad[6.33265,14.18110,24.15372,24.60083,45.84505],\\
\text{checker}:&\quad[6.20080,11.13933,20.19455,21.38170,36.04130].
\end{aligned}
\]
Stiffness consistency와 mass approximation을 구분한 근거다.

\section{산출물}
\label{sec:r1-chapter-6}
\label{sec:r1-artifacts}

\begin{itemize}
  \item Versioned \code{TeacherPhysicsRegistry}와 material/traction/solver/refinement config.
  \item Converged \code{TeacherPackage}: topology hash, quadrature, trajectory, applied traction/force/work,
  spatial/temporal convergence report와 replay command.
  \item Independent \code{StaticGSPackage} 두 개 이상과 source-object group/provenance.
  \item \code{CanonicalProbeMap}: probe payload, \(S_T,S_r\), common-valid mask, mapping report/noise floor와 hash.
  \item Traction/exact-ZOH analytic fixture report와 method-identity equality report.
  \item GS별 corrected \code{OracleGlobalPackage}와 oracle-measure ID/hash.
  \item Oracle chosen path, fit/re-diagonalization report와 paired probe error.
  \item \code{BaseTransportReport}: base mean/covariance transport fixture, rendered smoke output와 failure denominator report.
  \item \code{TeacherSpectrumStatus}: converged 여부, stable peak available/unavailable reason와 extraction-config hash.
  \item \code{R1PreflightReport}: 모든 threshold ID, command/source/config/artifact hash와 R2 진입 판정.
\end{itemize}

Teacher raw topology/correspondence와 oracle fitting payload는 training-side artifact다.
이를 target \code{ResponsePackage}에 복사한 산출물은 유효한 R1 결과가 아니다.

\subsection{구현 계보와 provenance}

Workspace 상대 경로 \code{code/}, \code{ideas/}, \code{experiments/}는 독립 저장소다.
GPU checkpoint는 code \code{7010682}, experiments \code{61d972e}에 이미 보존돼 있었다.
그 뒤 구조·수렴·샘플·보완 작업은 code \code{894993c}, experiments \code{c7f4c74}에 보존했다.
최종 commit 연결은 양쪽 \code{sessions/2026-09-09_01_teacher_checkpoint.md}와
ideas \code{sessions/2026-09-09_01_teacher_evidence_integration.md}를 따른다.
Run 생성 시에는 dirty source의 파일별 content hash가 실제 executable provenance다.
나중에 만든 commit ID가 당시 run의 source hash를 대신하지 않는다.

연구 기록의 수치는 각 report의 단위와 분모를 따른다. 초기 GPU fine-run RMS 상대 오차와
후속 고정 \(A,A\omega_*\) scale 오차는 다른 metric이며 직접 하나의 개선율로 합치지 않는다.
각 단계의 ``관련 테스트 수''는 범위가 겹치므로 누적 합산하지 않는다.
정적 fixture, 시간 오차, 공간 오차, replay, 데이터 무결성, R1 acceptance도 별도 판정이다.

2026-09-09 checkpoint에서는 관련 14개 Teacher test module의 234개 검사를 실제 재실행해
95.133초/OK/종료 코드 0을 확인했다. Source hash 271개, compact hash 60개, JSON 91개와
Python 41개·shell 13개의 구문을 검사했고, 원본 9개 run의 inventory 7,777개/
5,283,975,129 byte가 저장된 byte/SHA-256과 일치했다.
이는 해당 checkpoint의 누적 QA·원본 검산이며 이번 문서 감사의 새 물리 실험, 추가 GPU 실행 또는
repository 전체 test discovery가 아니다. 정확한 명령과 범위는
\code{code/sessions/2026-09-09_01_teacher_checkpoint.md} 및
\code{experiments/sessions/2026-09-09_01_teacher_checkpoint.md}에 보존한다.

\subsection{재현 명령·hash와 원본 보존}

명령은 workspace root에서 실행한다. Launcher가 \code{code/}로 이동하므로 launcher 인자 경로는
code 기준이다. 모든 새 run은 사용하지 않은 output 이름을 요구한다. 기존 artifact를 덮어쓰거나
자동 재개하지 않는다. 전체 물리 재실행과 inventory/source/gate 검산을 구분한다.

\begin{verbatim}
bash code/scripts/check_teacher_gpu.sh
bash code/scripts/generate_teacher_sample_dataset.sh \
  --output ../experiments/artifacts/runs/teacher_sample_dataset/my_run \
  --dataset ../experiments/artifacts/datasets/teacher_samples/my_dataset
bash code/scripts/audit_teacher_plate_spatial_remediation.sh \
  --source-run ../experiments/artifacts/runs/teacher_shell_linear_spatial/20260908_reference_v1 \
  --output ../experiments/artifacts/runs/teacher_spatial_remediation/my_reference
bash code/scripts/audit_teacher_plate_cubic_refinement.sh \
  --source-run ../experiments/artifacts/runs/teacher_spatial_remediation/20260908_reference_v1 \
  --output ../experiments/artifacts/runs/teacher_spatial_remediation/my_cubic
bash code/scripts/audit_teacher_plate_cubic_refinement.sh \
  --verify ../experiments/artifacts/runs/teacher_spatial_remediation/my_cubic
\end{verbatim}

전체 시간 refinement는 네 선행 run이 필요하다. 정확한 인자 및 재검산 명령은
\code{experiments/R1_teacher_shell_full_refinement/README.md}에 있다.
다른 단계의 launcher/인자도 아래 evidence index의 README가 소유한다.
P3의 \code{--verify}는 inventory/current source/판정 검산이며 전체 modal 물리 재계산은
별도 fresh run으로 실행했다.

\begin{longtable}{L{0.27\linewidth}L{0.65\linewidth}}
\toprule Artifact & SHA-256 (전체 값)\\\midrule
Full temporal report & \code{14700c5346b3f9502f724abd32febcc7a7a6266771ab08fbb750dae531e985f5}\\
Linear spatial report & \code{3a0d8e7ee5be505b190008d387c67a113a027ce0651595cb43b0812c8c33ba1a}\\
Dataset manifest & \code{1ab229207fa8ec60c282aa433c49d7a5f8a288deb1575688794b72b2afdae8f5}\\
Remediation reference report & \code{b5c3929f2e65160f9f8127fe4ea9612e41a0871a5554c4efd5bcec0e77249fd7}\\
P3 report & \code{799a48aa00e9eb40085ad030298d90a409959332092fbfdc784b312bc2cb46ac}\\
\bottomrule\end{longtable}

Remediation raw root는 \code{experiments/artifacts/runs/teacher_spatial_remediation/}다.
Reference는 \code{20260908_reference_v1}와 \code{20260908_reference_v1_replay},
P3는 \code{20260908_cubic_v1}와 \code{20260908_cubic_v1_replay}다.
Reference run은 각각 45.784/41.826 s, inventory 55개/약 45.16 MB;
P3 run은 104.291/106.025 s, inventory 16개/146,317,681 byte다.
각 pair는 runtime 외 각각 54/15개 inventory가 byte-identical이다.
모델 배열과 전체 comparison curve는 독립 계산으로 재생성했다.
모든 native nodal trajectory를 따로 저장한 것은 아니며 보존한 modal model과 계수로 재구성한다.
Preceding spatial source 30개와 새 source snapshot, config/environment 및 모든 inventory hash를 연결한다.

Git에는 compact report/config/environment/provenance/verification과 선택한 PNG를 보존한다.
Ignored raw의 모든 state/trace/vector, modal arrays 및 dataset NPZ는 Git clone만으로 복구되지 않는다.
Raw artifact를 별도 보존하거나 기록된 source/config로 재생성해야 한다. Compact evidence만으로
NPZ가 복원된다고 주장하지 않는다. 개인 절대 경로·secret은 기록에서 제외한다.

\subsection{상세 기록 index와 논문 작성 시 사용법}

Code 기능별 설계·승인·수정·검증 과정은 아래 session에 남아 있다. 날짜의 각 번호는 해당 repository
내 순서다. 2026-09-01 viewer 기록은 인수한 선행 맥락이며 이후 작업과 구분한다.
\begin{longtable}{L{0.20\linewidth}L{0.70\linewidth}}
\toprule Code session 날짜/번호 & 소유 내용\\\midrule
09-06 01, 02 & Dataset handoff; TeacherPhysicsRegistry native SI/schema/실제 model 검사\\
09-07 01, 02, 03 & Trajectory writer/reader/replay; wind sequence; authored-rest 초기 변위\\
09-07 04, 05, 06, 07 & Common probe/adjoint; convergence comparator; GPU launcher; GPU checkpoint\\
09-07 08, 09, 10 & Native bending 감사; area-weighted hinge; quadratic patch plate 대안\\
09-07 11, 12 & 3D shell structure; CPU nonlinear dynamics\\
09-08 01, 02, 03, 04 & Temporal reference/실패; acceleration precision; short 및 full refinement\\
09-08 05, 06, 07 & Linear spatial 실패; sample 생성·검증; stencil/P2/P3 보완\\
09-09 01 & 누적 코드 회귀·Git checkpoint\\
09-09 02 & 변화 바람/velocity-reset의 prefix replay·공력·개입 검산과 해상도 실패\\
\bottomrule\end{longtable}

\code{experiments/} 아래 다음 각 폴더의 \code{README.md}에는 해당 입력, 결과 표, raw 경로,
명령, source/config hash와 재검산 근거가 있다. 별도의 experiment session은 실행 과정과 실패를 보존한다.
\begin{longtable}{L{0.47\linewidth}L{0.43\linewidth}}
\toprule Evidence folder & 논문에서 참고할 근거\\\midrule
\code{R1_teacher_velocity_reset} & 25개 trace 검산, 속도 수렴 실패와 회전 위주 상태 반례\\
\code{R1_teacher_smoke} & 사용자 GPU 실행과 물리 수렴 미판정\\
\code{R1_teacher_bending_audit} & Native coefficient의 mesh scaling 반례\\
\code{R1_teacher_bending_mapping} & Area hinge의 방향 비등방성\\
\code{R1_teacher_plate_reference} & Static KL polynomial/방향 검증의 범위\\
\code{R1_teacher_shell_structure} & Objectivity/미분, 오차 상쇄와 thin locking 진단\\
\code{R1_teacher_shell_dynamics} & 수치 계약 통과와 초기 시간 실패\\
\code{R1_teacher_shell_temporal} & 독립 DOP853, 고주파와 line-search 실패\\
\code{R1_teacher_shell_precision} & Acceleration unknown의 정밀도 보정 근거\\
\code{R1_teacher_shell_refinement} & Short 시간 해상도 ladder\\
\code{R1_teacher_shell_full_refinement} & Full 시간 수렴·전체 trace 검산\\
\code{R1_teacher_shell_linear_spatial} & Exact temporal response에서도 남는 공간/방향 실패\\
\code{R1_teacher_sample_dataset} & 실제 15 window, source/replay/loader 검증\\
\code{R1_teacher_spatial_remediation} & 내부 force 반례, 독립 spline, P2/P3와 입력별 차이\\
\bottomrule\end{longtable}

논문 method에는 mass/shape/quadrature, attachment와 force sampling을 정확히 옮기고,
appendix에는 실패한 native/area-hinge/patch와 보완 근거를 선택해 인용할 수 있다.
Results에서는 실행 성공, numerical convergence와 accepted Teacher를 분리하고 normalized error의
분모와 측정 시각/연속 상계 여부를 함께 쓴다. Limitations에는 broad-band displaced response,
비선형 wind의 큰 변형/장기 검증 미완료, 독립 GS/oracle 및 R1 전체 미완료를 남긴다.
이 기록 자체는 아직 수행하지 않은 learner 실험이나 physics validation의 대체 증거가 아니다.

\section{필수 fixture와 metric}
\label{sec:r1-chapter-7}
\label{sec:r1-fixtures}

\subsection{Teacher 가속 개발의 제한 종료와 후보 동결}
\label{sec:r1-bounded-closeout}
2026-09-16 사용자 결정으로 전체 FP32 및 새 solver 개발을 종료하고, 완료된 동일 구간의
대조만 근거로 다음 실험 후보를 동결한다. RTX5070의 정상 일반 바람/HL00은
M2와 force block 32/32/256, GTX1080Ti의 일반 바람/HL00은 M1과256/256/256이다.
GTX1080Ti의 C0 및 알려진 HL01은 R64를 유지한다. 순서는 volume/interior edge/boundary edge다.
M1은 보조 행렬 분해·apply만 FP32이며, M2는 큰 내부 선형 연산에 FP32를 사용하되
FP64 master와 원래 A64 참 잔차를 유지한다. 비선형 승인·상태 hi/lo·독립 검산과
Newmark/기존 Gauss retry 정책은 바꾸지 않는다.

5070 성공 근거는 driver API13000에 한정된다. 현재 읽기 전용 조회의 API13040과 다르며,
후속 Graph 충돌의 해결이나 현재 환경 재현 성공을 주장하지 않는다. 인프라 복구는 별도 작업이다.
5070 W30은 완료2쌍만 집계하며, 추가 M2 실행과 미완료 prefix로 빠진 R64를 대체하지 않는다.
겹치는 W1/W30/HL00 시간은 합산하지 않고, 검산 포함 생성 시간과 process wall을 분리한다.

추가 검증 범위는1080Ti M1 W30의 full/summary 기록 각1회 및 저장된 HL01 선형계의
R64/F64\_fresh 비교로 제한한다. 후자는 원본 frame225이며 최고비용 frame236 검증이 아니다.
기존 warm-up 뒤3회에서 조립·분해·풀이 총비용과 원래 A64 잔차를 비교하고,
상주한 old factor의 풀이 비용도 따로 표시한다. 총시간20\% 감소는 개발 자원 배분 기준이며
정확도 허용오차가 아니다. 통과해도 \code{promising_frozen_candidate}에서 멈춘다.
전역 P 갱신 정책이나 전체 HL01 재생으로 확대하지 않는다.
제한 시험은 두 기록 모드의 W30 검산과 선형계 각3회 검증을 완료했다.
두 기록 모드 사이에 비영 수치 차이가 있어 회귀 동등성은 미판정으로 유지한다.
저장된 frame225의 F64\_fresh는 원래 잔차와 비용 기준을 통과하여
\code{promising_frozen_candidate}로만 남겼으며 HL01 유지 선택은 여전히 R64다.

선택 manifest, 완료 쌍과 두 제한 시험의 실제 결과·미측정·hash·명령은
\code{experiments/R1_teacher_velocity_reset/timestep_search/precision_v3/closeout_report.md}가 소유한다.
\code{production_enabled=false}, \code{training_eligible=false}와 회귀 예산 미정 상태를 유지한다.
제한된 개발 검산 통과는 R1의 공간·시간 수렴, 장기 teacher 적격성이나 연구 Gate 완료가 아니다.

2026-09-17 후속 국소 진단은 동일 frame225 snapshot 및 저장된 두 보정만 재사용했다.
초기 FP64 RHS가 정확히 재현됐으며 기존 trial/evaluate/line\_decide를 유지했다.
R64는 전량 보정이 승인됐지만 fresh는 전량 보정이 거부되고 세 번 줄인 보정에서만 승인됐다.
fresh의 승인 후 비선형 잔차 감소는 작으므로 frozen linear의 가속을 Newton 전체 가속으로 승계하지 않는다.
후속 판정은 \code{line_search_checked_frozen_candidate}이며 완전한 substep Newton/audit는 미검증이다.
세부 근거는 \code{experiments/R1_teacher_velocity_reset/timestep_search/precision_v3/frozen_line_search_report.md}에 둔다.
HL01 운영 R64와 생산/학습 적격성 보류, 전체 FP32 개발 종료는 유지한다.

사용자 후속 결정으로 이 추가 최적화 분기도 종료한다. Frozen linear의 원래 A64 잔차 통과와
국소 가속, line search의 $\lambda=1/8$ 승인, 비선형 잔차 감소(R64 94.07\%, fresh 1.14\%)를
별도 판정으로 유지한다. 전체 Newton/substep/독립 audit는 미검증이고 운영 채택은 보류다.
이를 fresh의 수치 실패나 전체 성능 저하 확정으로 해석하지 않으며 선형 가속을 전체 teacher
가속으로 보고하지도 않는다. Summary/full 기록 옵션을 유지하고 추가 HL01 재생,
frame236 checkpoint, 전역 rebuild, fresh P32, line-search 계수 변경 및 새 solver로 확대하지 않는다.
5070 Graph 장애와 학습 적격성 정확도 예산은 별도 미해결 항목이다.
최종 종료 상태와 원본 CSV/NPZ/hash 연결은
\code{experiments/R1_teacher_velocity_reset/timestep_search/precision_v3/fresh_branch_closeout.md}에 둔다.

2026-09-17 사용자 요청으로 동결된 M1/M2/R64 선택을 GPU 자동 탐지 세 씬 실행에 연결했다.
기존 굽힘1/500 입력과 Newmark/Gauss retry·독립 검산은 유지하고,1080Ti의 알려진 HL01은
R64를 유지한다.5070의 환경 차이는 실행 보류로 분리하며 생산/학습 적격성은 승격하지 않는다.
제한 smoke와 세 씬 본 실행 준비는 전체 가속·장기 정확도 완료가 아니다.
설정·명령·검증 근거는
\code{experiments/R1_teacher_velocity_reset/timestep_search/gpu_auto_scenes/README.md}를 따른다.

2026-09-18 사용자 요청으로 기존 Newmark 절반 $\Delta t$ 복구를 GPU 자동 선택의 별도 실행 옵션으로 연결했다.
기본 M1/M2/R64 선택은 유지하되 실패한 기본 구간만 FP64 Newmark 두 half step으로 재계산한다.
기존 검산·허용오차와 실패 복원은 유지하며 재귀 세분화나 추가 Gauss 복구는 하지 않는다.
제어 경로 검사와 본 고하중/성능 검증을 구분하고 학습 적격성은 승격하지 않는다.
실행 근거는 \code{experiments/R1_teacher_velocity_reset/timestep_search/gpu_auto_scenes/half_retry.md}에 둔다.

별도 국소 비교에서는 Newmark 수렴 실패 뒤 FP64 half2를 먼저 시도하고, 그것도 수렴하지
않으면 원래 승인 상태로 복원하여 Gauss6차8분할을 수행한다. 실패한 half의 부분 진행은
채택 궤적에 섞지 않으며 물리/비유한/독립 검산 오류를 다음 적분기로 우회하지 않는다.
저장된 프레임의 동일 시작 상태와 외력에서 기존 직접 Gauss 복구와 비교하며,
수렴 및 독립 검산 통과와 두 적분 정책의 끝 상태 차이를 분리한다.
국소 시간 감소를 동일 정확도의 가속이나 전체 teacher 적격성으로 승계하지 않는다.
기존 실행 기본값은 유지하며 실제 표본 범위와 원본은
\code{experiments/R1_teacher_velocity_reset/timestep_search/cascade_retry/report.md}를 따른다.

2026-09-20 후속 성능 표본에서 저부하 직접 Gauss는 Newmark보다 느렸고,
반복 복구가 발생한 고부하에서는 프레임 전체 Gauss 재계산이 빨랐다. 이 결과를
고정 frame 번호가 아닌 (1) 기존 code1/2/3의 복구 가능한 수치 실패, (2) 첫16 Newmark
substep 안의 실패3회, (3) 현재 장치에서 실측한 Newmark 잔여 시간이 전체 Gauss
예상 시간보다 큰 조건으로 연결한다. 세 조건을 모두 만족해도 프레임 시작의
FP64 hi/lo 상태로 복원하여 Gauss6차512단계를 다시 풀며, 버린 Newmark/복구 비용도
실행 시간에 포함한다. 수치 검산·기하·에너지 실패는 적분기 전환으로 우회하지 않는다.

동일 Gauss512 고부하3프레임에서 RTX5070 mixed32는 R64 대비22.83\% 짧았고
저부하도13.92\% 짧았으며 검산 실패와 FP64 복귀는0이었다. GTX1080Ti의 mixed32는
저부하3.59\% 느리고 고부하 합계 이득2.66\%가 프레임별로 혼재했다. 따라서 바람
Newmark는 각각 M2/M1을 사용하되 직접 Gauss는 RTX5070 mixed32, GTX1080Ti R64로 나눈다.
Mixed32는 FP64 master·원래 A64 참 잔차·비선형 승인·최종 상태/에너지·독립 검산을
유지한다. 이는 세 장면 후보 실행 계약이며 긴 궤적 teacher 적격성은 미검증이다.
세부 선택·실행 근거는
\code{code/docs/gpu_runtime_selection.md}와
\code{experiments/R1_teacher_velocity_reset/timestep_search/gpu_auto_scenes/adaptive_integrator.md}가 소유한다.

\subsection{3단계6차 Gauss의 GPU 상주 개발 비교}
2026-09-14 사용자 요청으로 기존 Newmark와 별도의 GPU Gauss 시험 경로를 구현했다.
3개 내부 단계를 결합하며, 과거 파일명 Gauss6의6단계12차 후보와 구분한다.
단계별 원래 힘/HVP와 consistent mass를 유지하고 Newton/GMRES/line search,
보조 행렬 갱신과 분해, hi/lo 상태 갱신을 GPU graph에서 수행한다.
공통 접선을 쓰는 단계 결합 보조 풀이만 실수 block으로 표현하며 실제 물리 연산자를 바꾸지 않는다.

기존 설정의 GPU Newmark, 세분화 GPU Newmark, CPU 중심 Gauss 및 GPU Gauss를
동일 저장 상태와 고정 외력으로 대조한다. 초기 준비, 구간 풀이 완료, 기록 replay,
독립 검산 비용을 분리한다. Gauss 독립 검산은 별도 CPU longdouble 식과
P3$\times$cubic 구간 기하를 기준으로 GPU의 독립 buffer/힘 재평가에도 연결했다.
검산은 본 계산의 잔차 cache를 재사용하지 않으며 CPU 판정과 손상 입력 거부를 대조한다.
명시한 수치 재현성 기준을 검사하고 bitwise 일치를 일반 계약으로 요구하지 않는다.
원본, 세부 오차와 비용, 실패 분모 및 적용 범위는
\code{timestep_search/cloth_coarse/gauss_gpu_comparison.md}가 소유한다.
후속 동일 국소 구간의 시간 간격 탐색에서는 단계 수 증가보다 반복 수 감소가 커
Newmark의 총 풀이 비용이 낮아지는 조건을 확인했다. 빠른 반복 수렴과
정밀 참조 대비 움직임 차이는 별도로 판정한다. 인접 시간 분할의 차이와
GPU 검산을 포함한 비용은 \code{cloth_coarse/gauss_timestep_sweep.md}를 따른다.
사용자 선택의8분할 샘플은 세 씬의 중력 준비와 무풍/바람 분기에 별도로 연결한다.
프레임 외력은 기존60Hz 정책을 유지하고 모든 적분 단계를 GPU에서 검산한다.
이 샘플의 raw hi/lo 상태 기록은60Hz 프레임 경계이며 모든 Gauss stage를 보존하는
teacher 원본 형식과 구분한다. 실행/기록 계약과 검증 범위는
\code{cloth_coarse/gauss6_bend500_three_scenes.md}를 따른다.
이는 개발 비교이며 기본 Newmark/FP64 hi/lo, 물성, 학습 label과 Gate를 변경하지 않는다.
전체 궤적 시간 수렴과 공간/GS 매핑, 접촉 및 학습 적격성은 미완료다.

\subsection{Gauss FP32 hi/lo 허용오차 개발 진단}
2026-09-15 동일 저장 상태와 프레임 외력, 같은 시간 간격에서 Gauss의 GPU 힘/HVP,
반복 풀이와 상태 갱신을 FP32 hi/lo 후보로 비교하는 별도 runtime을 연결했다.
반복 종료 기준 완화와 정밀도 변경의 효과는 같은 허용오차의 FP64 대조군으로 구분한다.
원래 FP64 식의 독립 검산은 모든 저장 stage를 다시 평가하고 초과를 기록한다.
약한 외력이 허용오차 아래로 사라진 결과와 실제 반복 풀이 실패는 유효한 가속으로 세지 않는다.
단일 프레임의 움직임 차이가 작다는 근거로 장기 안정성이나 학습 label 적격성을 선언하지 않는다.
설정, 실패 분모, 속도와 오차 및 측정 범위는
\code{cloth_coarse/gauss_fp32_trial.md}가 소유한다.
이는 개발 후보이며 공식 검산 허용오차, 기본 FP64 hi/lo와 Gate를 변경하지 않는다.

후속 행·열 스케일링은 실수 stage basis의 FP32 보조 행렬과 RHS/해를
일관되게 변환하며 원래 물리 연산자와 잔차 기준은 유지한다.
별도 보정 후보는 FP64 상태·힘·참 선형 잔차와 FP32 HVP/GMRES/보조 풀이를 결합한다.
작아진 보정 RHS의 정규화/역변환을 수행하고, 원 단위의 엄격 기준을 만족하지 못하면 실패로 기록한다.
스케일링 단독과 보정 추가의 비용/효과는 별도 lane에서 비교한다.
검증은 독립 FP64 식과 CPU longdouble fixture를 따르며,
구현 및 비교 분모는 \code{cloth_coarse/gauss_scaling_trial.md}가 소유한다.
이 후보의 구현과 국소 검증은 장기 시간/공간 수렴 또는 teacher 채택을 뜻하지 않는다.

\subsection{Newmark 우선과 Gauss 원상 복원 재계산의 개발 후보}
2026-09-15 사용자 승인으로 두 FP64 hi/lo 적분기를 프레임 경계에서 선택하는 별도 경로를 구현했다.
기본 Newmark64분할과 같은 원본의128분할을 각각 원래 식으로 독립 검산하고,
공통 중간 시각의 consistent mass 가중 위치/속도 차이를 전환 지표로 사용한다.
지표 초과 또는 풀이/검산 실패 시 원래 프레임 시작 상태와 같은 held 외력에서
Gauss6차512단계를 다시 계산한다. Gauss 자체 검산을 통과한 뒤에만 상태를 채택하며,
실패한 Newmark 상태에서 이어가거나 두 적분기의 검산 공식을 혼용하지 않는다.
전환 지표는 개발 설정으로서 공식 teacher 정확도 기준이 아니며, 두 저차 결과가
공통 오차를 갖는 경우를 별도 고차 참조와 비교한다. 추가 시험/검산/재계산 비용을
포함한 실측과 실패 분모는 \code{cloth_coarse/integrator_switch_trial.md}가 소유한다.
두 국소 저장 상태의 반복 비교는 원래 검산을 통과했으나 모두 고차 재계산을 선택했고,
이 표본에서는 가속 근거를 얻지 못했다. 작은 연속 fixture의 외력 변경/양방향 적분기 전달과
원상 복원 검증을 전체 생성기의 장기 검증으로 승계하지 않는다.
기하 투영 충분조건의 실패와 시간 적분 오차를 구분한다. 기본 솔버, 물리식, force clock,
R1 완료 체크와 Gate는 유지하며 연속 궤적/학습 적격성은 미완료다.

후속 사용자 확인에서는 추가128분할 시간 지표를 제외하고, Newmark64의 기존
국소 기하 검산 flag16 단독 실패에만 Gauss로 원상 복원 재계산하는 별도 시험을 요청했다.
다른 물리/솔버 실패는 재계산으로 숨기지 않고 실패로 보존한다.
이 기하 전용 정책은 시간 적분 오차를 감시하지 않으며, 앞선 시간 지표 시험과 구분한다.
기하 충분조건 실패가 고차 적분만으로 해결된다고 가정하지 않는다.
재시험과 분기 검증의 근거는 \code{cloth_coarse/geometry_switch_trial.md}가 소유하며
기본 생성기와 teacher 정확도/Gate는 변경하지 않는다.

\subsection{Newmark 수렴 실패 구간의 절반 시간 간격 재시도}
2026-09-15 사용자 요청의 별도 성능 시험은 FP64 hi/lo Newmark를 고정하고,
직사각형/손수건/삼각형 각각에서 고정 시간 간격과 수렴 실패 시 절반 간격 재시도를 비교한다.
기본 간격은1/3840초이며 정상 구간에는 추가 시간 오차 추정을 수행하지 않는다.
수렴 실패 직전 정상 상태와 프레임 유지 외력을 복원해1/7680초 두 단계로 대체하고,
독립 검산 통과 후 다음 기본 구간으로 진행한다. 절반 간격에서도 실패하면 중단한다.
기하 충분조건 경고나 비유한/물리 검산 실패를 이 재시도로 숨기지 않는다.
실패 시도와 재시도 비용을 포함한 실행 설정/분모는
\code{cloth_coarse/newmark_dt_suite.md}가 소유한다.
이 시험은 수렴 실패 복구용이며 정상 수렴 구간의 시간 정확도 인증을 뜻하지 않는다.
기본 생성기, 물리식, 완료 체크와 Gate는 유지한다.

후속 사용자 선택은 위 half-dt 대신 실패한 기본 구간을 Gauss6차8단계로 복구하는
별도 \code{newmark_gauss_retry} backend다. 성공 후 다음 구간은 기본 Newmark로 돌아간다.
Gauss 단계는 자체 stage/끝점/에너지/cubic 기하 식으로 독립 검산하며,
저장된 실제 dt와 적분기 식별자를 함께 기록한다. 기본 허용오차와 Gate는 유지한다.
반복적인 전환의 비용과 전체 안정성은 국소 복구 성공만으로 판단하지 않는다.
실행/검증 근거는 \code{cloth_coarse/newmark_gauss_retry.md}가 소유한다.

\subsection{1/500 학습 정확도의 후속 시간 세분화}
2026-09-14 동일2초 seed의1프레임 Newmark16/32세분화에서도 독립 검산은 통과했지만,
인접 시간 해상도의 순간 속도 상대 질량--시간 RMS 차이는28.59\%,20.18\%로 남았다.
프레임 순이동의 평균 속도 차이가 작다는 사실을 순간 속도 label의 시간 수렴으로 해석하지 않는다.

추가 고차 진단은 원래1substep인1/3840초 구간만 대조했다. 이 짧은 구간의
Newmark128/256분할은 독립 검산 통과 및 끝 속도 상대 질량 RMS 차이0.1056\%를 보였다.
256분할을 비교 기준으로 Newmark32분할의 끝 속도 차이는2.0733\%, 기존6차Gauss8분할은0.1499\%였다.
기준 자체가 exact solution인 것은 아니며 두 구간/오차 분모를 혼합하지 않는다.
당시 Gauss 시제품은 CPU 보조 분해 비용으로 더 빠르지 않았고 해당1/500 진단의 독립
stage/힘/에너지/기하 검산은 미연결이었다. 후속 GPU 구현과 독립 검산은 앞 절에서 구분한다.
따라서 이 선행 결과는 고차 적분의 정확도 개선 가능성을 보인 개발 근거로만 보존한다.
정확도 기준·물성·감쇠·학습 label·기본 Newmark/FP64 hi/lo는 변경하지 않았다.
동일 rest부터 전체 시간/공간 수렴, 실제 학습 label과 GS 매핑 및 접촉 검증은 미완료다.
상세 수치·명령·원본/hash는 \code{timestep_search/cloth_coarse/bend500_accuracy.md}가 소유한다.

\subsection{1/500 실패 프레임의 시간 간격 진단}
2026-09-14 별도 GPU 재현에서 바람2.01484375초의 GMRES720회 한도 실패를 확인했다.
보조 행렬 갱신16회와 GMRES1440회 후보는 실패했고, dt1/7680은2.00--2.05초의
384 substep을 독립 검산과 함께 통과했다. 원본 dt1/3840이나 실패 처리 정책은 변경하지 않았다.
다만 dt1/7680 대1/15360의 같은1프레임 위치 성분 차이는 최대0.364mm, 속도는1.605m/s이며,
더 작은 dt에서도 속도 차이가 남았다. 이는 풀이 실패를 피한 개발 후보이지 teacher 시간 수렴
또는 장기 안정성 완료가 아니다. 적응 substep/rollback과 접촉 모델은 미구현이다.
설정·타이머·원본 근거는 \code{timestep_search/cloth_coarse/bend500_failure_diagnosis.md}를 따른다.

\subsection{후속 GPU와 굽힘 비교의 검증 범위}
2026-09-14 확인한 저장 근거다. 이번 문서 갱신에서 시뮬레이션을 재실행하지 않았다.
동일한 과거10프레임/640단계에서 생성332.401초에서45.050초로 개선된 관측과
기존 경로 대조 통과는 \code{timestep_search/gpu_resident/optimization_summary.md}가 소유한다.
전체10초 또는 다른 메시 성능으로 일반화하지 않는다. 현재 매 프레임 GPU 완료 확인과
별도 시간 측정 경계 때문에 과거 비동기 기록과 직접 비교하지 않는다.

깃발1/100과1/300은 각각 중력2초+무풍4초+바람4초,38400 substep의 국소/물리 flag가0이다.
중력만의 최대 하강은 약0.132mm로 가시적 처짐 준비는 달성하지 못했다.
1/300은 바람 궤적이 달라졌고 사용자는 더 자연스럽다고 평가했지만 잔주름 정량 개선이나
자기 교차 부재를 검증한 것은 아니다. 굽힘 조건별 실행 시간은 부하/궤적 차이를 포함하므로
최적화 가속률로 해석하지 않는다. 별도 최초 하이브리드 재비교는 사용자 종료로 보류되어
양쪽 완주 성능/동등성 근거가 없다.

상세 설정·원본/hash·수치는 \code{experiments/R1_teacher_velocity_reset/timestep_search/cloth_coarse/}의
\code{gravity_flag_results.md}, \code{gravity_bend300_results.md},
\code{gravity_bend500.md}, \code{gravity_speed_compare.md}를 따른다.


\begin{longtable}{L{0.22\linewidth}L{0.46\linewidth}L{0.22\linewidth}}
\toprule
Fixture & 합격 조건 & 실패 owner\\
\midrule
\endfirsthead
\toprule
Fixture & 합격 조건 & 실패 owner\\
\midrule
\endhead
Teacher units/quadrature & analytic patch의 area/mass/force/work와 SI sign이 일치 & physics registry/quadrature\\
Spatial convergence & fixed accepted timestep에서 probe velocity/tip/work/FRF가 threshold 통과 & mesh/model\\
Temporal convergence & fixed accepted mesh에서 같은 metric이 threshold 통과 & integrator 또는 timestep\\
Map reproduction & partition, constant/affine field, coverage와 common-valid mask가 threshold 통과 & probe map\\
Adjoint work & total force와 traction 두 payload 각각의 virtual work identity 통과 & map/payload type\\
Traction identity & teacher/GS law/effective-scale/guard/sign hash 일치, zero guard-activation과 normal/sign/drag fixture 통과 & traction registry\\
% [논문 작성 참고 -- 수정 이유]
% 이 fixture는 여러 empirical coefficient의 개별 정확도가 아니라 동일 effective load와 비활성 safety guard를 검증한다.
Exact ZOH & random held-load reference, one-update trace와 aero-off passive decay 통과 & reduced evaluator\\
Oracle invariants & attachment, mass Gram, positive pole/damping, fit response와 free decay 통과 & oracle, measure 또는 evaluator\\
Independent-GS paired oracle & 같은 common mask에서 teacher 및 두 GS 차이가 declared reference+mapping floor 안 & map/oracle\\
Base transport & rest identity, proper rotation, SPD/finite covariance, attachment zero와 velocity consistency 통과 & transport branch\\
Leakage/replay & target artifact에 topology가 없고 동일 command/seed에서 semantic hash 재현 & artifact boundary\\
\bottomrule
\end{longtable}

각 fixture는 teacher spatial error, teacher temporal error, mapping error, oracle fit error와 transport error를
한 숫자로 합치지 않고 component별로 기록한다.

\subsection{사용자 GPU smoke: 실행 가능성과 수렴}

사용자가 \code{bash code/scripts/check_teacher_gpu.sh}를 실행했다.
GTX 1080 Ti 11 GiB, sm\_61, Warp 1.17.0, CUDA Toolkit 12.9/Driver 13.0에서
wind 5개, decay 2개, compare 3개, replay 2개의 \textbf{12/12 단계}가 통과했다.
Raw/probe 7쌍이며 pin drift와 guard activation은 0이다. Agent는 원본을 검토·재계산했고
추가 GPU simulation을 수행한 것으로 보고하지 않았다.
Raw: \code{experiments/artifacts/runs/teacher_gpu_check/20260907_042524_687651}.

공간 velocity 상대 RMS는 mesh 2$\to$4 10.07\%, 4$\to$8 24.90\%; temporal substeps
4$\to$8 23.99\%, 8$\to$16 11.41\%; aero-off decay mesh 4$\to$8 193.25\%였다.
분모는 fine run의 질량·시간 RMS 속도다. 짧은 0.2 s smoke이고
\code{convergence_status=not_assessed}다. 이는 발산률이나 변위 오차가 아니다.
같은 초기 probe 변위인데 native bending energy가 mesh별로 다른 관찰이 다음 감사를 이끌었다.

\subsection{시간 적분 정밀도와 시간 해상도 검증}

\subsubsection{초기 dynamics와 독립 reference}

Triangle-lumped total mass, position-only pin, Newmark/GMRES/line search를 갖는 CPU shell solver를
만들고 reaction, EOM residual, energy/work 및 모든 성공 state를 기록했다.
19 rollout, 2,105 step과 독립 재실행은 수치 계약을 통과했지만 nonlinear velocity의
160$\to$320 시간 대조는 11.394\%로 실패했다.
독립 DOP853 reference는 자체 대조를 통과했다.
추가 refinement에서 full N2560 velocity 11.7491\%, short N5120 5.50795\%였고,
short N10240은 5번째 step line search에서 실패했다. XY 고주파 시간 오차와
가까운 위치값을 빼서 가속도를 얻는 precision cancellation을 구분했다.
동일 시작 state에서 acceleration-form의 새 step 5 잔차는
\(3.12274\times10^{-10}\) m/s\(^2\)로 기존 bound \(1.34886\times10^{-8}\) m/s\(^2\)보다 작았다.
같은 저장 위치에 기존 위치 역산식을 적용하면 잔차가 \(2.45592\times10^{-8}\) m/s\(^2\)로 bound를 넘는다.
Full N2560에서 두 경로의 정규화 속도 차이는 \(3.31711\times10^{-9}\)였다.
계산 중단을 일으킨 cancellation의 보정과 실제 고주파 시간 오차 감소는 별도 문제다.

\subsubsection{교체한 acceleration-form의 refinement 결과}

현행 갱신식과 precision 보정 이유는 \S\ref{sec:r1-current-integrator}의
식~\eqref{eq:r1-acceleration-newmark}를 따른다. 같은 물리/입력을 유지한 실제 결과는 다음과 같다.

\begin{longtable}{L{0.20\linewidth}L{0.30\linewidth}L{0.40\linewidth}}
\toprule 구간 & refinement & velocity 차이와 판정\\\midrule
Short \(T_*/20\) & N10240/20480/40960 & 2.20620\% $\to$ 0.612308\% $\to$ 0.154044\%; 통과\\
Full \(T_*\) & N20480/40960/81920 & 6.059758\% $\to$ 2.974723\% $\to$ 0.868281\%; 통과\\
\bottomrule\end{longtable}

Full fixture의 \(T_*=1.0719083180935054\) s,
\(\omega_*=5.861681639298341\) rad/s는 n4 forward의 첫 positive normal mode 기준이다.
초기 amplitude는 0.001 m다. Short run은 full 기준 N의 분할 구간이며 새 계산 3,072 step/run이다.
Full은 143,360 step/run, 두 독립 run을 수행했다. 각 143,363 state, 560 chunk,
1,576,960 iteration/trial vector를 보존하고 state chain, force/reaction, EOM, Newmark update,
correction/Armijo/HVP와 비교 보고서를 재검산했다. Finest energy drift 기준은 \(10^{-3}\)이다.
각 run 약 45.7--45.8분, inventory 약 2.183 GB, 두 run 합 약 4.37 GB였다.
CPU/I/O 경합을 포함하므로 순수 solver 성능 benchmark가 아니다.
\textbf{이 통과는 해당 기존 shell과 초기 상태의 시간 진단이며 P3 teacher의 시간 통과가 아니다.}

\subsection{선형 공간 응답 실패를 확정한 실험}

\subsubsection{입력·경계·측정 고정}

1 m XY square, \(E=10^6\) Pa, \(\nu=0.3\), \(h=0.01\) m,
\(M_{\rm ref}=0.1\) kg, \(n=4,8,16\), forward/backward/checkerboard 세 diagonal이다.
\(x=0\)의 position만 고정하고 slope는 자유여서 normal rigid rotation null 하나가 남는다.
Clamped cantilever와 다른 경계다. \(u_z^0=10^{-3}x^2\), \(v^0=0\),
gravity/aerodynamics/damping 모두 off이며 동일 physical time \(T_*\), 10,241개 시각을 쓴다.

기존 shell의 rest-normal HVP에서 \(K\)를 조립하고 lumped \(M\)를 사용했다.
Mass-whitened generalized modal basis 전체를 써서 cos/sin exact response를 계산했다.
초기 상태를 mesh별 모드에 맞춰 다시 고르거나 high mode를 버리지 않았다.
Eigenvalue 원본을 보존하고 \(10^{-9}\max|\lambda|\) cutoff의 rigid null 하나를 별도 처리했다.
Time integration error를 제거한 상태에서도 공간 실패가 남는지 보는 진단이다.

Overlay16의 각 square를 네 triangle로 나누고 degree-2 3점 quadrature를 써 3,072개 probe를 만들었다.
이 규칙은 squared P1 difference에 정확하며 P3/spline 비교에는 충분하지 않다.
XY에서 기존 XZ mapper로 보낼 때 \((x,y,z)\mapsto(x,-z,y)\) proper rotation을 사용하고
dyadic rest position의 float32 표현을 확인했다. 평가 후 field를 원래 frame으로 되돌렸다.
Velocity는 \(A\omega_*\), displacement는 \(A=0.001\) m로 normalize하고 finest 1\%와
오차 감소를 사전 기준으로 썼다. Numerical floor는 별도 \(10^{-10}\)이다.

\begin{longtable}{L{0.27\linewidth}L{0.28\linewidth}L{0.35\linewidth}}
\toprule 비교 & displacement [\%] & velocity [\%]\\\midrule
Forward/backward 4$\to$8 & 5.956244 & 22.169172\\
Forward/backward 8$\to$16 & 2.927207 & 23.209961\\
Checkerboard 4$\to$8 & 3.857861 & 21.043244\\
Checkerboard 8$\to$16 & 1.395752 & 22.685509\\
n16 forward/backward & 0.570938 & 8.558328\\
n16 forward(or backward)/checker & 21.203509 & 51.546194\\
\bottomrule\end{longtable}

Maximum energy drift \(6.0021\times10^{-12}\), EOM \(1.4526\times10^{-12}\)로
수치 계산은 통과했지만 공간/방향 기준은 실패했다. 각 run 92,169 frame, 720 chunk,
1,509 inventory 약 266 MB다. 독립 두 run에서 runtime 외 1,508개 파일이 동일했다.
이 실패 뒤 dataset을 확대해 물리 결함을 덮지 않고 개발 sample만 분리 생성했다.

\subsection{두 입력군과 공간 norm·연속 시간 상계}
\label{sec:r1-pressure-comparison}

원래 \(w^0=10^{-3}x^2,v^0=0\)를 유지하고, 별도 진단으로 rest-start uniform pressure
\(p(t)=0.01\sin(\pi t/0.2)\) Pa를 \(0\le t\le0.2\) s에 가했다.
그 뒤 pressure는 0이다. 이 prescribed pressure는 relative wind quadratic traction도,
canonical frame-start held aerodynamic force도 아니다. 새 wind dataset로 소개하지 않는다.

모든 finite modal equation \(q''+\omega^2q=g\sin(at)\), \(a=\pi/0.2\)를 정확히 계산하고,
zero/resonant mode에는 안정적인 sinc/극한식, off 구간에는 homogeneous solution을 썼다.
기존 10,241개 시각과 \(A=0.001\) m, \(A\omega_*\) scale을 그대로 사용했다.
P3--P3 difference는 overlay16 Duffy4로 degree-6 squared field를 적분한다.
P3--tensor cubic spline의 교차항 \(u_{\mathrm{P3}}u_{\mathrm{spline}}\)은
overlay32 Duffy6로 total degree 최대 9의 곱을 적분한다.
각 모델의 self-norm은 자체 mass로 계산하며 spline squared field 전체를 Duffy6으로 적분한다는 뜻이 아니다.
P1 전용 3,072 probe 규칙을 그대로 재사용하지 않았다.

Mass-normalized continuous surface RMS를 \(\|\cdot\|\)라 하면 sampled maximum만 보고
시각 사이 peak를 놓치지 않도록 다음 finite-modal continuous-time 상계를 추가했다.
\[
 \max_{0\le t\le T_*}\|e(t)\|
 \le \max_j\|e(t_j)\|+L_e\frac{\Delta t}{2},\qquad
 L_u\ge\max_t\|\dot e_u(t)\|,\quad L_v\ge\max_t\|\dot e_v(t)\|.
\]
Pair의 각 model velocity/acceleration bound를 합하고 consistent-mass orthonormal basis와
\(\sqrt{M_{\rm ref}}\) normalization을 썼다. Coarse sampled lower와 fine continuous upper로
refinement 감소도 검사한다. 이는 주어진 finite modal 모델의 시간 sample 사이 상계이며
미해결 continuum error나 floating-point uncertainty 전체의 수학적 인증은 아니다.

Pressure mode의 \(|v|\)에는 \(2|g|/a\), on-interval \(|q|\)에는
\(\min(|g|\tau^2/2, |g|(1+a/\omega)/|\omega^2-a^2|)\) 형태의 유효 bound를 쓰고,
zero/resonance에서는 특수 fallback을 사용한다. Acceleration은 on에서
\(|g|+\omega^2|q|\), off에서 \(\operatorname{hypot}(\omega^2q_\tau,\omega v_\tau)\)다.
고주파를 제거하거나 입력별로 norm 분모를 재조정하지 않았다.

\begin{longtable}{L{0.39\linewidth}L{0.20\linewidth}L{0.29\linewidth}}
\toprule Pressure 비교 & displacement [\%] & velocity [\%]\\\midrule
Spline 8$\to$16 sampled & 0.0176344 & 0.191323\\
Spline 16$\to$32 sampled & 0.00105483 & 0.0234081\\
P2 8$\to$16 sampled 최대 & 2.245875 & 6.079196 (실패)\\
P3 8$\to$16 sampled 최대 & 0.016764 & 0.136376\\
P3 8$\to$16 continuous upper 최대 & 0.149654 & 0.526612 (통과)\\
P3 n16 direction continuous upper 최대 & 0.133208 & 0.398941 (통과)\\
P3 n16/spline32 continuous upper 최대 & 0.134202 & 0.408232 (통과)\\
\bottomrule\end{longtable}

P2 direction velocity도 1.106636\%로 실패했다. P3의 pressure 공간/방향/기준 비교는 고정 1\%를
통과했다. P3 free-energy drift 최대는 \(1.0684320\times10^{-8}\), pressure의 trapezoid
work/energy maximum defect는 \(1.5848926\times10^{-7}\)이다.
Work 수치는 정확 연속 시간 work 적분을 계산한 값으로 표기하지 않는다.

반면 원래 \(x^2\) 초기 상태의 P3 8$\to$16 velocity sampled error는
17.5895--17.9281\%, continuous upper는 약 24.27--24.80\%로 실패한다.
Displacement sampled error는 0.3017--0.3042\%다.
Spline의 같은 입력 velocity는 4$\to$8 27.3165\%, 8$\to$16 19.1290\%,
16$\to$32 14.9829\%로 감소하지만 아직 1\%에 못 미친다.
Spline32 초기 에너지의 16.8013\%가 100 rad/s 위, 8.4490\%가 500 rad/s 위에 있다.
이 초기 상태는 유효한 finite-energy 상태이며 ``불법 초기 조건''으로 제외하지 않는다.
이 입력은 \(u_z^0=10^{-3}x^2,\ v^0=0\)이다. 정상 pressure 응답과 고주파 성분이 있는 초기 변위에서
발생한 속도 응답의 수렴 난이도 차이로 해석하고 추가 검증 대상으로 둔다.

\section{종료 조건과 실패 경로}
\label{sec:r1-chapter-8}
\label{sec:r1-exit}

\subsection{R1 종료 조건}

다음이 모두 참일 때만 R2가 소유하는 T1 measure/token initialization prephase를 시작한다.
T1 산출물이 동결된 뒤에만 R2a fixed-seconds overfit으로 진행한다.
\begin{itemize}
  \item[$\square$] Teacher physics/material/solver/traction registry와 replay가 닫혔고
  teacher quadrature의 total area/mass가 \(A_{\mathrm{ref}},M_{\mathrm{ref}}\)와 일치한다.
  \item[$\square$] Spatial과 temporal convergence가 독립적으로 frozen threshold를 통과했다.
  \item[$\square$] 모든 declared GS variant가 하나의 common-valid probe denominator와 mapping fixture를 통과했다.
  \item[$\square$] Teacher/GS traction identity, reduced exact-ZOH와 aero-off free-decay fixture가 통과했다.
  \item[$\square$] Oracle canonical path 하나가 동결됐고 corrected oracle의 attachment/mass/response/paired fixture가 통과했다.
  \item[$\square$] Base mean/covariance transport가 rest/proper-rigid/SPD/attachment/velocity fixture를 통과했다.
  \item[$\square$] Spectrum convergence와 peak availability status를 분리해 기록했다.
\end{itemize}

\begin{stopbox}{R1 kill과 owner별 routing}
Teacher spatial/temporal reference가 수렴하지 않으면 physics model/mesh/timestep owner로 돌아간다.
Map reproduction/coverage가 실패하면 probe/map 또는 해당 GS package를 predeclared rule로 reject한다.
Traction/work/exact-ZOH가 실패하면 law/unit/evaluator owner를 고친다.
Corrected oracle가 mapping-noise floor 밖에서 실패하면 measure, rank, field fit, pole identification과 evaluator를
분리 진단한다. Base transport가 실패하면 normal/covariance branch를 고친다.
어느 경우에도 network, dataset 규모, Local 또는 renderer appearance를 추가하지 않는다.
\end{stopbox}

Converged spectrum에 stable peak가 없다는 사실만으로 R1을 실패시키지 않는다.
그 경우 \code{T_dom=unavailable}과 reason을 R2로 넘기며, R2b는 fixed-seconds fallback과
curriculum 비주장 경로를 사용한다.

\subsection{현행 계약의 채택 경계와 다음 진입 조건}

\begin{longtable}{L{0.23\linewidth}L{0.33\linewidth}L{0.34\linewidth}}
\toprule 항목 & 반영한 판단 & 아직 완료하지 않은 것\\\midrule
Teacher sample/DOF & 고차 DOF와 positive quadrature를 구분. Consistent mass는 식~\eqref{eq:r1-consistent-mass} 적용 & 최종 registry의 고차 law 및 actual model validator\\
Map sign & GS/기존 P1은 nonnegative 유지. 고차 Teacher에만 signed shape와 별도 version 허용 & P3 10-shape public map, coverage/noise/adjoint와 trajectory 연결\\
Convergence identity & 같은 backend/mass/BC/load에서 spatial/time acceptance 요구. 약한 바람 개발 비교 통과 & 4배 마지막 reset 공간 미달,10초 후보 미확정, 독립 기준·FRF/tip/work 검증\\
학습 적격성 & 물리 수식 구현·검증 전 추가 발행 보류. 기존 sample은 개발 근거로 보존 & Accepted Teacher와 source split, independent GS 및 oracle\\
첫 input family & P3 작은 굽힘 wind/reset의 수렴·독립 재실행과76개 sample 검증 & 큰 변형/새 입력 범위·oracle/GS 연결; native 실패와 기존 decay 요구 유지\\
Claim & P3는 training-only reference 후보, target method 변경 없음 & Learned held-out 개선이나 full Teacher 적격성 주장 불가\\
\bottomrule\end{longtable}

개발 full shell에는 objective membrane+bending, finite rotation, consistent mass,
attachment/reaction, nonlinear tangent/solve와 two-sided relative wind,
frame-start hold 및 work quadrature를 구현했다. 현재 구조 감쇠는0이며 상대 공기저항에 의한 감쇠를
구조/modal damping 채택과 혼동하지 않는다. 최종 damping/guard/고차 public registry 계약과의 일치는 별도로 닫아야 한다.
선형 P3 pressure 모듈을 기존 3-node/lumped-mass nonlinear solver에 이름만 바꿔 넣지 않는다.
개발 backend의 일부 시간·공간 진단 통과만으로 canonical producer를 채택하지 않는다.
Independent GS 두 개, common mask, spectrum convergence와 network-free oracle/base transport는 그 후에도
별도 R1 종료 조건이다. Student 학습이나 대규모 dataset 생성은 아직 수행하지 않았다.

\end{document}
